% Fonctions usuelles — Mathématiques Tle Tech — Cours signé OnBuch+
% Programme officiel MINESEC. Compiler avec : tectonic N04-fonctions-usuelles.tex
\newcommand{\DOCMATIERE}{Mathématiques}
\newcommand{\DOCNIVEAU}{Tle Tech}
\newcommand{\DOCMODULE}{Mathématiques – classe de Terminale technique}
\newcommand{\DOCLECON}{N04}
\newcommand{\DOCTITRE}{Fonctions usuelles}
% OnBuch+ — préambule « cours signés » ENRICHI (compilé avec tectonic / XeLaTeX)
% Système visuel « Pop » d'OnBuch+ : fond crème (jamais blanc), bordures encre
% épaisses, ombres « dures » décalées, titres Archivo Black en capitales.
% Version MATHÉMATIQUES : graphes (pgfplots/tikz), boîtes pédagogiques
% (définition, propriété, méthode, attention, le savais-tu, exercice résolu,
% exercice, TP, bilan), page de garde et signature OnBuch+ ancrée en bas.
%
% Macros attendues dans main.tex AVANT \input{preamble} :
%   \DOCMATIERE  \DOCNIVEAU  \DOCMODULE  \DOCLECON (numéro)  \DOCTITRE
\documentclass[11pt]{article}

\usepackage{fontspec}
\usepackage[french]{babel}
\usepackage[a4paper,margin=2cm,top=3.4cm,bottom=2.8cm,headheight=1.4cm,footskip=1.3cm]{geometry}
\usepackage{amsmath,amssymb,amsfonts}
\usepackage{mathtools} % \xmapsto, \coloneqq... souvent utilisés par les modèles
\usepackage{cancel} % simplifications barrées
% \abs{z} (module, valeur absolue) et \norm{u} (norme), utilisés par les modèles
\providecommand{\abs}{}\let\abs\relax\DeclarePairedDelimiter{\abs}{\lvert}{\rvert}
\providecommand{\norm}{}\let\norm\relax\DeclarePairedDelimiter{\norm}{\lVert}{\rVert}
\usepackage[table]{xcolor} % [table] : fournit aussi \rowcolors (alternance de lignes)
\usepackage{pagecolor}
\usepackage{tikz}
\usetikzlibrary{arrows.meta,positioning,calc,shapes.geometric,shapes.misc,shapes.arrows,decorations.pathmorphing,decorations.pathreplacing,decorations.markings,patterns,shadows,mindmap,trees,fit,backgrounds,matrix,intersections,angles,quotes}
\usepackage{pgfplots}
\usepackage{pgf-pie} % diagrammes circulaires \pie (statistiques)
\pgfplotsset{compat=1.18}
\usepgfplotslibrary{fillbetween,groupplots} % aires entre courbes (calcul intégral)
\usepackage{enumitem}
\usepackage{fancyhdr}
% [most] charge toutes les bibliothèques tcolorbox (listings, minted, raster,
% documentation...) et gonfle inutilement la mémoire TeX sur un document long
% (~300 boîtes) : on ne charge que ce qui est réellement utilisé.
\usepackage[skins,breakable]{tcolorbox}
\usepackage{graphicx}
\usepackage{colortbl}
\usepackage{booktabs}
\usepackage{tabularx}
\usepackage{diagbox} % case d'en-tête diagonale des tableaux à double entrée
% Colonnes extensibles centrée / gauche / droite (C, L, R), souvent utilisées par les modèles
\newcolumntype{C}{>{\centering\arraybackslash}X}
\newcolumntype{L}{>{\raggedright\arraybackslash}X}
\newcolumntype{R}{>{\raggedleft\arraybackslash}X}
\usepackage{array}
\usepackage{multicol}
\usepackage{siunitx}
% parse-numbers=false : accepte aussi \SI{2,5 \times 10^{-3}}{...}
\sisetup{locale=FR,output-decimal-marker={,},group-minimum-digits=4,parse-numbers=false}
\DeclareSIUnit\ohm{\ensuremath{\symup{\Omega}}} % U+2126 absent de la police
\usepackage{qrcode}
% \degree : commande fréquemment utilisée par les modèles pour un angle ou une
% température en dehors de \SI{}{} (non fournie par les paquets chargés).
\providecommand{\degree}{^{\circ}}
% \qed (fin de démonstration), utilisé par les modèles sans amsthm
\providecommand{\qed}{\hfill$\square$}
% \barre : double barre des valeurs interdites dans un tableau de variations (tkz-tab)
\providecommand{\barre}{\ensuremath{\|}}
% environnement proof (amsthm non chargé) utilisé par les modèles
\ifdefined\proof\else\newenvironment{proof}[1][Démonstration]{\par\noindent\textit{#1.}\ }{\qed\par}\fi
\usepackage{lastpage}
\usepackage{hyperref}
\hypersetup{hidelinks}

% ============================================================
% Palette « Pop » OnBuch+ — mêmes hex que la classe P (plus_ui.dart)
% ============================================================
\definecolor{poporange}{HTML}{F59321}
\definecolor{popdark}{HTML}{E07A0C}
\definecolor{popcream}{HTML}{FFF6E8}
\definecolor{popink}{HTML}{1C1714}
\definecolor{poppurple}{HTML}{7A5AE0}
\definecolor{popgreen}{HTML}{2E9E6A}
\definecolor{poppink}{HTML}{E0457B}
\definecolor{popblue}{HTML}{2F7DE1}
\definecolor{popgold}{HTML}{C98A12}
\definecolor{popmuted}{HTML}{8A7F76}
\definecolor{popline}{HTML}{EADFD2}
\definecolor{popgray}{HTML}{8A7F76}
\colorlet{popgrey}{popgray}
% Alias tolérants (le modèle nomme parfois les couleurs différemment)
\colorlet{popwhite}{white}
\colorlet{popblack}{popink}
\colorlet{popred}{poppink}
\colorlet{popyellow}{popgold}
% Teintes claires pour fonds de tableaux / zones
\colorlet{popblueL}{popblue!10!white}
\colorlet{poporangeL}{poporange!14!white}
\colorlet{popgreenL}{popgreen!12!white}
\colorlet{poppurpleL}{poppurple!12!white}
\colorlet{poppinkL}{poppink!12!white}
\colorlet{popgoldL}{popgold!14!white}

\pagecolor{popcream}

% ============================================================
% Typographie
% ============================================================
\defaultfontfeatures{Ligatures=TeX}
\setmainfont{PlusJakartaSans-Medium.ttf}[Path=../../../fonts/,BoldFont=PlusJakartaSans-Bold.ttf]
% Maths en sans-serif (Fira Math) avec chiffres et lettres droites de Plus
% Jakarta, pour que les formules se fondent dans le texte.
\usepackage[math-style=ISO,bold-style=ISO]{unicode-math}
\setmathfont{FiraMath-Regular.otf}[Path=../../../fonts/]
\setmathfont{PlusJakartaSans-Medium.ttf}[Path=../../../fonts/,range={up/num,up/{Latin,latin},\mathup}]
\usepackage{icomma} % virgule décimale sans espace en mode math
% Caractères mathématiques Unicode absents des polices de texte (Plus Jakarta,
% Archivo Black) : rendus par leur équivalent mathématique, en texte comme en
% formule — sinon ils s'affichent en carré vide (ex. « Arithmétique dans ℤ »).
\usepackage{newunicodechar}
\newunicodechar{ℝ}{\ensuremath{\mathbb{R}}}
\newunicodechar{ℤ}{\ensuremath{\mathbb{Z}}}
\newunicodechar{ℂ}{\ensuremath{\mathbb{C}}}
\newunicodechar{ℕ}{\ensuremath{\mathbb{N}}}
\newunicodechar{ℚ}{\ensuremath{\mathbb{Q}}}
\newunicodechar{𝔻}{\ensuremath{\mathbb{D}}}
\newunicodechar{α}{\ensuremath{\alpha}}
\newunicodechar{β}{\ensuremath{\beta}}
\newunicodechar{γ}{\ensuremath{\gamma}}
\newunicodechar{δ}{\ensuremath{\delta}}
\newunicodechar{ε}{\ensuremath{\varepsilon}}
\newunicodechar{θ}{\ensuremath{\theta}}
\newunicodechar{λ}{\ensuremath{\lambda}}
\newunicodechar{μ}{\ensuremath{\mu}}
\newunicodechar{σ}{\ensuremath{\sigma}}
\newunicodechar{τ}{\ensuremath{\tau}}
\newunicodechar{φ}{\ensuremath{\varphi}}
\newunicodechar{ψ}{\ensuremath{\psi}}
\newunicodechar{ω}{\ensuremath{\omega}}
\newunicodechar{Σ}{\ensuremath{\Sigma}}
% majuscules produites par \MakeUppercase dans les titres (α-aminés -> Α)
\newunicodechar{Α}{\ensuremath{\alpha}}
\newunicodechar{Β}{\ensuremath{\beta}}
\newunicodechar{Γ}{\ensuremath{\Gamma}}
\newunicodechar{Δ}{\ensuremath{\Delta}}
\newunicodechar{Ε}{\ensuremath{\varepsilon}}
\newunicodechar{Θ}{\ensuremath{\Theta}}
\newunicodechar{Λ}{\ensuremath{\Lambda}}
\newunicodechar{Μ}{\ensuremath{\mu}}
\newunicodechar{Τ}{\ensuremath{\tau}}
\newunicodechar{Φ}{\ensuremath{\Phi}}
\newunicodechar{Ψ}{\ensuremath{\Psi}}
\newunicodechar{Ω}{\ensuremath{\Omega}}
\newunicodechar{↦}{\ensuremath{\mapsto}}
\newunicodechar{∈}{\ensuremath{\in}}
\newunicodechar{∉}{\ensuremath{\notin}}
\newunicodechar{∧}{\ensuremath{\wedge}}
\newunicodechar{⇔}{\ensuremath{\Leftrightarrow}}
\newunicodechar{⟺}{\ensuremath{\Longleftrightarrow}}
\newunicodechar{⇒}{\ensuremath{\Rightarrow}}
\newunicodechar{⟹}{\ensuremath{\Longrightarrow}}
\newunicodechar{∘}{\ensuremath{\circ}}
\newunicodechar{∪}{\ensuremath{\cup}}
\newunicodechar{∩}{\ensuremath{\cap}}
\newunicodechar{≡}{\ensuremath{\equiv}}
\newunicodechar{⊂}{\ensuremath{\subset}}
\newunicodechar{⊥}{\ensuremath{\perp}}
\newunicodechar{∃}{\ensuremath{\exists}}
\newunicodechar{∀}{\ensuremath{\forall}}
\newunicodechar{∛}{\ensuremath{\sqrt[3]{\,}}}
\newunicodechar{∜}{\ensuremath{\sqrt[4]{\,}}}
\newunicodechar{⃗}{\ensuremath{{}^{\to}}}
\newunicodechar{ⁿ}{\ifmmode^{n}\else\textsuperscript{\ensuremath{n}}\fi}
\newunicodechar{ˣ}{\ifmmode^{x}\else\textsuperscript{\ensuremath{x}}\fi}
\newunicodechar{ᵇ}{\ifmmode^{b}\else\textsuperscript{\ensuremath{b}}\fi}
\newunicodechar{ʳ}{\ifmmode^{r}\else\textsuperscript{\ensuremath{r}}\fi}
\newunicodechar{ᵘ}{\ifmmode^{u}\else\textsuperscript{\ensuremath{u}}\fi}
\newunicodechar{ʸ}{\ifmmode^{y}\else\textsuperscript{\ensuremath{y}}\fi}
\newunicodechar{ᵃ}{\ifmmode^{a}\else\textsuperscript{\ensuremath{a}}\fi}
\newunicodechar{ᵉ}{\ifmmode^{e}\else\textsuperscript{\ensuremath{e}}\fi}
\newunicodechar{ᵏ}{\ifmmode^{k}\else\textsuperscript{\ensuremath{k}}\fi}
\newunicodechar{ᵖ}{\ifmmode^{p}\else\textsuperscript{\ensuremath{p}}\fi}
\newunicodechar{ᴺ}{\ifmmode^{N}\else\textsuperscript{\ensuremath{N}}\fi}
\newunicodechar{ᶜ}{\ifmmode^{c}\else\textsuperscript{\ensuremath{c}}\fi}
\newunicodechar{ʰ}{\ifmmode^{h}\else\textsuperscript{\ensuremath{h}}\fi}
\newunicodechar{ᵠ}{\ifmmode^{q}\else\textsuperscript{\ensuremath{q}}\fi}
\newunicodechar{ₙ}{\ifmmode_{n}\else\textsubscript{\ensuremath{n}}\fi}
\newunicodechar{ₐ}{\ifmmode_{a}\else\textsubscript{\ensuremath{a}}\fi}
\newunicodechar{ᵢ}{\ifmmode_{i}\else\textsubscript{\ensuremath{i}}\fi}
\newunicodechar{ᵣ}{\ifmmode_{r}\else\textsubscript{\ensuremath{r}}\fi}
\newunicodechar{ₖ}{\ifmmode_{k}\else\textsubscript{\ensuremath{k}}\fi}
\newunicodechar{ᵦ}{\ifmmode_{\beta}\else\textsubscript{\ensuremath{\beta}}\fi}
\newunicodechar{ₓ}{\ifmmode_{x}\else\textsubscript{\ensuremath{x}}\fi}
\newfontfamily\popdisplay{ArchivoBlack-Regular.ttf}[Path=../../../fonts/]
\newfontfamily\poplogo{SpaceGrotesk-Bold.ttf}[Path=../../../fonts/]
\newfontfamily\popsemi{PlusJakartaSans-SemiBold.ttf}[Path=../../../fonts/]
\newfontfamily\popxbold{PlusJakartaSans-ExtraBold.ttf}[Path=../../../fonts/]
\newfontfamily\popxboldls{PlusJakartaSans-ExtraBold.ttf}[Path=../../../fonts/,LetterSpace=3.0]

\setlist[enumerate]{leftmargin=1.7em,itemsep=5pt,topsep=4pt}
\setlist[itemize]{leftmargin=1.7em,itemsep=4pt,topsep=4pt,label={\color{poporange}\textbullet}}
\setlength{\parindent}{0pt}
\setlength{\parskip}{5pt}
\renewcommand{\arraystretch}{1.25}

% pgfplots : style Pop par défaut
\pgfplotsset{
  every axis/.append style={axis line style={popink,line width=1pt},tick style={popink},
    grid style={popline,line width=0.5pt},label style={font=\small},tick label style={font=\footnotesize}},
  popcurve/.style={poporange,line width=1.8pt,no markers,smooth},
}
% Même style disponible dans un \draw TikZ simple (hors pgfplots)
\tikzset{popcurve/.style={poporange,line width=1.8pt}}

% ============================================================
% Pill / squiggle
% ============================================================
\newcommand{\poppill}[3]{%
  \tikz[baseline=-0.55ex]\node[draw=popink,line width=1.2pt,rounded corners=2.6pt,%
    fill=#2,inner xsep=6.5pt,inner ysep=3pt]{%
    \popxboldls\fontsize{7}{7}\selectfont\color{#3}\MakeUppercase{#1}};%
}
\newcommand{\popsquiggle}[1]{%
  \tikz[baseline=-0.4ex]{
    \draw[#1,line width=2.6pt,line cap=round]
      plot [smooth] coordinates {(0,0.09) (0.34,0.01) (0.68,0.21) (1.02,0.01) (1.36,0.10)};
  }%
}
% Mot-clé surligné (marqueur orange sous le texte)
\newcommand{\cle}[1]{{\popxbold\color{popdark}#1}}

% ============================================================
% En-tête / pied
% ============================================================
\pagestyle{fancy}
\fancyhf{}
\renewcommand{\headrulewidth}{0pt}
\renewcommand{\footrulewidth}{0pt}
\lhead{\raisebox{-0.15ex}{\poplogo\fontsize{13}{13}\selectfont\color{popink}OnBuch\color{poporange}+}}
\rhead{\poppill{\DOCMATIERE\ \textbullet\ \DOCNIVEAU\ \textbullet\ Leçon \DOCLECON}{white}{popink}}
\lfoot{\popxbold\fontsize{7.6}{7.6}\selectfont\color{popmuted}\MakeUppercase{Cours signé OnBuch+}\ \textbullet\ {\popsemi\DOCTITRE}}
\rfoot{\tikz[baseline=-0.6ex]\node[rounded corners=8.5pt,fill=popink,minimum height=17pt,minimum width=17pt,inner xsep=5pt,inner sep=0]{\popxbold\fontsize{8}{8}\selectfont\color{popcream}\thepage\,/\,\pageref*{LastPage}};}

% ============================================================
% Titres
% ============================================================
\newcommand{\chaptitre}[2]{%
  \begin{tcolorbox}[enhanced,colback=poporange,colframe=popink,boxrule=2.6pt,arc=11pt,
      drop shadow=popink,left=15pt,right=15pt,top=12pt,bottom=11pt,before skip=3pt,after skip=18pt]
    \poppill{#2}{popink}{popcream}\par\vspace{9pt}
    {\raggedright\hyphenpenalty=10000\exhyphenpenalty=10000\popdisplay\fontsize{22}{24}\selectfont\color{popcream}\MakeUppercase{#1}\par}
  \end{tcolorbox}
}
% Section numérotée : pastille orange avec le numéro + titre Archivo
\newcounter{popsec}
\newcommand{\coursec}[1]{%
  \stepcounter{popsec}\setcounter{popsub}{0}%
  \par\vspace{16pt}\noindent
  \begin{minipage}{\linewidth}
  \tikz[baseline=-0.7ex]\node[draw=popink,line width=1.6pt,fill=poporange,rounded corners=4pt,minimum size=22pt,inner sep=0]{\popdisplay\fontsize{12}{12}\selectfont\color{popcream}\thepopsec};%
  \hspace{8pt}{\popdisplay\fontsize{15}{17}\selectfont\color{popink}\MakeUppercase{#1}}\par
  \vspace{4pt}\noindent{\color{poporange}\rule{36pt}{3.6pt}}
  \end{minipage}\par\nopagebreak\vspace{8pt}
}
\newcounter{popsub}
\newcommand{\courssub}[1]{%
  \stepcounter{popsub}%
  \par\vspace{9pt}\noindent{\popxbold\fontsize{11.5}{13}\selectfont\color{popink}{\color{poporange}\thepopsec.\thepopsub}\hspace{6pt}#1}\par\nopagebreak\vspace{4pt}
}

% ============================================================
% Boîtes pédagogiques — même recette : fond blanc, bordure épaisse colorée,
% ombre dure encre, badge plein. Titre optionnel [#1].
% ============================================================
\tcbset{popbase/.style={breakable,enhanced,colback=white,boxrule=2.3pt,arc=9pt,
  shadow={3pt}{-3pt}{0pt}{popink},left=14pt,right=14pt,top=11pt,bottom=11pt,before skip=11pt,after skip=15pt,
  fontupper=\popsemi\fontsize{10.6}{14.5}\selectfont\color{popink},
  fontlower=\popsemi\fontsize{10.6}{14.5}\selectfont\color{popink}}}
% Coche dessinée (le glyphe ✓ n'existe pas dans les polices du document)
\newcommand{\popcheck}[1]{\tikz[baseline=-0.5ex]\draw[#1,line width=1.6pt,line cap=round,line join=round](-0.13,0.0)--(-0.04,-0.09)--(0.13,0.1);}
\newcommand{\popboxhead}[3]{\poppill{#1}{#2}{white}\ifx\relax#3\relax\else\hspace{6pt}{\popxbold\fontsize{10.6}{12}\selectfont\color{popink}#3}\fi\par\vspace{7pt}}

\newtcolorbox{aretenir}[1][]{popbase,colframe=popgreen,before upper={\popboxhead{À retenir}{popgreen}{#1}}}
\newtcolorbox{exemplebox}[1][]{popbase,colframe=poporange,before upper={\popboxhead{Exemple}{poporange}{#1}}}
\newtcolorbox{definition}[1][]{popbase,colframe=popblue,before upper={\popboxhead{Définition}{popblue}{#1}}}
\newtcolorbox{propriete}[1][]{popbase,colframe=poppurple,before upper={\popboxhead{Propriété}{poppurple}{#1}}}
\newtcolorbox{methode}[1][]{popbase,colframe=popgold,before upper={\popboxhead{Méthode}{popgold}{#1}}}
\newtcolorbox{attention}[1][]{popbase,colframe=poppink,before upper={\popboxhead{Attention}{poppink}{#1}}}
\newtcolorbox{experience}[1][]{popbase,colframe=popblue,colback=popblueL,before upper={\popboxhead{Expérience}{popblue}{#1}}}
% « Le savais-tu ? » : carte violette pleine (contexte, culture, Cameroun)
\newtcolorbox{savaistu}[1][]{popbase,colback=poppurple,colframe=popink,
  fontupper=\popsemi\fontsize{10.4}{14.2}\selectfont\color{white},
  before upper={\poppill{Le savais-tu\,?}{white}{poppurple}\ifx\relax#1\relax\else\hspace{6pt}{\popxbold\color{white}#1}\fi\par\vspace{7pt}}}
% Exercice résolu : énoncé en haut, \tcblower puis la solution sur fond crème
\newcounter{popexo}\newcounter{popexr}
\newtcolorbox{exoresolu}[1][]{popbase,colframe=poporange,colbacklower=popcream,
  segmentation style={popink,line width=1.2pt,dashed},
  before upper={\stepcounter{popexr}\popboxhead{Exercice résolu \thepopexr}{poporange}{#1}},
  before lower={\poppill{Solution}{popink}{popcream}\par\vspace{6pt}}}
% Exercice (non résolu en place) — la correction est dans la partie Corrigés
\newtcolorbox{exercice}[1][]{popbase,colframe=popink,
  before upper={\stepcounter{popexo}\popboxhead{Exercice \thepopexo}{popink}{#1}}}
% Corrigé
\newtcolorbox{corrige}[1][]{popbase,colframe=popgreen,colback=popgreenL,
  before upper={\popboxhead{Corrigé}{popgreen}{#1}}}
% Objectifs / prérequis / bilan
\newtcolorbox{objectifs}{popbase,colframe=popink,colback=poporangeL,
  before upper={\popboxhead{Objectifs de la leçon}{popink}{}}}
\newtcolorbox{prerequis}{popbase,colframe=popink,colback=popblueL,
  before upper={\popboxhead{Prérequis}{popblue}{}}}
\newtcolorbox{bilan}{popbase,colframe=popink,colback=white,boxrule=2.8pt,
  before upper={\popboxhead{Fiche bilan}{poporange}{L'essentiel à connaître pour l'examen}}}
% Figure encadrée avec légende
\newtcolorbox{popfigure}[1][]{enhanced,breakable=false,colback=white,colframe=popline,boxrule=1.4pt,arc=8pt,
  left=10pt,right=10pt,top=10pt,bottom=8pt,before skip=10pt,after skip=12pt,halign=center,
  lower separated=false,halign lower=center,
  fontlower=\popsemi\fontsize{9}{11.5}\selectfont\color{popmuted},
  #1}
\newcommand{\legende}[1]{\par\vspace{6pt}{\popsemi\fontsize{9}{11.5}\selectfont\color{popmuted}#1}}

% ============================================================
% Signature OnBuch+
% ============================================================
\newcommand{\poplogoinline}[1]{{\poplogo\fontsize{#1}{#1}\selectfont\color{popink}OnBuch\color{poporange}+}}
\newcommand{\signatureonbuch}{%
  \begin{tcolorbox}[enhanced,colback=popink,colframe=popink,boxrule=2.2pt,arc=10pt,
      drop shadow=poporange,left=14pt,right=14pt,top=10pt,bottom=10pt,before skip=0pt,after skip=0pt]
    \tikz[baseline=-0.7ex]\node[circle,fill=popgreen,draw=popcream,line width=1.3pt,minimum size=22pt,inner sep=0]{\popcheck{white}};%
    \hspace{9pt}\begin{minipage}[c]{0.62\linewidth}
      {\popxbold\fontsize{12}{13}\selectfont\color{popcream}Signé\ \textbullet\ {\poplogo OnBuch\color{poporange}+}}\par\vspace{2pt}
      {\raggedright\popsemi\fontsize{8}{10.5}\selectfont\color{popline}\MakeUppercase{Équipe pédagogique OnBuch+}\\\MakeUppercase{Conforme au programme officiel MINESEC}\par}
    \end{minipage}\hfill
    {\poplogo\fontsize{22}{22}\selectfont\color{popcream}OnBuch\color{poporange}+}
  \end{tcolorbox}%
}

% ============================================================
% Page de garde
% #1 titre · #2 matière · #3 niveau · #4 module · #5 n° leçon · #6 durée ·
% #7 description / objectifs (texte ou itemize)
% ============================================================
\newcommand{\pagegarde}[7]{%
  \thispagestyle{empty}
  \newgeometry{a4paper,margin=1.7cm,top=1.6cm,bottom=1.4cm}%
  \begin{tikzpicture}[remember picture,overlay]
    \fill[poporange!18] ([xshift=-3.2cm,yshift=-3.6cm]current page.north east) circle (4.6cm);
    \fill[poppurple!14] ([xshift=2.4cm,yshift=7.6cm]current page.south west) circle (3.8cm);
  \end{tikzpicture}%
  {\poplogo\fontsize{30}{30}\selectfont\color{popink}OnBuch\color{poporange}+}\hfill
  \raisebox{0.6ex}{\poppill{Cours signé \textbullet\ Premium}{popink}{popcream}}\par
  \vspace{1pt}\popsquiggle{poppurple}\par
  \vspace{8pt}
  \begin{tcolorbox}[enhanced,colback=poporange,colframe=popink,boxrule=2.8pt,arc=13pt,
      drop shadow=popink,left=18pt,right=18pt,top=12pt,bottom=13pt,after skip=10pt]
    \poppill{#2 \textbullet\ #3}{popink}{popcream}\hspace{5pt}\poppill{#4}{white}{popink}\par\vspace{8pt}
    {\popdisplay\fontsize{14}{14}\selectfont\color{popink}LEÇON #5}\par\vspace{4pt}
    {\raggedright\hyphenpenalty=10000\exhyphenpenalty=10000\popdisplay\fontsize{25}{27}\selectfont\color{popcream}\MakeUppercase{#1}\par}
    \vspace{8pt}
    {\popxbold\fontsize{9.5}{9.5}\selectfont\color{popink}\MakeUppercase{Durée conseillée : #6}}
  \end{tcolorbox}
  \begin{tcolorbox}[enhanced,colback=white,colframe=popink,boxrule=2.2pt,arc=10pt,
      drop shadow=popink,left=16pt,right=16pt,top=10pt,bottom=10pt,after skip=8pt]
    \poppill{Dans cette leçon}{poppurple}{white}\par\vspace{5pt}
    \popsemi\fontsize{9.3}{12.6}\selectfont\color{popink}#7
  \end{tcolorbox}
  \noindent\begin{minipage}[t]{0.487\linewidth}
    \begin{tcolorbox}[enhanced,colback=popgreen,colframe=popink,boxrule=2.2pt,arc=10pt,
        drop shadow=popink,left=11pt,right=11pt,top=10pt,bottom=10pt,height=3.1cm,valign=center]
      \tcbox[colback=white,colframe=popink,boxrule=1.3pt,arc=5pt,boxsep=0pt,nobeforeafter,
        left=3pt,right=3pt,top=3pt,bottom=3pt,valign=center]{\qrcode[height=1.9cm]{https://play.google.com/store/apps/details?id=com.luvvix.onbuch&hl=fr}}%
      \hspace{8pt}\begin{minipage}[c]{0.5\linewidth}
        {\popdisplay\fontsize{11}{12.5}\selectfont\color{white}\MakeUppercase{Télécharge OnBuch}}\par\vspace{4pt}
        {\raggedright\popsemi\fontsize{8.4}{11}\selectfont\color{white}Gratuit \textbullet\ Google Play\\Tuteur IA, annales, concours, résultats\par}
      \end{minipage}
    \end{tcolorbox}
  \end{minipage}\hfill
  \begin{minipage}[t]{0.487\linewidth}
    \begin{tcolorbox}[enhanced,colback=poppurple,colframe=popink,boxrule=2.2pt,arc=10pt,
        drop shadow=popink,left=11pt,right=11pt,top=10pt,bottom=10pt,height=3.1cm,valign=center]
      \tikz[baseline=-0.7ex]\node[circle,fill=white,draw=popink,line width=1.3pt,minimum size=1.6cm,inner sep=0]{\popdisplay\fontsize{24}{24}\selectfont\color{poppurple}+};%
      \hspace{8pt}\begin{minipage}[c]{0.6\linewidth}
        {\popdisplay\fontsize{11}{12.5}\selectfont\color{white}\MakeUppercase{Passe à OnBuch+}}\par\vspace{4pt}
        {\raggedright\popsemi\fontsize{8.4}{11}\selectfont\color{white}Tous les cours signés, Tuteur IA illimité, fiches PDF — onglet OnBuch+ de l'app\par}
      \end{minipage}
    \end{tcolorbox}
  \end{minipage}
  \vspace{8pt}
  \popcontenu
  \vfill
  \signatureonbuch
  \restoregeometry
  \newpage
}

% Rangée « Ce cours contient » (page de garde)
\newcommand{\poptile}[3]{% #1 couleur #2 gros texte #3 légende
  \begin{tcolorbox}[enhanced,colback=#1,colframe=popink,boxrule=2pt,arc=9pt,shadow={2.5pt}{-2.5pt}{0pt}{popink},
      left=6pt,right=6pt,top=9pt,bottom=9pt,halign=center,valign=center,height=1.75cm,width=\linewidth,nobeforeafter]
    {\popdisplay\fontsize{12.5}{13}\selectfont\color{white}\MakeUppercase{#2}}\par\vspace{3pt}
    {\centering\popsemi\fontsize{7.8}{9.5}\selectfont\color{white}#3\par}
  \end{tcolorbox}}
\newcommand{\popcontenu}{%
  \par\noindent{\popxboldls\fontsize{7.5}{7.5}\selectfont\color{popmuted}CE COURS SIGNÉ CONTIENT}\par\vspace{7pt}
  \noindent\begin{minipage}[t]{0.235\linewidth}\poptile{popblue}{Cours}{complet, illustré, pas à pas}\end{minipage}\hfill
  \begin{minipage}[t]{0.235\linewidth}\poptile{poporange}{Méthodes}{et exercices résolus}\end{minipage}\hfill
  \begin{minipage}[t]{0.235\linewidth}\poptile{poppink}{Exercices}{type examen + corrigés}\end{minipage}\hfill
  \begin{minipage}[t]{0.235\linewidth}\poptile{popgreen}{Fiche bilan}{l'essentiel à retenir}\end{minipage}\par}

% Sommaire
\newtcolorbox{sommaire}{enhanced,colback=white,colframe=popink,boxrule=2.2pt,arc=10pt,
  drop shadow=popink,left=18pt,right=18pt,top=13pt,bottom=13pt,before skip=6pt,after skip=20pt,
  before upper={\poppill{Sommaire}{poppurple}{white}\par\vspace{9pt}\popsemi\fontsize{10.6}{14.5}\selectfont\color{popink}}}

% Signature de fin de cours — poussée tout en bas de la dernière page
\newcommand{\findecours}{%
  \par\vfill\nopagebreak
  \begin{center}\popsquiggle{poporange}\end{center}\vspace{4pt}
  \signatureonbuch
}
\AtBeginDocument{\def\llbracket{\mathopen{[\mkern-3.5mu[}}\def\rrbracket{\mathclose{]\mkern-3.5mu]}}} % intervalles d'entiers (glyphe ⟦ absent de la police)
% Glyphes absents de Fira Math : redéfinis à partir de glyphes disponibles
\AtBeginDocument{%
  \def\setminus{\mathbin{\backslash}}\let\smallsetminus\setminus
  \def\vdots{\mathord{\vbox{\baselineskip3.2pt\lineskiplimit0pt\kern4pt\hbox{.}\hbox{.}\hbox{.}}}}%
  \def\ddots{\mathinner{\mkern1mu\raise6.5pt\hbox{.}\mkern2mu\raise3.6pt\hbox{.}\mkern2mu\raise0.7pt\hbox{.}\mkern1mu}}%
  \def\triangle{\mathord{\scalebox{0.9}{$\symup{\Delta}$}}}%
  \def\nabla{\mathord{\raisebox{\depth}{\scalebox{1}[-1]{$\symup{\Delta}$}}}}%
  \def\checkmark{\popcheck{popdark}}%
  \def\star{\mathbin{\ast}}\def\bigstar{\mathord{\ast}}%
  \def\diamond{\mathbin{\scalebox{0.6}{$\Diamond$}}}%
  \def\bigcup{\mathop{\mathchoice{\raisebox{-0.3em}{\scalebox{1.7}{$\cup$}}}{\scalebox{1.25}{$\cup$}}{\cup}{\cup}}\displaylimits}%
  \def\bigcap{\mathop{\mathchoice{\raisebox{-0.3em}{\scalebox{1.7}{$\cap$}}}{\scalebox{1.25}{$\cap$}}{\cap}{\cap}}\displaylimits}%
  \def\lesseqqgtr{\mathrel{\lessgtr}}\def\bigoplus{\mathop{\oplus}\displaylimits}%
}
\providecommand{\ssi}{si et seulement si} % abréviation fréquente
\AtBeginDocument{\providecommand{\wideparen}{\overparen}} % arc de cercle

%% FIGFIX-BEGIN v3 — correctifs globaux des figures (docs/onbuch-generation/tools/figfix)
\usepackage{adjustbox}\usepackage{etoolbox}
% toute figure TikZ/pgfplots plus large que la ligne est réduite à la largeur disponible
\BeforeBeginEnvironment{tikzpicture}{\begin{adjustbox}{max width=\linewidth}}
\AfterEndEnvironment{tikzpicture}{\end{adjustbox}}
% légendes pgfplots lisibles par-dessus les courbes
\pgfplotsset{every axis/.append style={legend style={fill=white,fill opacity=0.92,text opacity=1,draw=black!15,inner sep=3pt}}}
% étiquettes d'arêtes (syntaxe edge["..."]) avec fond blanc
\tikzset{every edge quotes/.append style={fill=white,inner sep=1.2pt}}
% bulles de cartes mentales : texte à la ligne dans la bulle, bulle un peu plus grande
\tikzset{every concept/.append style={text width=2.0cm,align=center,minimum size=2.3cm,inner sep=2pt}}
%% FIGFIX-END
\begin{document}

\pagegarde{Fonctions usuelles}{Mathématiques}{Tle Tech}{Mathématiques – classe de Terminale technique}{N04}{8 séances (fiche de progression harmonisée)}{Cette leçon complète la boîte à outils des fonctions usuelles de Terminale : fonctions rationnelles, logarithme népérien et exponentielle népérienne. Elle fournit les méthodes d'étude, de résolution d'équations et d'inéquations, et de recherche de primitives indispensables pour le Probatoire et le Baccalauréat techniques.
\vspace{4pt}
\begin{multicols}{2}\small
\begin{itemize}
\item À la fin de cette leçon, je sais déterminer l'ensemble de définition, les limites et les asymptotes d'une fonction rationnelle
\item À la fin de cette leçon, je sais dresser le tableau de variations complet d'une fonction homographique ou rationnelle simple et tracer sa courbe
\item À la fin de cette leçon, je sais utiliser les propriétés algébriques de ln et de exp pour transformer des expressions
\item À la fin de cette leçon, je sais étudier les variations et les limites des fonctions ln, exp, x ↦ ln(ax+b) et x ↦ exp(ax+b)
\item À la fin de cette leçon, je sais résoudre des équations et inéquations faisant intervenir ln ou exp
\item À la fin de cette leçon, je sais déterminer les primitives des fonctions homographiques et des fonctions usuelles
\end{itemize}\end{multicols}}

\begin{sommaire}
\begin{multicols}{2}
\begin{enumerate}
\item Fonctions rationnelles et homographiques
\item La fonction logarithme népérien
\item La fonction exponentielle népérienne
\item Équations et inéquations avec ln et exp
\item Primitives des fonctions homographiques et usuelles
\item Étude des fonctions x ↦ ln(ax+b) et x ↦ exp(ax+b)
\item Applications à des situations concrètes
\item Activité d'intégration
\item Méthodes et exercices résolus
\item Exercices
\item Corrigés des exercices
\item Fiche bilan
\end{enumerate}
\end{multicols}
\end{sommaire}

\begin{objectifs}
\begin{itemize}
\item À la fin de cette leçon, je sais déterminer l'ensemble de définition, les limites et les asymptotes d'une fonction rationnelle
\item À la fin de cette leçon, je sais dresser le tableau de variations complet d'une fonction homographique ou rationnelle simple et tracer sa courbe
\item À la fin de cette leçon, je sais utiliser les propriétés algébriques de ln et de exp pour transformer des expressions
\item À la fin de cette leçon, je sais étudier les variations et les limites des fonctions ln, exp, x ↦ ln(ax+b) et x ↦ exp(ax+b)
\item À la fin de cette leçon, je sais résoudre des équations et inéquations faisant intervenir ln ou exp
\item À la fin de cette leçon, je sais déterminer les primitives des fonctions homographiques et des fonctions usuelles
\item À la fin de cette leçon, je sais modéliser une situation concrète de croissance ou de décroissance par une fonction exponentielle ou logarithme
\item À la fin de cette leçon, je sais rédiger et justifier une démarche complète d'étude de fonction et interpréter les résultats dans un contexte réel
\end{itemize}
\end{objectifs}

\begin{prerequis}
\begin{itemize}
\item Dérivation des fonctions usuelles et théorèmes sur les limites (classe de Première)
\item Fonctions polynômes du second degré et étude de signe
\item Notion de primitive et primitives des fonctions de référence (début d'année de Terminale)
\item Résolution d'équations et d'inéquations du premier et du second degré
\item Lecture graphique : asymptotes, tangentes, tableaux de variations
\end{itemize}
\end{prerequis}

\newpage

\coursec{Fonctions rationnelles et homographiques}

\noindent
\textbf{Situation d'accroche.} Mama Ngo Bell gère un dépôt de boissons à Douala. Chaque mois, le nombre de clients fidélisés double presque, tandis que le coût unitaire de stockage diminue quand le volume augmente. Comment prévoir dans combien de mois sa clientèle dépassera 10\,000 personnes ? Comment modéliser le bénéfice en fonction du prix de vente quand les coûts sont de type « $a + b/x$ » ? Ces questions, qui relèvent de la croissance exponentielle et des fonctions rationnelles, sont au cœur de cette leçon. Les fonctions logarithme népérien et exponentielle népérienne, outils puissants de l'analyse, permettront de répondre rigoureusement à ces situations de la vie courante. Commençons par maîtriser les fonctions rationnelles et leurs cas particuliers : les fonctions homographiques.

\courssub{Rappels : fonctions rationnelles et ensemble de définition}

\begin{definition}[Fonction rationnelle]
Une \cle{fonction rationnelle} est une fonction qui s'écrit sous la forme
\[
f(x) = \frac{P(x)}{Q(x)}
\]
où $P$ et $Q$ sont des fonctions polynomiales et $Q$ n'est pas le polynôme nul.
\end{definition}

\begin{propriete}[Ensemble de définition]
L'ensemble de définition $D_f$ d'une fonction rationnelle $f(x) = \frac{P(x)}{Q(x)}$ est l'ensemble des réels pour lesquels le dénominateur est non nul :
\[
D_f = \{ x \in \mathbb{R} \mid Q(x) \neq 0 \}.
\]
Les valeurs racines de $Q(x)=0$ sont des \cle{valeurs interdites} (elles ne font pas partie du domaine).
\end{propriete}

\begin{attention}[Piège classique : l'oubli des valeurs interdites]
\emph{Erreur fréquente :} Simplifier une expression rationnelle \textbf{avant} d'avoir déterminé l'ensemble de définition.
\par
Exemple : $f(x) = \frac{x^2-1}{x-1}$.
\begin{itemize}
\item \textbf{Mauvaise démarche :} On simplifie $f(x)=x+1$ et on dit $D_f=\mathbb{R}$.
\item \textbf{Bonne démarche :} $Q(x)=x-1$. La valeur interdite est $x=1$. Donc $D_f = \mathbb{R}\setminus\{1\}$.
Sur $D_f$, on a bien $f(x)=x+1$, mais la fonction n'est \textbf{pas définie} en $1$. Sa courbe est la droite $y=x+1$ \emph{privée du point} d'abscisse $1$ (trou).
\end{itemize}
\emph{Règle d'or :} Déterminer $D_f$ \textbf{avant} toute simplification.
\end{attention}

\begin{exemplebox}[Détermination de l'ensemble de définition]
Déterminer $D_f$ pour $f(x) = \frac{2x+3}{x^2-4}$.
\begin{itemize}
\item Dénominateur : $Q(x) = x^2-4 = (x-2)(x+2)$.
\item $Q(x)=0 \Leftrightarrow x=2$ ou $x=-2$.
\item $D_f = \mathbb{R} \setminus \{-2; 2\} = ]-\infty;-2[ \cup ]-2;2[ \cup ]2;+\infty[$.
\end{itemize}
\end{exemplebox}

\courssub{Les fonctions homographiques : définition et forme réduite}

\begin{definition}[Fonction homographique]
Une \cle{fonction homographique} est une fonction rationnelle de la forme
\[
f(x) = \frac{ax+b}{cx+d}
\]
avec $a,b,c,d \in \mathbb{R}$, $c \neq 0$ (sinon c'est une fonction affine) et $ad-bc \neq 0$ (sinon c'est une fonction constante).
\end{definition}

\begin{propriete}[Écriture réduite (forme canonique)]
Toute fonction homographique $f(x) = \frac{ax+b}{cx+d}$ ($c \neq 0$) peut s'écrire sous la forme
\[
f(x) = \alpha + \frac{\beta}{x - x_0}
\]
où
\[
\alpha = \frac{a}{c}, \quad x_0 = -\frac{d}{c}, \quad \beta = \frac{bc-ad}{c^2} = -\frac{ad-bc}{c^2}.
\]
Cette écriture s'obtient par une division euclidienne du numérateur par le dénominateur.
\end{propriete}

\begin{methode}[Passer de la forme générale à la forme réduite]
Pour $f(x) = \frac{ax+b}{cx+d}$ ($c \neq 0$) :
\begin{enumerate}
\item Effectuer la division euclidienne de $ax+b$ par $cx+d$.
\item Le quotient est $\alpha = \frac{a}{c}$.
\item Le reste est $b - \frac{ad}{c} = \frac{bc-ad}{c}$.
\item On obtient $f(x) = \frac{a}{c} + \frac{bc-ad}{c(cx+d)} = \frac{a}{c} + \frac{bc-ad}{c^2(x+\frac{d}{c})}$.
\item Identifier $\alpha = \frac{a}{c}$, $x_0 = -\frac{d}{c}$, $\beta = \frac{bc-ad}{c^2}$.
\end{enumerate}
\end{methode}

\begin{exemplebox}[Mise sous forme réduite]
Soit $f(x) = \frac{2x+3}{x-1}$.
\begin{itemize}
\item Division : $(2x+3) \div (x-1)$. Quotient $2$, Reste $5$.
\item $f(x) = 2 + \frac{5}{x-1}$.
\item Ici $\alpha = 2$, $\beta = 5$, $x_0 = 1$.
\end{itemize}
\end{exemplebox}

\courssub{Interprétation géométrique : hyperbole et centre de symétrie}

\begin{propriete}[Courbe représentative et centre de symétrie]
La courbe $\mathcal{C}_f$ de la fonction homographique $f(x) = \alpha + \frac{\beta}{x-x_0}$ est une \cle{hyperbole} dont les asymptotes sont les droites d'équations $x = x_0$ (verticale) et $y = \alpha$ (horizontale).
Le point $\Omega(x_0; \alpha)$ est le \cle{centre de symétrie} de l'hyperbole : pour tout $x \in D_f$, le point $M(x; f(x))$ a pour symétrique par rapport à $\Omega$ le point $M'(2x_0-x; 2\alpha-f(x))$ qui appartient aussi à $\mathcal{C}_f$.
\end{propriete}

\begin{popfigure}
\begin{tikzpicture}[scale=0.9, >=Stealth]
% Axes
\draw[->] (-4,0) -- (5,0) node[right] {$x$};
\draw[->] (0,-3) -- (0,5) node[above] {$y$};
% Graduations
\foreach \x in {-3,-2,-1,1,2,3,4} \draw (\x,0.1) -- (\x,-0.1) node[below] {\small $\x$};
\foreach \y in {-2,-1,1,2,3,4} \draw (0.1,\y) -- (-0.1,\y) node[left] {\small $\y$};
% Asymptotes
\draw[popblue, dashed, thick] (1,-3) -- (1,5) node[above right] {$x=1$};
\draw[popblue, dashed, thick] (-4,2) -- (5,2) node[right] {$y=2$};
% Centre de symétrie
\fill[poporange] (1,2) circle (2.5pt) node[above right] {$\Omega(1;2)$};
% Branche gauche (x < 1)
\draw[popcurve, thick, domain=-3:0.6, samples=100] plot (\x, {2 + 5/(\x-1)});
% Branche droite (x > 1)
\draw[popcurve, thick, domain=1.4:4, samples=100] plot (\x, {2 + 5/(\x-1)});
% Points remarquables
\fill[popgreen] (0, -3) circle (2pt) node[below left] {$A(0;-3)$};
\fill[popgreen] (-1.5, 0) circle (2pt) node[left] {$B(-1.5;0)$};
\fill[popgreen] (2, 7) circle (2pt) node[above right] {$C(2;7)$}; % hors cadre mais indiqué
% Vecteurs symétrie
\draw[poppurple, <->] (0.5, 0.5) -- (1,2) -- (1.5, 3.5);
\node[poppurple] at (0.2, 1.5) {\small Symétrie centrale};
\end{tikzpicture}
\legende{Courbe de $f(x) = 2 + \frac{5}{x-1} = \frac{2x+3}{x-1}$. Asymptotes $x=1$ et $y=2$, centre $\Omega(1;2)$.}
\end{popfigure}

\courssub{Limites, asymptotes et variations}

\begin{methode}[Déterminer les asymptotes d'une fonction homographique]
Pour $f(x) = \frac{ax+b}{cx+d}$ ($c \neq 0$, $ad-bc \neq 0$) :
\begin{enumerate}
\item \textbf{Asymptote verticale :} Chercher la valeur interdite $x_0 = -\frac{d}{c}$ (racine du dénominateur). Calculer $\lim_{x \to x_0^-} f(x)$ et $\lim_{x \to x_0^+} f(x)$. L'une tend vers $+\infty$, l'autre vers $-\infty$. La droite $x=x_0$ est asymptote verticale.
\item \textbf{Asymptote horizontale :} Calculer $\lim_{x \to -\infty} f(x)$ et $\lim_{x \to +\infty} f(x)$. Les deux limites valent $\alpha = \frac{a}{c}$. La droite $y = \frac{a}{c}$ est asymptote horizontale.
\end{enumerate}
\end{methode}

\begin{propriete}[Dérivée et sens de variation]
Soit $f(x) = \frac{ax+b}{cx+d}$ définie sur $D_f = \mathbb{R} \setminus \{-\frac{d}{c}\}$.
Sa dérivée est
\[
f'(x) = \frac{ad-bc}{(cx+d)^2}.
\]
Le signe de $f'(x)$ est constant sur chaque intervalle de $D_f$ (le dénominateur est un carré, toujours $>0$).
\begin{itemize}
\item Si $ad-bc > 0$, alors $f'(x) > 0$ : $f$ est \textbf{strictement croissante} sur $]-\infty; -\frac{d}{c}[$ et sur $]-\frac{d}{c}; +\infty[$.
\item Si $ad-bc < 0$, alors $f'(x) < 0$ : $f$ est \textbf{strictement décroissante} sur $]-\infty; -\frac{d}{c}[$ et sur $]-\frac{d}{c}; +\infty[$.
\end{itemize}
\end{propriete}

\begin{experience}[Démonstration guidée : signe de la dérivée]
On veut étudier le sens de variation de $f(x) = \frac{ax+b}{cx+d}$.
\begin{enumerate}
\item Calculer $f'(x)$ en utilisant la formule de dérivation d'un quotient : $(\frac{u}{v})' = \frac{u'v-uv'}{v^2}$ avec $u=ax+b$, $v=cx+d$.
\item $u' = a$, $v' = c$.
\item $f'(x) = \frac{a(cx+d) - c(ax+b)}{(cx+d)^2} = \frac{acx+ad - acx - bc}{(cx+d)^2} = \frac{ad-bc}{(cx+d)^2}$.
\item Analyser le signe : $(cx+d)^2 > 0$ pour tout $x \neq -\frac{d}{c}$. Le signe de $f'$ est donc celui de $ad-bc$.
\item Conclure sur la monotonie sur chaque intervalle de définition.
\end{enumerate}
\end{experience}

\begin{aretenir}[Tableau de variations type — Fonction homographique]
\begin{center}
\begin{tabular}{|c|ccccc|}\hline
$x$ & $-\infty$ & & $-\frac{d}{c}$ & & $+\infty$ \\ \hline
\multicolumn{6}{|c|}{\textbf{Cas $ad-bc > 0$ (Croissante)}} \\ \hline
$f'(x)$ & & $+$ & $\not\exists$ & $+$ & \\ \hline
$f(x)$ & $\frac{a}{c}$ & $\nearrow$ & $+\infty$ & & $\frac{a}{c}$ \\
       &           &       & $-\infty$ & $\nearrow$ &       \\ \hline
\multicolumn{6}{|c|}{\textbf{Cas $ad-bc < 0$ (Décroissante)}} \\ \hline
$f'(x)$ & & $-$ & $\not\exists$ & $-$ & \\ \hline
$f(x)$ & $\frac{a}{c}$ & $\searrow$ & $+\infty$ & & $\frac{a}{c}$ \\
       &           &       & $-\infty$ & $\searrow$ &       \\ \hline
\end{tabular}
\end{center}
\emph{Note :} La double ligne pour $f(x)$ au niveau de la valeur interdite rappelle que la fonction n'y est pas définie et que les branches ne se rejoignent pas.
\end{aretenir}

\begin{popfigure}
\begin{tikzpicture}[scale=0.8]
% Tableau de variations visuel pour ad-bc > 0
\node[anchor=north west, align=center] at (0,0){
\begin{tabular}{|c|ccccc|}\hline
$x$ & $-\infty$ & & $x_0$ & & $+\infty$ \\ \hline
$f'(x)$ & & $+$ & $\not\exists$ & $+$ & \\ \hline
$f$ & $\alpha$ & $\nearrow$ & $+\infty$ & & $\alpha$ \\
    &          &        & $-\infty$ & $\nearrow$ &    \\ \hline
\end{tabular}
};
\node[anchor=north west, text=popgreen] at (0,-2.5) {\textbf{Cas $ad-bc > 0$ : $f$ croissante sur chaque intervalle}};

% Cas ad-bc < 0
\node[anchor=north west, align=center] at (0,-4){
\begin{tabular}{|c|ccccc|}\hline
$x$ & $-\infty$ & & $x_0$ & & $+\infty$ \\ \hline
$f'(x)$ & & $-$ & $\not\exists$ & $-$ & \\ \hline
$f$ & $\alpha$ & $\searrow$ & $+\infty$ & & $\alpha$ \\
    &          &        & $-\infty$ & $\searrow$ &    \\ \hline
\end{tabular}
};
\node[anchor=north west, text=poporange] at (0,-6.5) {\textbf{Cas $ad-bc < 0$ : $f$ décroissante sur chaque intervalle}};
\end{tikzpicture}
\legende{Tableaux de variations types selon le signe de $ad-bc$.}
\end{popfigure}

\courssub{Étude complète : exemple résolu}

\begin{exoresolu}[Étude complète de $f(x) = \frac{x-2}{x+1}$]
\textbf{Énoncé :} Étudier la fonction $f(x) = \frac{x-2}{x+1}$ et tracer sa courbe $\mathcal{C}_f$.

\textbf{Solution détaillée :}
\begin{enumerate}
\item \textbf{Ensemble de définition :} $D_f = \mathbb{R} \setminus \{-1\}$ (valeur interdite $x=-1$).
\item \textbf{Forme réduite :} Division euclidienne de $x-2$ par $x+1$.
Quotient $1$, Reste $-3$.
$f(x) = 1 - \frac{3}{x+1}$.
Identification : $\alpha = 1$, $\beta = -3$, $x_0 = -1$.
\item \textbf{Asymptotes :}
\begin{itemize}
\item Verticale : $x = -1$.
$\lim_{x \to -1^-} f(x) = +\infty$ (car $x+1 \to 0^-$, $\frac{-3}{0^-}=+\infty$).
$\lim_{x \to -1^+} f(x) = -\infty$.
\item Horizontale : $\lim_{x \to \pm\infty} f(x) = 1$. Asymptote $y=1$.
\end{itemize}
\item \textbf{Dérivée et variations :}
$a=1, b=-2, c=1, d=1 \Rightarrow ad-bc = 1(1) - (-2)(1) = 3 > 0$.
$f'(x) = \frac{3}{(x+1)^2} > 0$ sur $D_f$.
$f$ est strictement croissante sur $]-\infty;-1[$ et sur $]-1;+\infty[$.
\item \textbf{Tableau de variations :}
\begin{center}
\begin{tabular}{|c|ccccc|}\hline
$x$ & $-\infty$ & & $-1$ & & $+\infty$ \\ \hline
$f'(x)$ & & $+$ & $\not\exists$ & $+$ & \\ \hline
$f(x)$ & $1$ & $\nearrow$ & $+\infty$ & & $1$ \\
       &     &       & $-\infty$ & $\nearrow$ &   \\ \hline
\end{tabular}
\end{center}
\item \textbf{Points remarquables :}
\begin{itemize}
\item Intersection avec l'axe des ordonnées : $x=0 \Rightarrow f(0)=-2$. Point $A(0;-2)$.
\item Intersection avec l'axe des abscisses : $f(x)=0 \Leftrightarrow x-2=0 \Rightarrow x=2$. Point $B(2;0)$.
\item Centre de symétrie : $\Omega(-1; 1)$.
\end{itemize}
\item \textbf{Tracé de la courbe :} Voir figure ci-dessous.
\end{enumerate}
\end{exoresolu}

\begin{popfigure}
\begin{tikzpicture}[scale=0.85, >=Stealth]
% Axes
\draw[->] (-5,0) -- (4,0) node[right] {$x$};
\draw[->] (0,-5) -- (0,4) node[above] {$y$};
% Graduations
\foreach \x in {-4,-3,-2,-1,1,2,3} \draw (\x,0.1) -- (\x,-0.1) node[below] {\small $\x$};
\foreach \y in {-4,-3,-2,-1,1,2,3} \draw (0.1,\y) -- (-0.1,\y) node[left] {\small $\y$};
% Asymptotes
\draw[popblue, dashed, thick] (-1,-5) -- (-1,4) node[above] {$x=-1$};
\draw[popblue, dashed, thick] (-5,1) -- (4,1) node[right] {$y=1$};
% Centre
\fill[poporange] (-1,1) circle (2.5pt) node[above right] {$\Omega(-1;1)$};
% Courbe branche gauche
\draw[popcurve, thick, domain=-4.5:-1.3, samples=100] plot (\x, {1 - 3/(\x+1)});
% Courbe branche droite
\draw[popcurve, thick, domain=-0.7:3.5, samples=100] plot (\x, {1 - 3/(\x+1)});
% Points
\fill[popgreen] (0,-2) circle (2pt) node[below right] {$A(0;-2)$};
\fill[popgreen] (2,0) circle (2pt) node[above right] {$B(2;0)$};
% Flèches variations
\draw[popgreen, thick, <->] (-3, -0.5) -- (-1.5, 2.5); % croissante gauche
\draw[popgreen, thick, <->] (0, -1) -- (2, 1); % croissante droite
\end{tikzpicture}
\legende{Courbe $\mathcal{C}_f$ de $f(x) = \frac{x-2}{x+1} = 1 - \frac{3}{x+1}$.}
\end{popfigure}

\courssub{Complément : Fonction rationnelle à dénominateur du second degré}

\begin{savaistu}[Au-delà des homographiques : dénominateur degré 2]
En Terminale Technique, on rencontre parfois des fonctions rationnelles de la forme $f(x) = \frac{P(x)}{Q(x)}$ où $\deg Q = 2$.
L'étude suit le même schéma mais demande plus de rigueur :
\begin{enumerate}
\item $D_f$ : résoudre $Q(x)=0$ (0, 1 ou 2 valeurs interdites).
\item Limites aux bornes de $D_f$ (asymptotes verticales aux racines de $Q$).
\item Limites à $\pm\infty$ : comparer degrés de $P$ et $Q$ pour asymptote horizontale ou oblique.
\item Dérivée : $f'(x) = \frac{P'Q - PQ'}{Q^2}$. Le signe se fait par un \textbf{tableau de signes} du numérateur $N(x)=P'Q-PQ'$ (polynôme) et du dénominateur $Q^2$ (toujours $\ge 0$).
\item Tableau de variations et courbe.
\end{enumerate}
\emph{Exemple type :} $f(x) = \frac{x}{x^2+1}$ (définie sur $\mathbb{R}$, paire, max en $1$, min en $-1$, asymptote $y=0$).
\end{savaistu}

\courssub{Application concrète : Coût unitaire moyen}

\begin{exemplebox}[Modélisation économique : Coût fixe et coût variable]
Mama Ngo Bell a un coût fixe mensuel de $500\,000$ FCFA (loyer, salaires) et un coût variable de $200$ FCFA par caisse de boissons stockée.
Si elle stocke $x$ caisses ($x>0$), le coût total est $C_T(x) = 500\,000 + 200x$.
Le \textbf{coût unitaire moyen} est $C_m(x) = \frac{C_T(x)}{x} = \frac{500\,000}{x} + 200$.
C'est une fonction homographique (de la forme $\alpha + \frac{\beta}{x}$ avec $\alpha=200$, $\beta=500\,000$).
\begin{itemize}
\item Lorsque $x$ augmente, $\frac{500\,000}{x}$ diminue : le coût fixe se « dilue » dans la production.
\item $C_m(x)$ est strictement décroissante (car $\beta > 0$ et $x>0$, dérivée $-\frac{500\,000}{x^2} < 0$).
\item Asymptote horizontale $y=200$ : le coût unitaire tend vers le coût variable pur quand la production devient très grande.
\end{itemize}
Cela explique pourquoi les gros volumes permettent des prix unitaires plus bas !
\end{exemplebox}

\begin{exercice}[Entraînement]
Soit la fonction $g(x) = \frac{3x-1}{2x+4}$.
\begin{enumerate}
\item Déterminer $D_g$ et écrire $g$ sous forme réduite.
\item Trouver les asymptotes et le centre de symétrie.
\item Calculer $g'(x)$ et dresser le tableau de variations.
\item Tracer la courbe $\mathcal{C}_g$ dans un repère.
\end{enumerate}
\end{exercice}

\begin{corrige}[Exercice 1]
\begin{enumerate}
\item $D_g = \mathbb{R} \setminus \{-2\}$. Division : $g(x) = \frac{3}{2} - \frac{7}{2(x+2)}$. ($\alpha=1.5, \beta=-3.5, x_0=-2$).
\item Asymptotes : $x=-2$ (verticale), $y=1.5$ (horizontale). Centre $\Omega(-2; 1.5)$.
\item $a=3, b=-1, c=2, d=4 \Rightarrow ad-bc = 12 - (-2) = 14 > 0$. $g'(x) = \frac{14}{(2x+4)^2} > 0$. $g$ croissante sur $]-\infty;-2[$ et $]-2;+\infty[$.
\item Tableau de variations standard pour $ad-bc>0$ avec $x_0=-2, \alpha=1.5$. Courbe : hyperbole centre $\Omega$, branches croissantes.
\end{enumerate}
\end{corrige}

\coursec{La fonction logarithme népérien}

Après avoir étudié les fonctions rationnelles et homographiques, nous nous sommes heurtés à une limite : la fonction $x \mapsto \frac{1}{x}$ n'admet pas de primitive sous forme de fonction rationnelle sur $\mathbb{R}^*$. Pour combler ce vide, les mathématiciens ont \emph{défini} une nouvelle fonction à partir de l'intégrale de $\frac{1}{x}$. C'est la naissance du logarithme népérien, outil indispensable pour modéliser des phénomènes où la croissance ralentit (chimie, acoustique, sismologie, finance).

\courssub{Définition et propriétés fondamentales}

\begin{definition}[Logarithme népérien $\ln x$]
On appelle \textbf{logarithme népérien} (ou logarithme naturel) la fonction $\ln$ définie sur $]0\,; +\infty[$ par :
\[
\forall x > 0,\quad \ln x = \int_1^x \frac{1}{t}\,\mathrm{d}t.
\]
C'est l'unique fonction dérivable sur $]0\,; +\infty[$ qui vérifie :
\begin{itemize}
    \item $(\ln x)' = \dfrac{1}{x}$ pour tout $x > 0$,
    \item $\ln 1 = 0$.
\end{itemize}
Le nombre $e$ (constante d'Euler, $e \approx 2,71828\ldots$) est l'unique nombre tel que $\ln e = 1$.
\end{definition}

\begin{propriete}[Conséquences immédiates]
\begin{itemize}
    \item \textbf{Valeurs remarquables :} $\ln 1 = 0$ ; $\ln e = 1$.
    \item \textbf{Dérivée :} $\forall x > 0,\quad (\ln x)' = \dfrac{1}{x}$.
    \item \textbf{Sens de variation :} Comme $\frac{1}{x} > 0$ sur $]0\,; +\infty[$, la fonction $\ln$ est \textbf{strictement croissante} sur $]0\,; +\infty[$.
    \item \textbf{Limites :}
    \[
    \lim_{x \to 0^+} \ln x = -\infty \quad\text{(asymptote verticale : l'axe des ordonnées $x=0$)},\qquad
    \lim_{x \to +\infty} \ln x = +\infty.
    \]
    \item \textbf{Signe de $\ln x$ :}
    \[
    \ln x \begin{cases}
    < 0 & \text{si } 0 < x < 1,\\
    = 0 & \text{si } x = 1,\\
    > 0 & \text{si } x > 1.
    \end{cases}
    \]
\end{itemize}
\end{propriete}

\begin{experience}[Démonstration de la stricte croissance]
On sait que $\ln$ est dérivable sur $]0\,;+\infty[$ et que $(\ln x)' = 1/x$.
Pour tout $x > 0$, $1/x > 0$. Donc $(\ln x)' > 0$ sur tout l'intervalle.
D'après le théorème de la variation (si $f' > 0$ alors $f$ est strictement croissante), $\ln$ est strictement croissante sur $]0\,;+\infty[$.
\end{experience}

\begin{popfigure}
\begin{tikzpicture}
\begin{axis}[
    width=\linewidth, height=8cm,
    axis lines=middle,
    xlabel=$x$, ylabel=$y$,
    xmin=-0.5, xmax=5,
    ymin=-3, ymax=2.5,
    xtick={1,2.718}, xticklabels={$1$,$e$},
    ytick={-2,-1,0,1,2}, yticklabels={$-2$,$-1$,$0$,$1$,$2$},
    domain=0.02:5, samples=200,
    clip=false,
]
% Asymptote verticale
\draw[popmuted, dashed, thick] (axis cs:0,-3) -- (axis cs:0,2.5) node[above, pos=0.9, popmuted] {$x=0$};
% Courbe ln(x)
\addplot[popcurve, thick] {ln(x)};
% Tangente en (1,0) : y = x - 1
\addplot[poporange, dashed, domain=-0.5:3] {x-1} node[right, pos=0.9, poporange] {$y=x-1$};
% Points remarquables
\addplot[only marks, mark=*, popblue, mark size=2.5] coordinates {(1,0) (2.718,1)};
\node[above right, popblue] at (axis cs:1,0) {$(1\,;\,0)$};
\node[above right, popblue] at (axis cs:2.718,1) {$(e\,;\,1)$};
% Flèche croissance lente
\draw[popgreen, ->, thick] (axis cs:4.5, 1.4) -- (axis cs:4.5, 1.6) node[above] {Croissance lente};
\end{axis}
\end{tikzpicture}
\legende{Courbe représentative $y = \ln x$ : passage par $(1;0)$ et $(e;1)$, tangente en $1$ d'équation $y=x-1$ (coefficient directeur $1$), asymptote verticale $x=0$.}
\end{popfigure}

\courssub{Propriétés algébriques du logarithme népérien}

Ces propriétés transforment des produits en sommes, des quotients en différences et des puissances en produits. Elles sont la raison historique de l'invention des logarithmes (calculs astronomiques).

\begin{propriete}[Relation fonctionnelle fondamentale]
Pour tous $a > 0$ et $b > 0$ :
\[
\ln(ab) = \ln a + \ln b.
\]
\end{propriete}

\begin{experience}[Démonstration guidée de $\ln(ab)=\ln a+\ln b$]
Fixons $a > 0$. Considérons la fonction $f(x) = \ln(ax) - \ln x$ définie sur $]0\,;+\infty[$.
\begin{enumerate}
    \item $f$ est dérivable comme différence de fonctions dérivables.
    \item $f'(x) = \frac{a}{ax} - \frac{1}{x} = \frac{1}{x} - \frac{1}{x} = 0$.
    \item $f$ est donc constante sur $]0\,;+\infty[$.
    \item Calculons cette constante pour $x=1$ : $f(1) = \ln(a \times 1) - \ln 1 = \ln a - 0 = \ln a$.
    \item Ainsi $\forall x > 0,\; \ln(ax) - \ln x = \ln a$, soit $\ln(ax) = \ln a + \ln x$.
    \item En posant $x = b$, on obtient $\ln(ab) = \ln a + \ln b$.
\end{enumerate}
\end{experience}

\begin{propriete}[Autres propriétés algébriques (conséquences directes)]
Pour $a>0$, $b>0$, $n \in \mathbb{Z}$ (et $q \in \mathbb{N}^*$ pour les racines) :
\begin{align*}
    \ln\left(\frac{a}{b}\right) &= \ln a - \ln b, \\
    \ln(a^n) &= n \ln a, \\
    \ln(\sqrt[q]{a}) &= \frac{1}{q}\ln a.
\end{align*}
\end{propriete}

\begin{attention}[Piège classique : $\ln(a+b) \neq \ln a + \ln b$]
\textbf{Jamais} on n'écrit $\ln(a+b) = \ln a + \ln b$. C'est \textbf{FAUX}.
Exemple : $\ln(2+3) = \ln 5 \approx 1,61$ alors que $\ln 2 + \ln 3 \approx 0,69 + 1,10 = 1,79$.
Le logarithme transforme le \textbf{produit} en somme, pas la somme en somme.
De même : $\ln(a-b) \neq \ln a - \ln b$.
\end{attention}

\begin{exemplebox}[Simplification d'expressions logarithmiques]
Simplifier $A = \ln\left(\frac{e^3 \sqrt{e}}{e^2}\right)$.
\begin{align*}
A &= \ln(e^3) + \ln(\sqrt{e}) - \ln(e^2) \quad\text{(propriétés quotient et produit)}\\
  &= 3\ln e + \frac{1}{2}\ln e - 2\ln e \quad\text{(propriété puissance et racine)}\\
  &= 3 \times 1 + \frac{1}{2} \times 1 - 2 \times 1 \quad\text{($\ln e = 1$)}\\
  &= 1,5 = \frac{3}{2}.
\end{align*}
\end{exemplebox}

\courssub{Dérivée de la composée $\ln(u)$}

C'est un outil technique majeur pour l'étude de fonctions composées faisant intervenir $\ln$.

\begin{propriete}[Dérivée de $\ln(u)$]
Soit $u$ une fonction dérivable et strictement positive sur un intervalle $I$.
La fonction $\ln \circ u$ est dérivable sur $I$ et :
\[
(\ln \circ u)' = \frac{u'}{u}.
\]
\end{propriete}

\begin{methode}[Dériver $\ln(u(x))$]
\begin{enumerate}
    \item Vérifier que $u(x) > 0$ sur l'intervalle d'étude.
    \item Calculer $u'(x)$.
    \item Écrire $(\ln u)' = \dfrac{u'}{u}$.
    \item Simplifier l'expression obtenue.
\end{enumerate}
\end{methode}

\begin{exoresolu}[Dérivée de fonctions composées]
Dériver les fonctions suivantes sur leur ensemble de définition.
\begin{enumerate}
    \item $f(x) = \ln(3x^2 + 1)$ sur $\mathbb{R}$.
    \item $g(x) = \ln\left(\frac{x-1}{x+2}\right)$.
\end{enumerate}
\tcblower
\textbf{1.} $u(x) = 3x^2+1 > 0$ sur $\mathbb{R}$. $u'(x) = 6x$.
$f'(x) = \dfrac{6x}{3x^2+1}$.

\textbf{2.} On détermine d'abord l'ensemble de définition $D_g$ :
$\frac{x-1}{x+2} > 0 \iff (x-1)(x+2) > 0 \iff x \in ]-\infty\,;-2[ \cup ]1\,;+\infty[$.
\emph{Méthode 1 (décomposition) :} Sur $]1\,;+\infty[$, on peut écrire $g(x) = \ln(x-1) - \ln(x+2)$.
$g'(x) = \dfrac{1}{x-1} - \dfrac{1}{x+2} = \dfrac{(x+2)-(x-1)}{(x-1)(x+2)} = \dfrac{3}{(x-1)(x+2)}$.
\emph{Méthode 2 (formule directe, valable sur tout $D_g$) :} $u(x)=\frac{x-1}{x+2}$, $u'(x)=\frac{3}{(x+2)^2}$.
$g'(x)=\frac{u'}{u}=\frac{3}{(x+2)^2}\cdot\frac{x+2}{x-1}=\frac{3}{(x-1)(x+2)}$.
Cette expression est valable sur $]-\infty\,;-2[ \cup ]1\,;+\infty[$.
\end{exoresolu}

\courssub{Étude de la fonction $x \mapsto \ln(ax+b)$}

Cette étude est un classique du Probatoire : elle mobilise domaine, dérivée, variations, limites et courbe.

\begin{methode}[Étude complète de $f(x) = \ln(ax+b)$ avec $a \neq 0$]
\begin{enumerate}
    \item \textbf{Domaine de définition :} Résoudre $ax+b > 0$.
    \begin{itemize}
        \item Si $a > 0$ : $D_f = ]-\frac{b}{a}\,; +\infty[$.
        \item Si $a < 0$ : $D_f = ]-\infty\,; -\frac{b}{a}[$.
    \end{itemize}
    \item \textbf{Dérivée :} $f'(x) = \dfrac{a}{ax+b}$.
    \item \textbf{Sens de variation :} Le signe de $f'$ est le signe de $a$ (car $ax+b > 0$ sur $D_f$).
    \begin{itemize}
        \item $a > 0 \Rightarrow f' > 0 \Rightarrow f$ strictement croissante.
        \item $a < 0 \Rightarrow f' < 0 \Rightarrow f$ strictement décroissante.
    \end{itemize}
    \item \textbf{Limites aux bornes du domaine :}
    \begin{itemize}
        \item Si $a > 0$ : $\lim_{x \to (-\frac{b}{a})^+} f(x) = -\infty$ (asymptote verticale $x = -\frac{b}{a}$), $\lim_{x \to +\infty} f(x) = +\infty$.
        \item Si $a < 0$ : $\lim_{x \to -\infty} f(x) = +\infty$, $\lim_{x \to (-\frac{b}{a})^-} f(x) = -\infty$ (asymptote verticale $x = -\frac{b}{a}$).
    \end{itemize}
    \item \textbf{Tableau de variations et courbe.}
\end{enumerate}
\end{methode}

\begin{exoresolu}[Étude de $f(x) = \ln(2x-3)$]
\begin{enumerate}
    \item Déterminer $D_f$, $f'(x)$, le tableau de variations.
    \item Tracer la courbe $\mathcal{C}_f$.
\end{enumerate}
\tcblower
\textbf{1. Domaine :} $2x-3 > 0 \iff x > \frac{3}{2}$. Donc $D_f = ]\frac{3}{2}\,; +\infty[$.
\textbf{Dérivée :} $f'(x) = \dfrac{2}{2x-3}$. Sur $D_f$, $2x-3 > 0$, donc $f'(x) > 0$.
\textbf{Variations :} $f$ est strictement croissante sur $]\frac{3}{2}\,; +\infty[$.
\textbf{Limites :}
$\lim_{x \to (\frac{3}{2})^+} \ln(2x-3) = \ln(0^+) = -\infty$ (asymptote verticale $x = \frac{3}{2}$).
$\lim_{x \to +\infty} \ln(2x-3) = +\infty$.
\textbf{Tableau de variations :}
\begin{center}
\begin{tabular}{|c|ccc|}\hline
$x$ & $\frac{3}{2}$ & & $+\infty$ \\ \hline
$f'(x)$ & & $+$ & \\ \hline
$f$ & $-\infty$ & $\nearrow$ & $+\infty$ \\ \hline
\end{tabular}
\end{center}

\textbf{2.} La courbe passe par le point d'abscisse $x=2$ (car $2\times2-3=1$, $\ln 1=0$) : point $(2\,;\,0)$.
Tangente en $x=2$ : $f'(2)=2$, équation $y = 2(x-2)$.
\end{exoresolu}

\courssub{Comparaison de croissances et limite de référence}

La fonction $\ln x$ tend vers $+\infty$ quand $x \to +\infty$, mais elle le fait \emph{extrêmement lentement} : plus lentement que toute puissance $x^\alpha$ ($\alpha>0$).

\begin{aretenir}[Limite de référence (admise au Probatoire)]
\[
\lim_{x \to +\infty} \frac{\ln x}{x} = 0.
\]
Plus généralement, pour tout $\alpha > 0$ : $\lim_{x \to +\infty} \frac{\ln x}{x^\alpha} = 0$.
On dit que « $x^\alpha$ domine $\ln x$ à l'infini ».
\end{aretenir}

\begin{popfigure}
\begin{tikzpicture}
\begin{axis}[
    width=\linewidth, height=7cm,
    axis lines=middle,
    xlabel=$x$, ylabel=$y$,
    xmin=0, xmax=10,
    ymin=0, ymax=4,
    xtick={1,2,5,10}, ytick={1,2,3},
    domain=0.1:10, samples=200,
    legend pos=north west,
]
\addplot[popblue, thick] {ln(x)}; \addlegendentry{$y=\ln x$}
\addplot[poporange, thick] {sqrt(max(0,x))}; \addlegendentry{$y=\sqrt{x}$}
\addplot[popgreen, thick] {x/3}; \addlegendentry{$y=x/3$}
\node[popblue, anchor=south east] at (axis cs:9, 2.2) {$\ln x$};
\node[poporange, anchor=south west] at (axis cs:9, 3) {$\sqrt{x}$};
\node[popgreen, anchor=south west] at (axis cs:9, 3.3) {$x/3$};
\end{axis}
\end{tikzpicture}
\legende{Comparaison de croissances : $\ln x$ est dominée par $\sqrt{x}$, elle-même dominée par $x$. Pour $x$ grand, $\ln x$ est négligeable devant $x$.}
\end{popfigure}

\begin{savaistu}[Histoire et applications : du calcul astronomique au pH]
John \textbf{Neper} (1550--1617), Écossais, publie en 1614 ses \emph{Mirifici Logarithmorum Canonis Descriptio}. Son but : transformer les multiplications fastidieuses des astronomes (calculs de positions planétaires) en additions simples grâce à $\ln(ab)=\ln a+\ln b$. Les tables de logarithmes ont été l'outil principal des scientifiques et ingénieurs pendant 350 ans, jusqu'à la calculatrice électronique (années 1970).
Aujourd'hui, l'échelle logarithmique est partout :
\begin{itemize}
    \item \textbf{pH (chimie) :} $\mathrm{pH} = -\log_{10}[\mathrm{H}^+] = -\frac{\ln[\mathrm{H}^+]}{\ln 10}$. Une variation de 1 unité de pH correspond à un facteur 10 en concentration.
    \item \textbf{Échelle de Richter (sismologie) :} Magnitude $M = \log_{10}(A/A_0)$. Un séisme de magnitude 6 libère $\approx 32$ fois plus d'énergie qu'un séisme de magnitude 5.
    \item \textbf{Décibels (acoustique) :} $L = 10 \log_{10}(I/I_0)$.
\end{itemize}
Le logarithme népérien ($\ln$) est le logarithme « naturel » en mathématiques pures car sa dérivée est $1/x$ (coefficient 1), ce qui simplifie toutes les formules du calcul différentiel et intégral.
\end{savaistu}

\courssub{Résolution d'équations logarithmiques simples}

L'injectivité (strictement croissante) de $\ln$ sur $]0\,;+\infty[$ permet de résoudre les équations de la forme $\ln u = \ln v$.

\begin{methode}[Résoudre $\ln u = \ln v$]
\begin{enumerate}
    \item \textbf{Conditions d'existence (C.E.) :} Imposer $u > 0$ et $v > 0$. Noter l'ensemble $D$.
    \item \textbf{Équation :} Sur $D$, $\ln$ est injective, donc $\ln u = \ln v \iff u = v$.
    \item \textbf{Résoudre} $u = v$.
    \item \textbf{Sélectionner} les solutions appartenant à $D$.
    \item \textbf{Conclure} par l'ensemble des solutions $\mathcal{S}$.
\end{enumerate}
\end{methode}

\begin{exoresolu}[Équation logarithmique]
Résoudre dans $\mathbb{R}$ : $\ln(2x-1) = \ln(3-x)$.
\tcblower
\textbf{1. Conditions d'existence :}
$\begin{cases} 2x-1 > 0 \\ 3-x > 0 \end{cases} \iff \begin{cases} x > \frac{1}{2} \\ x < 3 \end{cases} \iff x \in ]\frac{1}{2}\,;\,3[$. C'est l'ensemble $D$.

\textbf{2. Équation :} Sur $D$, $\ln$ est injective, donc :
$2x - 1 = 3 - x \iff 3x = 4 \iff x = \frac{4}{3}$.

\textbf{3. Vérification :} $\frac{4}{3} \approx 1,33$ appartient bien à $]\frac{1}{2}\,;\,3[$.

\textbf{4. Conclusion :} $\mathcal{S} = \left\{\frac{4}{3}\right\}$.
\end{exoresolu}

\begin{exercice}[Entraînement]
\begin{enumerate}
    \item Calculer sans calculatrice : $\ln(e^2) - 2\ln(\sqrt{e}) + \ln(1/e)$.
    \item Déterminer la dérivée de $h(x) = \ln\left(\frac{x^2+1}{x}\right)$ sur $]0\,;+\infty[$.
    \item Résoudre : $\ln(x^2-3) = \ln(2x)$.
    \item Étudier la limite de $f(x) = \frac{\ln x}{\sqrt{x}}$ en $+\infty$ et en $0^+$.
    \item Étudier la fonction $f(x) = \ln(4-2x)$ (domaine, dérivée, variations, limites, tableau, courbe).
\end{enumerate}
\end{exercice}

\begin{corrige}[Exercice n°1]
\begin{enumerate}
    \item $\ln(e^2) - 2\ln(e^{1/2}) + \ln(e^{-1}) = 2 - 2(\frac{1}{2}) + (-1) = 2 - 1 - 1 = 0$.
    \item $h(x) = \ln(x^2+1) - \ln x$. $h'(x) = \frac{2x}{x^2+1} - \frac{1}{x} = \frac{2x^2 - (x^2+1)}{x(x^2+1)} = \frac{x^2-1}{x(x^2+1)}$.
    \item C.E. : $x^2-3>0 \iff x<-\sqrt{3}$ ou $x>\sqrt{3}$ ET $2x>0 \iff x>0$. Donc $D = ]\sqrt{3}\,;+\infty[$.
    Éq : $x^2-3 = 2x \iff x^2-2x-3=0 \iff (x-3)(x+1)=0 \iff x=3$ ou $x=-1$.
    Seul $x=3 \in D$. $\mathcal{S} = \{3\}$.
    \item $\lim_{x\to+\infty} \frac{\ln x}{\sqrt{x}} = 0$ (croissance : $\sqrt{x}$ domine $\ln x$).
    $\lim_{x\to 0^+} \frac{\ln x}{\sqrt{x}} = \frac{-\infty}{0^+} = -\infty$.
    \item $f(x) = \ln(4-2x)$. Domaine : $4-2x > 0 \iff x < 2$. $D_f = ]-\infty\,; 2[$.
    $f'(x) = \frac{-2}{4-2x} = \frac{-1}{2-x}$. Sur $D_f$, $2-x > 0$, donc $f'(x) < 0$.
    $f$ strictement décroissante sur $]-\infty\,; 2[$.
    Limites : $\lim_{x\to -\infty} \ln(4-2x) = +\infty$ ; $\lim_{x\to 2^-} \ln(4-2x) = \ln(0^+) = -\infty$ (asymptote verticale $x=2$).
    Point remarquable : $x= \frac{3}{2} \Rightarrow f(\frac{3}{2}) = \ln 1 = 0$ ; point $(\frac{3}{2}\,;\,0)$.
    Tangente en $\frac{3}{2}$ : $f'(\frac{3}{2}) = -1$, équation $y = -(x-\frac{3}{2})$.
\end{enumerate}
\end{corrige}

\coursec{La fonction exponentielle népérienne}

Dans la section précédente, nous avons découvert la fonction \cle{logarithme népérien} : à tout nombre strictement positif $x$, elle associe un unique réel $\ln x$, et elle transforme les produits en sommes. Mais on peut se poser la question inverse : étant donné un réel $y$, existe-t-il un nombre dont le logarithme vaut $y$ ? Par exemple, quel nombre $x$ vérifie $\ln x = 2$ ? Comme $\ln$ est strictement croissante et continue sur $]0\,;\,+\infty[$, avec des limites $-\infty$ en $0$ et $+\infty$ en $+\infty$, elle atteint \emph{une seule fois} chaque valeur réelle. Il existe donc un unique $x$ tel que $\ln x = 2$ : on le note $x = e^{2}$. Cette « remontée » du logarithme vers le nombre s'appelle la \cle{fonction exponentielle népérienne}. C'est elle qui modélise les croissances rapides : intérêts composés, désintégration radioactive, charge d'un condensateur, croissance d'une population.

\courssub{Définition : la réciproque du logarithme}

\begin{definition}[Fonction exponentielle népérienne]
La fonction logarithme népérien est une bijection de $]0\,;\,+\infty[$ sur $\mathbb{R}$. Sa \cle{fonction réciproque}, notée $\exp$ ou $x \longmapsto e^{x}$, est appelée \textbf{fonction exponentielle népérienne}. Pour tout réel $x$ et tout réel $y > 0$ :
\[ y = e^{x} \iff x = \ln y. \]
\end{definition}

Autrement dit, $e^{x}$ est \emph{le nombre dont le logarithme vaut $x$}. Le nombre $e$ est la base : $e = e^{1}$, c'est le nombre dont le logarithme vaut $1$, soit $e \approx \num{2,718}$.

\begin{propriete}[Relations fondamentales entre ln et exp]
\begin{itemize}
\item Pour tout réel $a > 0$ : $e^{\ln a} = a$.
\item Pour tout réel $x$ : $\ln\left(e^{x}\right) = x$.
\end{itemize}
Ces deux égalités traduisent exactement le fait que $\ln$ et $\exp$ sont réciproques l'une de l'autre : chacune « annule » l'effet de l'autre.
\end{propriete}

\begin{exemplebox}[Lecture directe]
$e^{0} = 1$ car $\ln 1 = 0$ ; $e^{1} = e \approx \num{2,718}$ car $\ln e = 1$ ; $e^{\ln 5} = 5$ ; $\ln\left(e^{-3}\right) = -3$.
\end{exemplebox}

\begin{propriete}[Domaine, image et signe]
La fonction $\exp$ est définie sur $\mathbb{R}$ tout entier, et à valeurs dans $]0\,;\,+\infty[$ :
\[ \text{pour tout réel } x, \quad e^{x} > 0. \]
\end{propriete}

En effet, $e^{x}$ est par définition un nombre de la forme « antécédent par ln », donc un élément de $]0\,;\,+\infty[$.

\begin{attention}[$e^{x}$ ne s'annule jamais]
Erreur très fréquente au Baccalauréat technique : chercher à résoudre $e^{x} = 0$ ou affirmer que $e^{x} < 0$ pour $x$ négatif. C'est \textbf{faux} : pour tout réel $x$, $e^{x} > 0$. Quand $x$ est très négatif, $e^{x}$ est très proche de $0$, mais ne l'atteint jamais. Conséquence pratique : dans une équation produit du type $e^{x}\,(x - 3) = 0$, le facteur $e^{x}$ ne peut pas être nul, il reste seulement $x = 3$.
\end{attention}

\courssub{Dérivée et sens de variation}

Voici la propriété qui rend la fonction exponentielle si précieuse en physique, en électricité et en économie : elle est \emph{sa propre dérivée}. Le taux de croissance de $e^{x}$ en un point est égal à sa valeur en ce point : plus la quantité est grande, plus elle croît vite. C'est exactement le comportement d'un capital placé à intérêts composés.

\begin{propriete}[Dérivabilité de exp]
La fonction $\exp$ est dérivable sur $\mathbb{R}$ et, pour tout réel $x$ :
\[ \left(e^{x}\right)' = e^{x}. \]
C'est l'unique fonction dérivable sur $\mathbb{R}$ égale à sa dérivée et valant $1$ en $0$ (résultat admis).
\end{propriete}

\begin{experience}[Démonstration par la dérivation des fonctions réciproques]
Admettons que $\exp$ soit dérivable sur $\mathbb{R}$ et retrouvons sa dérivée, en utilisant le fait que $\ln$ et $\exp$ sont réciproques.
\begin{enumerate}
\item Partons de l'identité valable pour tout réel $x$ : $\ln\left(e^{x}\right) = x$.
\item Dérivons les deux membres. À gauche, on dérive une composée $\ln \circ\, \exp$ : la dérivée de $\ln(u)$ est $\dfrac{u'}{u}$, donc
\[ \frac{\left(e^{x}\right)'}{e^{x}} = 1. \]
\item Comme $e^{x} > 0$, on peut multiplier les deux membres par $e^{x}$ :
\[ \left(e^{x}\right)' = e^{x}. \]
\end{enumerate}
On retrouve bien que l'exponentielle est sa propre dérivée.
\end{experience}

\begin{propriete}[Sens de variation]
La fonction $\exp$ est \textbf{strictement croissante} sur $\mathbb{R}$.
\end{propriete}

En effet, sa dérivée est $e^{x}$, qui est strictement positive pour tout $x$ : la fonction est donc strictement croissante. Conséquence immédiate et très utile :
\[ e^{a} = e^{b} \iff a = b \qquad \text{et} \qquad e^{a} < e^{b} \iff a < b. \]

\begin{propriete}[Tableau de variations de la fonction exponentielle]
\begin{center}
\begin{tabular}{|c|ccc|}\hline
$x$ & $-\infty$ & & $+\infty$ \\ \hline
$f'(x)$ & & $+$ & \\ \hline
$f(x)=e^x$ & $0$ & $\nearrow$ & $+\infty$ \\ \hline
\end{tabular}
\end{center}
La droite d'équation $y=0$ (axe des abscisses) est asymptote horizontale en $-\infty$.
\end{propriete}

\begin{propriete}[Dérivée de la composée $e^{u}$]
Si $u$ est une fonction dérivable sur un intervalle $I$, alors la fonction $e^{u}$ est dérivable sur $I$ et :
\[ \left(e^{u}\right)' = u' \times e^{u}. \]
\end{propriete}

\begin{exemplebox}[Dérivées de composées]
\begin{itemize}
\item Si $f(x) = e^{3x+1}$, alors $u(x) = 3x+1$, $u'(x) = 3$, donc $f'(x) = 3e^{3x+1}$.
\item Si $g(x) = e^{-x}$, alors $g'(x) = -e^{-x}$.
\item Si $h(x) = e^{x^{2}}$, alors $h'(x) = 2x\,e^{x^{2}}$.
\end{itemize}
\end{exemplebox}

\courssub{Primitives}

\begin{propriete}[Primitives de l'exponentielle et de ses composées affines]
\begin{itemize}
\item Sur $\mathbb{R}$, les primitives de $x \mapsto e^{x}$ sont les fonctions $x \mapsto e^{x} + C$, $C \in \mathbb{R}$.
\item Pour tout réel $a \neq 0$ et tout réel $b$, sur $\mathbb{R}$, les primitives de $x \mapsto e^{ax+b}$ sont les fonctions $x \mapsto \dfrac{1}{a}e^{ax+b} + C$, $C \in \mathbb{R}$.
\end{itemize}
\end{propriete}

\begin{attention}[Condition sur $a$]
La formule $\int e^{ax+b}\,dx = \frac{1}{a}e^{ax+b}+C$ n'est valable que si $a \neq 0$. Si $a=0$, l'intégrale est $\int e^{b}\,dx = xe^{b}+C$.
\end{attention}

\courssub{Limites et croissances comparées}

\begin{propriete}[Limites de exp aux bornes]
\[ \lim_{x \to -\infty} e^{x} = 0 \qquad \text{et} \qquad \lim_{x \to +\infty} e^{x} = +\infty. \]
La première limite signifie que la droite d'équation $y = 0$ (l'axe des abscisses) est \cle{asymptote horizontale} à la courbe de $\exp$ en $-\infty$.
\end{propriete}

Ces limites se déduisent de celles de $\ln$ par le passage à la réciproque : comme $\ln y \to -\infty$ quand $y \to 0^{+}$, la valeur $e^{x}$ tend vers $0$ quand $x \to -\infty$.

\begin{propriete}[Croissances comparées — limite de référence, admise]
\[ \lim_{x \to +\infty} \frac{e^{x}}{x} = +\infty. \]
Autrement dit, en $+\infty$, l'exponentielle « l'emporte » sur toute puissance de $x$ : elle croît infiniment plus vite que $x$.
\end{propriete}

\begin{exoresolu}[Calcul de limite avec croissances comparées]
Déterminer $\displaystyle\lim_{x \to +\infty} \left(e^{x} - 3x\right)$.
\tcblower
On met $e^{x}$ en facteur, car on est face à une forme indéterminée « $+\infty - \infty$ » :
\[ e^{x} - 3x = e^{x}\left(1 - \frac{3x}{e^{x}}\right) = e^{x}\left(1 - \frac{3}{\dfrac{e^{x}}{x}}\right). \]
Or $\dfrac{e^{x}}{x} \to +\infty$, donc $\dfrac{3}{e^{x}/x} \to 0$ et la parenthèse tend vers $1$. Par produit :
\[ \lim_{x \to +\infty} \left(e^{x} - 3x\right) = +\infty. \]
\end{exoresolu}

\courssub{Propriétés algébriques}

L'exponentielle hérite des propriétés du logarithme, « à l'envers » : là où $\ln$ transforme les produits en sommes, $\exp$ transforme les sommes en produits.

\begin{propriete}[Relation fonctionnelle et conséquences]
Pour tous réels $a$ et $b$, et pour tout entier relatif $n$ :
\[ e^{a+b} = e^{a} \times e^{b}, \qquad e^{a-b} = \frac{e^{a}}{e^{b}}, \qquad e^{-a} = \frac{1}{e^{a}}, \qquad \left(e^{a}\right)^{n} = e^{na}. \]
\end{propriete}

\begin{experience}[Démonstration de la relation fondamentale]
Montrons que $e^{a+b} = e^{a} \times e^{b}$. Posons $u = e^{a}$ et $v = e^{b}$ (ce sont des réels strictement positifs).
\begin{enumerate}
\item Par définition de la réciproque : $a = \ln u$ et $b = \ln v$.
\item Donc $a + b = \ln u + \ln v = \ln(uv)$, d'après la propriété fondamentale du logarithme.
\item En repassant par l'exponentielle : $e^{a+b} = e^{\ln(uv)} = uv = e^{a} \times e^{b}$.
\end{enumerate}
Les autres formules s'en déduisent : $e^{a-b} = e^{a} \times e^{-b} = \dfrac{e^{a}}{e^{b}}$, et $\left(e^{a}\right)^{n} = e^{a} \times \cdots \times e^{a} = e^{na}$.
\end{experience}

\begin{exemplebox}[Simplifications]
$e^{2} \times e^{3} = e^{5}$ ; $\dfrac{e^{7}}{e^{4}} = e^{3}$ ; $\left(e^{2}\right)^{3} = e^{6}$ ; $e^{-2} = \dfrac{1}{e^{2}} \approx \num{0,135}$.
\end{exemplebox}

\begin{attention}[Pièges de calcul]
\begin{itemize}
\item $e^{a} + e^{b}$ ne se simplifie \textbf{pas} : en particulier $e^{a} + e^{b} \neq e^{a+b}$. Par exemple $e^{1} + e^{1} = 2e \approx \num{5,44}$, alors que $e^{2} \approx \num{7,39}$.
\item $e^{a} \times e^{b} = e^{a+b}$, et non $e^{ab}$.
\item $\left(e^{a}\right)^{n} = e^{na}$, et non $e^{a^{n}}$.
\end{itemize}
\end{attention}

\courssub{Courbe représentative et tangente remarquable}

Puisque $\exp$ et $\ln$ sont réciproques, leurs courbes sont \textbf{symétriques par rapport à la droite d'équation $y = x$} (la première bissectrice) : le point $(a\,;\,b)$ est sur la courbe de $\exp$ si et seulement si $(b\,;\,a)$ est sur celle de $\ln$.

\begin{popfigure}
\begin{center}
\begin{tikzpicture}
\begin{axis}[width=0.85\linewidth, axis lines=middle, xlabel={$x$}, ylabel={$y$}, xmin=-3.5, xmax=3.5, ymin=-3.5, ymax=5.5, xtick={-3,-2,-1,1,2,3}, ytick={-3,-2,-1,1,2,3,4,5}, legend style={at={(0.03,0.97)}, anchor=north west, font=\small}, samples=200]
\addplot[popcurve, domain=-3:1.6]{exp(x)};
\addlegendentry{$y = e^{x}$}
\addplot[poporange, thick, domain=0.03:3.5]{ln(x)};
\addlegendentry{$y = \ln x$}
\addplot[popmuted, dashed, domain=-3:3.5]{x};
\addlegendentry{$y = x$}
\end{axis}
\end{tikzpicture}
\end{center}
\legende{Les courbes de $\exp$ et de $\ln$ sont symétriques par rapport à la droite $y = x$.}
\end{popfigure}

Comme $\left(e^{x}\right)' = e^{x}$, le nombre dérivé en $0$ vaut $e^{0} = 1$. La tangente à la courbe au point $(0\,;\,1)$ a donc pour équation :
\[ y = e^{0}(x - 0) + e^{0} = x + 1. \]

\begin{popfigure}
\begin{center}
\begin{tikzpicture}
\begin{axis}[width=0.85\linewidth, axis lines=middle, xlabel={$x$}, ylabel={$y$}, xmin=-3, xmax=2.5, ymin=-0.5, ymax=6, xtick={-3,-2,-1,1,2}, ytick={1,2,3,4,5}, samples=200, legend style={at={(0.03,0.97)}, anchor=north west, font=\small}]
\addplot[popcurve, domain=-3:1.8]{exp(x)};
\addlegendentry{$y = e^{x}$}
\addplot[poppink, thick, domain=-2.5:2.2]{x+1};
\addlegendentry{Tangente : $y = x + 1$}
\addplot[only marks, mark=*, popdark] coordinates {(0,1)};
\end{axis}
\end{tikzpicture}
\end{center}
\legende{La courbe de $\exp$ et sa tangente au point $(0\,;\,1)$, d'équation $y = x+1$. L'axe des abscisses est asymptote en $-\infty$.}
\end{popfigure}

\begin{methode}[Étudier le signe de $e^{u} - 1$ ou comparer $e^{u}$ à $1$]
Comme $\exp$ est strictement croissante et que $e^{0} = 1$ :
\[ e^{u} \geq 1 \iff u \geq 0 \qquad \text{et} \qquad e^{u} \leq 1 \iff u \leq 0. \]
Pour étudier le signe de $e^{u(x)} - 1$ :
\begin{enumerate}
\item On écrit $e^{u(x)} - 1 = e^{u(x)} - e^{0}$.
\item Par stricte croissance de $\exp$, cette quantité a le même signe que $u(x) - 0 = u(x)$.
\item On étudie donc simplement le signe de $u(x)$.
\end{enumerate}
Exemple : $e^{2x-4} - 1 \geq 0 \iff 2x - 4 \geq 0 \iff x \geq 2$.
\end{methode}

\begin{exoresolu}[Étude complète d'une fonction exponentielle]
Soit $f$ définie sur $\mathbb{R}$ par $f(x) = e^{2x-4} - 1$.
\begin{enumerate}
\item Calculer $f'(x)$ et dresser le tableau de variations de $f$.
\item Déterminer les limites de $f$ aux bornes de $\mathbb{R}$.
\item Déterminer le signe de $f(x)$.
\end{enumerate}
\tcblower
\textbf{1.} $f$ est dérivable sur $\mathbb{R}$ comme composée de fonctions dérivables. $u(x)=2x-4$, $u'(x)=2$.
\[ f'(x) = u'(x) e^{u(x)} = 2e^{2x-4}. \]
Comme $e^{2x-4} > 0$ pour tout $x$ et que $2 > 0$, on a $f'(x) > 0$ sur $\mathbb{R}$ : $f$ est strictement croissante sur $\mathbb{R}$.

\textbf{Tableau de variations :}
\begin{center}
\begin{tabular}{|c|ccc|}\hline
$x$ & $-\infty$ & & $+\infty$ \\ \hline
$f'(x)$ & & $+$ & \\ \hline
$f(x)$ & $-1$ & $\nearrow$ & $+\infty$ \\ \hline
\end{tabular}
\end{center}

\textbf{2.} Limites aux bornes :
\begin{itemize}
\item Quand $x \to -\infty$, $2x-4 \to -\infty$, donc $e^{2x-4} \to 0$, ainsi $f(x) \to -1$.
\item Quand $x \to +\infty$, $2x-4 \to +\infty$, donc $e^{2x-4} \to +\infty$, ainsi $f(x) \to +\infty$.
\end{itemize}

\textbf{3.} $f(x) \geq 0 \iff e^{2x-4} \geq e^{0} \iff 2x - 4 \geq 0 \iff x \geq 2$.
Donc $f$ est négative sur $]-\infty\,;\,2[$, nulle en $2$, positive sur $]2\,;\,+\infty[$.
\end{exoresolu}

\begin{aretenir}[L'essentiel sur la fonction exponentielle]
\begin{itemize}
\item $y = e^{x} \iff x = \ln y$ ; \quad $e^{\ln a} = a$ ($a > 0$) ; \quad $\ln(e^{x}) = x$.
\item $\exp$ est définie sur $\mathbb{R}$, strictement croissante, et $e^{x} > 0$ pour tout $x$.
\item $\left(e^{x}\right)' = e^{x}$ et $\left(e^{u}\right)' = u'\,e^{u}$.
\item Primitives : $\int e^{x}\,dx = e^{x}+C$ ; $\int e^{ax+b}\,dx = \frac{1}{a}e^{ax+b}+C$ ($a \neq 0$).
\item $\displaystyle\lim_{x \to -\infty} e^{x} = 0$ (asymptote $y = 0$) ; $\displaystyle\lim_{x \to +\infty} e^{x} = +\infty$ ; $\displaystyle\lim_{x \to +\infty} \frac{e^{x}}{x} = +\infty$.
\item $e^{a+b} = e^{a}e^{b}$ ; $e^{-a} = \dfrac{1}{e^{a}}$ ; $\left(e^{a}\right)^{n} = e^{na}$.
\item Les courbes de $\exp$ et $\ln$ sont symétriques par rapport à $y = x$ ; tangente en $(0\,;\,1)$ : $y = x+1$.
\end{itemize}
\end{aretenir}

\begin{savaistu}[Le nombre $e$, des banques aux populations]
Le nombre $e \approx \num{2,71828}$ est né au XVII\textsuperscript{e} siècle dans l'étude des \textbf{intérêts composés}. Si un capital de 1~000~000 FCFA est placé à $100\,\%$ l'an et que les intérêts sont calculés de plus en plus souvent (chaque mois, chaque jour, chaque instant...), le capital au bout d'un an se rapproche de $1\,000\,000 \times e$ FCFA, sans jamais le dépasser. Plus généralement, un capital placé au taux $t$ en composition continue est multiplié par $e^{t}$ chaque année. La même fonction décrit la croissance d'une population de bactéries, la charge d'un condensateur dans un circuit électrique ou le refroidissement d'une pièce mécanique : partout où « la vitesse de croissance est proportionnelle à la quantité présente », l'exponentielle apparaît.
\end{savaistu}

\coursec{Équations et inéquations avec ln et exp}

Dans les sections précédentes, nous avons étudié la fonction exponentielle népérienne $x \mapsto e^x$, bijection strictement croissante de $\mathbb{R}$ sur $]0;+\infty[$, et sa réciproque, le logarithme népérien $x \mapsto \ln x$, bijection strictement croissante de $]0;+\infty[$ sur $\mathbb{R}$. Ces deux fonctions modélisent de nombreux phénomènes techniques : charge et décharge d'un condensateur, décroissance radioactive, intérêts composés, croissance d'une population. Résoudre des équations et inéquations les faisant intervenir repose sur deux idées maîtresses : \textbf{la bijectivité} (qui permet de « faire tomber » le $\ln$ ou l'exponentielle) et \textbf{la stricte croissance} (qui conserve le sens des inégalités).

\courssub{Équivalences fondamentales et conditions d'existence}

Avant toute manipulation, il est impératif de rappeler le domaine de définition du logarithme népérien.

\begin{propriete}[Domaine de définition de $\ln$]
La fonction $\ln$ n'est définie que sur $]0;+\infty[$. Pour toute expression $\ln(u(x))$, la \textbf{condition d'existence (C.E.)} est :
\[ u(x) > 0 \]
Toute solution candidate qui ne satisfait pas cette condition est \textbf{parasite} et doit être rejetée.
\end{propriete}

\begin{propriete}[Équivalences fondamentales]
\begin{itemize}
    \item Pour tous réels $a$ et $b$ \textbf{strictement positifs} :
    \[ \ln a = \ln b \iff a = b \qquad \text{et} \qquad \ln a = k \iff a = e^k \ (k \in \mathbb{R}). \]
    \item Pour tous réels $a$ et $b$ \textbf{quelconques} :
    \[ e^a = e^b \iff a = b \qquad \text{et} \qquad e^a = k \iff a = \ln k \ (\text{avec } k > 0). \]
\end{itemize}
Ces équivalences, conséquences directes de la bijectivité des deux fonctions, sont les \textbf{clés de résolution} : elles ramènent l'équation à une équation algébrique simple.
\end{propriete}

\begin{attention}[Piège classique : l'oubli de la condition d'existence]
L'erreur la plus fréquente en Terminale est de résoudre $\ln(2x-1) = \ln(x-2)$ en écrivant directement $2x-1 = x-2$, ce qui donne $x = -1$. Or $2\times(-1)-1 = -3 < 0$ : $\ln(2x-1)$ n'existe pas pour $x=-1$. L'équation n'a \textbf{aucune solution}.
\textbf{Règle d'or :} on pose \textbf{systématiquement} la C.E. \emph{avant} de transformer l'équation, et on confronte chaque solution trouvée à cette condition.
\end{attention}

\courssub{Résolution d'équations simples : $\ln u = k$ et $e^u = k$}

\begin{methode}[Résoudre une équation de la forme $\ln(u(x)) = k$ ou $e^{u(x)} = k$]
\textbf{Étape 1 : Condition d'existence (C.E.)}
\begin{itemize}
    \item Pour $\ln(u(x)) = k$ : imposer $u(x) > 0$.
    \item Pour $e^{u(x)} = k$ : exiger $k > 0$ (car une exponentielle est toujours strictement positive). Si $k \le 0$, il n'y a \textbf{pas de solution}.
\end{itemize}
\textbf{Étape 2 : Transformation équivalente}
\begin{itemize}
    \item $\ln(u(x)) = k \iff u(x) = e^k$.
    \item $e^{u(x)} = k \iff u(x) = \ln k$ (sachant $k>0$).
\end{itemize}
\textbf{Étape 3 : Résolution algébrique et vérification}
Résoudre l'équation obtenue, puis ne garder que les solutions compatibles avec la C.E.
\end{methode}

\begin{exemplebox}[Résolution de $\ln(3x-2) = 2$ et $e^{2x+1} = 5$]
\underline{Équation 1 :} $\ln(3x-2) = 2$
\begin{enumerate}
    \item \textbf{C.E.} : $3x-2 > 0 \iff x > \frac{2}{3}$.
    \item \textbf{Transformation} : $3x-2 = e^2$.
    \item \textbf{Résolution} : $3x = e^2 + 2 \iff x = \frac{e^2+2}{3}$.
    \item \textbf{Vérification} : $e^2 \approx 7,389$, donc $x \approx 3,13 > \frac{2}{3}$. \textbf{Solution valide}.
\end{enumerate}
$S = \left\{ \dfrac{e^2+2}{3} \right\}$.

\underline{Équation 2 :} $e^{2x+1} = 5$
\begin{enumerate}
    \item \textbf{C.E.} : $5 > 0$, condition satisfaite.
    \item \textbf{Transformation} : $2x+1 = \ln 5$.
    \item \textbf{Résolution} : $2x = \ln 5 - 1 \iff x = \frac{\ln 5 - 1}{2}$.
\end{enumerate}
$S = \left\{ \dfrac{\ln 5 - 1}{2} \right\}$.
\end{exemplebox}

\courssub{Équations faisant intervenir des sommes ou différences de logarithmes}

On utilise les propriétés algébriques du logarithme pour regrouper les termes avant d'appliquer l'équivalence fondamentale.
Rappel : pour $a>0$ et $b>0$, $\ln a + \ln b = \ln(ab)$ et $\ln a - \ln b = \ln\left(\frac{a}{b}\right)$.

\begin{exoresolu}[Résolution de $\ln(x+1) + \ln(x-1) = \ln 8$]
\textbf{Étape 1 : Conditions d'existence}
\begin{align*}
    x+1 &> 0 \iff x > -1 \\
    x-1 &> 0 \iff x > 1
\end{align*}
Les deux conditions doivent être satisfaites simultanément : la C.E. globale est \textbf{$x > 1$}.

\textbf{Étape 2 : Regroupement (valide sur le domaine de définition)}
$\ln(x+1) + \ln(x-1) = \ln\bigl((x+1)(x-1)\bigr) = \ln(x^2-1)$.
L'équation devient : $\ln(x^2-1) = \ln 8$.

\textbf{Étape 3 : Utilisation de l'injectivité}
$x^2 - 1 = 8 \iff x^2 = 9 \iff x = 3 \text{ ou } x = -3$.

\textbf{Étape 4 : Vérification par rapport à la C.E. ($x>1$)}
\begin{itemize}
    \item $x = 3$ : $3 > 1$, \textbf{acceptée}.
    \item $x = -3$ : $-3 \not> 1$, \textbf{rejetée} (solution parasite).
\end{itemize}
\textbf{Solution finale :} $S = \{3\}$.
\end{exoresolu}

\begin{attention}[Attention au regroupement prématuré]
On ne peut écrire $\ln(x+1)+\ln(x-1)=\ln(x^2-1)$ que si $x+1>0$ \textbf{et} $x-1>0$. Sans la C.E., on accepterait $x=-3$ car $(-3)^2-1=8>0$, alors que $\ln(-3+1)=\ln(-2)$ n'existe pas. La C.E. protège contre les solutions parasites.
\end{attention}

\courssub{Équations se ramenant au second degré par changement de variable}

Certaines équations mélangeant $e^x$ et $e^{2x}$ se transforment en équations du second degré, car $e^{2x} = (e^x)^2$.

\begin{methode}[Changement de variable $X = e^x$]
Pour une équation de type $a e^{2x} + b e^x + c = 0$ :
\begin{enumerate}
    \item Poser $X = e^x$. \textbf{Contrainte essentielle :} $X > 0$ (car $e^x > 0$ pour tout réel $x$).
    \item L'équation devient $a X^2 + b X + c = 0$.
    \item Résoudre cette équation du second degré d'inconnue $X$.
    \item Ne conserver que les racines \textbf{strictement positives}.
    \item Revenir à $x$ par $x = \ln X$ pour chaque racine positive conservée.
\end{enumerate}
\end{methode}

\begin{exoresolu}[Résolution de $e^{2x} - 5e^x + 6 = 0$]
\begin{enumerate}
    \item \textbf{Changement de variable :} $X = e^x$, avec $X > 0$.
    \item \textbf{Équation en $X$ :} $X^2 - 5X + 6 = 0$.
    \item \textbf{Résolution :} $\Delta = (-5)^2 - 4\times 1 \times 6 = 25 - 24 = 1 > 0$.
    \[ X_1 = \frac{5-1}{2} = 2 \qquad ; \qquad X_2 = \frac{5+1}{2} = 3. \]
    \item \textbf{Sélection :} $2 > 0$ et $3 > 0$ : les deux racines sont acceptées.
    \item \textbf{Retour à $x$ :}
    \begin{align*}
        e^x = 2 &\iff x = \ln 2 \\
        e^x = 3 &\iff x = \ln 3
    \end{align*}
\end{enumerate}
\textbf{Solution finale :} $S = \{\ln 2 ; \ln 3\}$.
\end{exoresolu}

\begin{popfigure}
\begin{tikzpicture}[scale=0.9]
\begin{axis}[
    width=\linewidth, height=8cm,
    axis lines=middle,
    xlabel=$x$, ylabel=$y$,
    xmin=-1, xmax=2.5,
    ymin=-1, ymax=8,
    xtick={0,0.693,1.099,2}, xticklabels={$0$, $\ln 2$, $\ln 3$, $2$},
    ytick={1,2,3,6}, yticklabels={$1$, $2$, $3$, $6$},
    grid=major, grid style={popline},
    clip=false
]
\addplot[popcurve, domain=-1:2.05, samples=200] {exp(x)};
\node[popdark, anchor=south west] at (axis cs:1.5, 5.2) {$y = e^x$};

\addplot[poporange, dashed, domain=-1:2.5] {3};
\node[poporange, anchor=south west] at (axis cs:1.6, 3.1) {$y = 3$};
\addplot[popgreen, dashed, domain=-1:2.5] {2};
\node[popgreen, anchor=north west] at (axis cs:1.6, 1.9) {$y = 2$};

\addplot[only marks, mark=*, mark size=2.5, popblue] coordinates {(1.0986, 3)};
\addplot[only marks, mark=*, mark size=2.5, popblue] coordinates {(0.6931, 2)};

\draw[popmuted, dashed] (axis cs:1.0986, 0) -- (axis cs:1.0986, 3);
\draw[popmuted, dashed] (axis cs:0.6931, 0) -- (axis cs:0.6931, 2);
\draw[popmuted, dashed] (axis cs:0, 3) -- (axis cs:1.0986, 3);
\draw[popmuted, dashed] (axis cs:0, 2) -- (axis cs:0.6931, 2);
\end{axis}
\end{tikzpicture}
\legende{Lecture graphique : les solutions de $e^x = 2$ et $e^x = 3$ (issues de $e^{2x}-5e^x+6=0$) sont les abscisses $\ln 2$ et $\ln 3$ des points d'intersection.}
\end{popfigure}

\courssub{Inéquations avec $\ln$ et $\exp$}

La résolution des inéquations repose sur la \textbf{stricte croissance} des deux fonctions : elles \emph{conservent} le sens des inégalités.

\begin{propriete}[Ordre et fonctions $\ln$ / $\exp$]
Soient $u(x)$ et $v(x)$ deux expressions réelles.
\begin{itemize}
    \item \textbf{Pour $\ln$ :} sous les conditions $u(x)>0$ et $v(x)>0$,
    \[ \ln u(x) \le \ln v(x) \iff u(x) \le v(x) \qquad \text{et} \qquad \ln u(x) \le k \iff u(x) \le e^k. \]
    \item \textbf{Pour $\exp$ :} pour tous réels $u(x)$, $v(x)$,
    \[ e^{u(x)} \le e^{v(x)} \iff u(x) \le v(x) \qquad \text{et, si } k>0, \qquad e^{u(x)} \le k \iff u(x) \le \ln k. \]
    \item Si $k \le 0$, l'inéquation $e^{u(x)} \le k$ n'a \textbf{aucune solution} ; l'inéquation $e^{u(x)} \ge k$ est vérifiée \textbf{pour tout $x$} du domaine.
\end{itemize}
Les mêmes équivalences valent avec $<$, $\ge$, $>$.
\end{propriete}

\begin{methode}[Résoudre une inéquation avec $\ln$]
\begin{enumerate}
    \item \textbf{Poser les C.E.} : chaque quantité sous un $\ln$ doit être strictement positive.
    \item \textbf{Transformer} en inéquation algébrique équivalente (la croissance conserve le sens).
    \item \textbf{Résoudre} l'inéquation algébrique (tableau de signes si produit ou quotient).
    \item \textbf{Intersecter} l'ensemble des solutions avec l'ensemble des C.E.
\end{enumerate}
\end{methode}

\begin{exemplebox}[Résolution de $\ln(2x-1) \ge 0$]
\underline{Écriture préalable :} $\ln(2x-1) \ge \ln 1$ (car $\ln 1 = 0$).
\begin{enumerate}
    \item \textbf{C.E.} : $2x-1 > 0 \iff x > \frac{1}{2}$.
    \item \textbf{Transformation (croissance de $\ln$)} : $2x-1 \ge 1$.
    \item \textbf{Résolution} : $2x \ge 2 \iff x \ge 1$.
    \item \textbf{Intersection avec la C.E.} : $[1; +\infty[ \cap \left]\frac{1}{2}; +\infty\right[ = [1; +\infty[$.
\end{enumerate}
$S = [1; +\infty[$.
\end{exemplebox}

\begin{popfigure}
\begin{tikzpicture}
\draw[popdark, thick, <->] (-0.5,0) -- (5.5,0);
\foreach \x/\label in {0/0, 0.5/$\frac{1}{2}$, 1/1, 2/2, 3/3, 4/4} {
    \draw (\x, -0.1) -- (\x, 0.1);
    \node[below] at (\x, -0.3) {\label};
}

\draw[popblue, very thick] (0.5, 0.5) -- (5, 0.5);
\draw[popblue, fill=white] (0.5, 0.5) circle (2.2pt);
\node[popblue, above] at (2.75, 0.6) {C.E. : $x > \frac{1}{2}$};

\draw[popgreen, very thick] (1, 1.2) -- (5, 1.2);
\draw[popgreen, fill=popgreen] (1, 1.2) circle (2.2pt);
\node[popgreen, above] at (3, 1.3) {$2x-1 \ge 1 \iff x \ge 1$};

\draw[poporange, ultra thick] (1, -0.9) -- (5, -0.9);
\draw[poporange, fill=poporange] (1, -0.9) circle (2.2pt);
\node[poporange, below] at (3, -1) {\textbf{Solution finale : $S = [1; +\infty[$}};

\draw[popblue, ->] (0.5, 0.45) -- (0.5, 0.15);
\draw[popgreen, ->] (1, 1.15) -- (1, 0.15);
\draw[poporange, ->] (1, -0.85) -- (1, -0.45);
\end{tikzpicture}
\legende{Droite des réels pour $\ln(2x-1) \ge 0$ : la solution est l'intersection de la C.E. ($x>\frac{1}{2}$, borne ouverte) et de l'inéquation transformée ($x\ge 1$, borne fermée).}
\end{popfigure}

\begin{exoresolu}[Résolution de $\dfrac{\ln x}{x-2} \le 0$]
Le quotient impose un \textbf{tableau de signes}.
\textbf{Domaine de définition (C.E.)} : $x > 0$ (pour $\ln x$) et $x \neq 2$ (dénominateur non nul).
Donc $D = ]0; 2[ \cup ]2; +\infty[$.

Signes : $\ln x < 0$ sur $]0;1[$, $\ln 1 = 0$, $\ln x > 0$ sur $]1;+\infty[$ ; $x-2 < 0$ sur $]0;2[$, nul en $2$, positif après.

\begin{center}
\begin{tabular}{|c|ccccccc|}\hline
$x$ & $0$ & & $1$ & & $2$ & & $+\infty$ \\ \hline
$\ln x$ & $||$ & $-$ & $0$ & $+$ & $|$ & $+$ & $+$ \\
$x-2$ & $|$ & $-$ & $-$ & $-$ & $0$ & $+$ & $+$ \\ \hline
$\dfrac{\ln x}{x-2}$ & $||$ & $+$ & $0$ & $-$ & $||$ & $+$ & $+$ \\ \hline
\end{tabular}
\end{center}
On cherche où le quotient est $\le 0$ : il est strictement négatif sur $]1; 2[$ et nul en $x=1$.
\textbf{Solution :} $S = [1; 2[$.
\end{exoresolu}

\courssub{Équations paramétrées : $a e^x + b = 0$}

Cette forme simple permet de discuter le nombre de solutions selon les paramètres, exercice fréquent au Probatoire.

\begin{propriete}[Discussion de $a e^x + b = 0$ ($a \neq 0$)]
L'équation s'écrit $e^x = -\frac{b}{a}$. Comme $e^x > 0$ pour tout réel $x$ :
\begin{itemize}
    \item si $-\frac{b}{a} \le 0$ (c'est-à-dire $ab \ge 0$) : \textbf{aucune solution} ($S = \varnothing$) ;
    \item si $-\frac{b}{a} > 0$ (c'est-à-dire $ab < 0$) : \textbf{une unique solution} $x = \ln\left(-\frac{b}{a}\right)$.
\end{itemize}
\end{propriete}

\begin{exemplebox}[Discussion de $(m-1)e^x + m = 0$ selon les valeurs du réel $m$]
\textbf{Cas particulier :} si $m = 1$, l'équation devient $0 \times e^x + 1 = 0$, impossible : $S = \varnothing$.

\textbf{Cas général ($m \neq 1$) :} $e^x = -\frac{m}{m-1}$. On étudie le signe de $-\frac{m}{m-1}$.

\begin{center}
\begin{tabular}{|c|ccccccc|}\hline
$m$ & $-\infty$ & & $0$ & & $1$ & & $+\infty$ \\ \hline
$m$ & $|$ & $-$ & $0$ & $+$ & $|$ & $+$ & $+$ \\
$m-1$ & $|$ & $-$ & $-$ & $-$ & $0$ & $+$ & $+$ \\ \hline
$-\dfrac{m}{m-1}$ & $|$ & $-$ & $0$ & $+$ & $||$ & $-$ & $-$ \\ \hline
\end{tabular}
\end{center}

\begin{itemize}
    \item Si $m \in ]-\infty; 0] \cup ]1; +\infty[$ : $-\frac{m}{m-1} \le 0$, donc $S = \varnothing$.
    \item Si $m = 1$ : $S = \varnothing$ (vu ci-dessus).
    \item Si $m \in ]0; 1[$ : $-\frac{m}{m-1} > 0$, donc $S = \left\{ \ln\left(-\frac{m}{m-1}\right) \right\}$.
\end{itemize}
\end{exemplebox}

\courssub{Interprétation concrète : temps de doublement et seuils}

Les équations exponentielles modélisent des phénomènes réels : charge d'un condensateur, croissance bactérienne, intérêts composés, décroissance radioactive.

\begin{savaistu}[La règle des 70 — Doublement d'un capital]
Un capital placé à un taux annuel de $t$ \% suit la loi $C_n = C_0 \left(1+\frac{t}{100}\right)^n$.
On cherche $n$ tel que $C_n = 2 C_0$ :
\[ \left(1+\frac{t}{100}\right)^n = 2 \iff n \ln\left(1+\frac{t}{100}\right) = \ln 2. \]
Pour $t$ petit (quelques \%), $\ln\left(1+\frac{t}{100}\right) \approx \frac{t}{100}$, donc $n \approx \frac{100 \ln 2}{t} \approx \frac{69,3}{t} \approx \frac{70}{t}$.
\textbf{Exemple :} à $3,5\%$ l'an, un capital double en environ $\frac{70}{3,5} = 20$ ans.
Cette estimation mentale est très utilisée en gestion, en finance et en biologie (temps de doublement d'une population).
\end{savaistu}

\begin{exoresolu}[Application : seuil de concentration dans un réacteur]
Dans un réacteur chimique, la concentration $C(t)$ (en mol/L) d'un produit suit la loi $C(t) = 10\left(1 - e^{-0,2t}\right)$ pour $t \ge 0$ ($t$ en minutes). Déterminer à partir de quel instant la concentration dépasse $8$ mol/L.
\begin{enumerate}
    \item \textbf{Inéquation :} $10\left(1 - e^{-0,2t}\right) > 8$.
    \item \textbf{Résolution :}
    \begin{align*}
        1 - e^{-0,2t} &> 0,8 \\
        -e^{-0,2t} &> -0,2 \\
        e^{-0,2t} &< 0,2 \quad \text{(division par } -1 \text{ : changement de sens)}
    \end{align*}
    \item \textbf{Application de $\ln$ (strictement croissante)} : $-0,2t < \ln 0,2$.
    \item \textbf{Division par $-0,2$ (négatif : changement de sens)} :
    \[ t > \frac{\ln 0,2}{-0,2} = \frac{-\ln 5}{-0,2} = 5 \ln 5. \]
    \item \textbf{Valeur approchée :} $\ln 5 \approx 1,609$, donc $t > 8,05$ min environ.
\end{enumerate}
\textbf{Conclusion :} la concentration dépasse $8$ mol/L à partir d'environ \textbf{8 minutes 3 secondes}.
\end{exoresolu}

\begin{attention}[Deux changements de sens à ne pas oublier]
Dans la résolution d'inéquations exponentielles du type $e^{-kt} < a$, on change le sens de l'inégalité \textbf{deux fois} : en isolant l'exponentielle (multiplication par $-1$) et en divisant par le coefficient négatif $-k$. Appliquer $\ln$ ne change \textbf{pas} le sens car $\ln$ est croissante. Vérifier toujours la cohérence du résultat avec le contexte (ici, la concentration augmente avec le temps, donc le seuil est bien un « $t >$ ... »).
\end{attention}

\courssub{Exercices d'entraînement}

\begin{exercice}[Équations et inéquations de base]
Résoudre dans $\mathbb{R}$ :
\begin{enumerate}
    \item $\ln(x^2 - 4) = \ln 5$
    \item $e^{2x-1} = 3$
    \item $\ln(2x+1) < 0$
    \item $e^x + 2e^{-x} = 3$ \quad \textit{(indication : multiplier par $e^x$ et poser $X=e^x$)}
\end{enumerate}
\end{exercice}

\begin{exercice}[Problème concret : décroissance radioactive]
L'activité $A(t)$ (en Bq) d'un échantillon radioactif est donnée par $A(t) = A_0 e^{-\lambda t}$ avec $\lambda = 0,01$ min$^{-1}$.
\begin{enumerate}
    \item Exprimer le temps $t_{1/2}$ (période radioactive) pour lequel l'activité est divisée par 2, en fonction de $\lambda$, puis calculer sa valeur.
    \item Au bout de combien de temps l'activité devient-elle inférieure à $10\%$ de l'activité initiale ?
\end{enumerate}
\end{exercice}

\begin{corrige}[Exercice 1]
\begin{enumerate}
    \item \textbf{C.E.} : $x^2-4>0 \iff x<-2$ ou $x>2$. \textbf{Équation} : $x^2-4=5 \iff x^2=9 \iff x=3$ ou $x=-3$. Les deux valeurs vérifient la C.E. ($3>2$ et $-3<-2$). $S=\{-3; 3\}$.
    \item $2x-1 = \ln 3 \iff x = \dfrac{\ln 3 + 1}{2}$. $S = \left\{\dfrac{\ln 3 + 1}{2}\right\}$.
    \item \textbf{C.E.} : $2x+1>0 \iff x>-\frac{1}{2}$. $\ln(2x+1) < \ln 1 \iff 2x+1 < 1 \iff x < 0$. Intersection : $S = \left]-\frac{1}{2}; 0\right[$.
    \item On multiplie par $e^x > 0$ : $e^{2x} + 2 = 3e^x$. On pose $X = e^x$ ($X>0$) : $X^2 - 3X + 2 = 0$, $\Delta = 9-8=1$, racines $X=1$ et $X=2$, toutes deux positives. $e^x=1 \iff x=0$ ; $e^x=2 \iff x=\ln 2$. $S=\{0; \ln 2\}$.
\end{enumerate}
\end{corrige}

\begin{corrige}[Exercice 2]
\begin{enumerate}
    \item $A_0 e^{-\lambda t_{1/2}} = \dfrac{A_0}{2} \iff e^{-\lambda t_{1/2}} = \dfrac{1}{2} \iff -\lambda t_{1/2} = -\ln 2 \iff t_{1/2} = \dfrac{\ln 2}{\lambda} = \dfrac{\ln 2}{0,01} = 100 \ln 2 \approx 69,3$ min.
    \item $A_0 e^{-0,01 t} < 0,1\, A_0 \iff e^{-0,01 t} < 0,1 \iff -0,01 t < \ln 0,1 = -\ln 10 \iff t > \dfrac{\ln 10}{0,01} = 100 \ln 10 \approx 230,3$ min, soit environ 3 h 50 min.
\end{enumerate}
\end{corrige}

\coursec{Primitives des fonctions homographiques et usuelles}

Après avoir maîtrisé la résolution d'équations et d'inéquations faisant intervenir le logarithme népérien et l'exponentielle, nous abordons maintenant une étape fondamentale : la recherche de \cle{primitives}. Cette notion est la clé pour calculer des aires sous des courbes, modéliser des croissances ou des décroissances en technique (charge et décharge d'un condensateur, refroidissement d'une pièce, amortissement), et préparer l'étude complète des fonctions $x \mapsto \ln(ax+b)$ et $x \mapsto e^{ax+b}$ de la section suivante.

\begin{definition}[Primitive]
Soit $f$ une fonction définie sur un intervalle $I$. On dit qu'une fonction $F$ est une \cle{primitive} de $f$ sur $I$ si $F$ est dérivable sur $I$ et si, pour tout $x$ de $I$ :
\[
F'(x) = f(x).
\]
Si $F$ est une primitive de $f$ sur $I$, alors toutes les primitives de $f$ sur $I$ sont les fonctions $x \mapsto F(x) + C$, où $C$ est une constante réelle quelconque.
\end{definition}

\courssub{Rappel : tableau des primitives des fonctions de référence}

Avant d'attaquer les fonctions homographiques et composées, il est indispensable de connaître par cœur les primitives des fonctions usuelles sur leurs intervalles de définition. Ce formulaire est votre « boîte à outils » pour toute l'année. Chaque ligne s'obtient en lisant le tableau des dérivées « à l'envers ».

\begin{popfigure}
\centering
\begin{tabularx}{\linewidth}{|l|X|l|}\hline
\rowcolor{popblueL}
\textbf{Fonction $f(x)$} & \textbf{Une primitive $F(x)$} & \textbf{Intervalle $I$} \\ \hline
$k$ (constante) & $kx + C$ & $\mathbb{R}$ \\ \hline
$x^n$ ($n \in \mathbb{N}$) & $\dfrac{x^{n+1}}{n+1} + C$ & $\mathbb{R}$ \\ \hline
$\dfrac{1}{x^2}$ & $-\dfrac{1}{x} + C$ & $]-\infty,0[$ ou $]0,+\infty[$ \\ \hline
$\dfrac{1}{\sqrt{x}}$ & $2\sqrt{x} + C$ & $]0, +\infty[$ \\ \hline
$\dfrac{1}{x}$ & $\ln x + C$ & $]0, +\infty[$ \\ \hline
$\dfrac{1}{x}$ & $\ln(-x) + C$ & $]-\infty, 0[$ \\ \hline
$e^x$ & $e^x + C$ & $\mathbb{R}$ \\ \hline
$\cos x$ & $\sin x + C$ & $\mathbb{R}$ \\ \hline
$\sin x$ & $-\cos x + C$ & $\mathbb{R}$ \\ \hline
\end{tabularx}
\legende{Tableau des primitives usuelles au programme de Terminale technique.}
\end{popfigure}

\begin{attention}[Constante d'intégration]
Toute primitive s'entend \emph{à une constante près}. On note systématiquement $+ C$ (avec $C \in \mathbb{R}$) à la fin de chaque résultat. Oublier $C$ est une erreur de rédaction sanctionnée à l'examen. De même, une primitive se détermine toujours \textbf{sur un intervalle} : précisez-le.
\end{attention}

\courssub{Primitive de la forme $u'/u$ : le logarithme népérien}

C'est la primitive la plus utilisée en Terminale technique. Elle permet de primitiver toute fraction dont le numérateur est la dérivée du dénominateur (ou un multiple constant de celle-ci).

\begin{propriete}[Primitive de $u'/u$]
Soit $u$ une fonction dérivable et \textbf{strictement positive} sur un intervalle $I$.
Une primitive de $x \mapsto \dfrac{u'(x)}{u(x)}$ sur $I$ est :
\[
F(x) = \ln\big(u(x)\big) + C.
\]
\end{propriete}

\begin{experience}[Démonstration guidée (exigible)]
On pose $F(x) = \ln\big(u(x)\big)$ avec $u(x) > 0$ sur $I$. Dérivons $F$ avec la formule de dérivation des fonctions composées : $(\ln \circ u)' = \dfrac{u'}{u}$.
\begin{enumerate}
\item Rappel : si $v(x) = \ln x$, alors $v'(x) = \dfrac{1}{x}$.
\item Dérivée d'une composée : $\big(v \circ u\big)'(x) = u'(x) \times v'\big(u(x)\big)$.
\item Donc $F'(x) = u'(x) \times \dfrac{1}{u(x)} = \dfrac{u'(x)}{u(x)} = f(x)$.
\end{enumerate}
Conclusion : $F$ est bien une primitive de $f$ sur $I$. Cette démonstration est exigible au Probatoire, souvent sous la forme « Vérifier que $F$ est une primitive de $f$ ».
\end{experience}

\begin{methode}[Primitiver une fraction de type $u'/u$]
\begin{enumerate}
\item Identifier le dénominateur $u(x)$.
\item Calculer sa dérivée $u'(x)$.
\item Vérifier que le numérateur est égal à $u'(x)$ (ou à un multiple constant $k \cdot u'(x)$).
\item Vérifier le signe de $u$ sur l'intervalle considéré.
\item Conclure : une primitive est $\ln\big(u(x)\big) + C$ (ou $k\ln\big(u(x)\big) + C$).
\end{enumerate}
\end{methode}

\begin{exoresolu}[Application directe]
Énoncé : Déterminer une primitive de $f(x) = \dfrac{2x+3}{x^2+3x+5}$ sur $\mathbb{R}$.
\tcblower
Solution détaillée :
\begin{enumerate}
\item On pose $u(x) = x^2+3x+5$. Alors $u'(x) = 2x+3$.
\item Le numérateur est exactement $u'(x)$ : on est dans le cas $\dfrac{u'}{u}$.
\item Signe de $u$ : le discriminant de $x^2+3x+5$ est $\Delta = 9 - 20 = -11 < 0$. Le trinôme est donc du signe de $a=1>0$ : $u(x) > 0$ pour tout réel $x$. La fonction $\ln(u)$ est bien définie sur $\mathbb{R}$.
\item Conclusion : une primitive de $f$ sur $\mathbb{R}$ est
\[
F(x) = \ln(x^2+3x+5) + C.
\]
\end{enumerate}
\end{exoresolu}

\begin{attention}[Piège : numérateur non proportionnel à $u'$]
Si le numérateur n'est \textbf{pas} un multiple constant de la dérivée du dénominateur, on \textbf{ne peut pas} appliquer la formule $\ln(u)$. Par exemple, $\dfrac{x+1}{x^2+4}$ n'est pas de la forme $k\,\dfrac{u'}{u}$ (car $u' = 2x$ et $x+1$ n'est pas proportionnel à $2x$). Il faut alors décomposer la fraction en somme de deux termes.
\end{attention}

\courssub{Primitives des fonctions homographiques : écriture sous la forme $\alpha + \dfrac{\beta}{x-x_0}$}

Une fonction \cle{homographique} s'écrit $f(x) = \dfrac{ax+b}{cx+d}$ avec $c \neq 0$ et $ad-bc \neq 0$. Pour la primitiver, on la réécrit d'abord comme somme d'une constante et d'une fraction de type $\dfrac{1}{x-x_0}$.

\begin{propriete}[Forme réduite d'une homographique]
Toute fonction homographique $f(x) = \dfrac{ax+b}{cx+d}$ ($c \neq 0$, $ad-bc \neq 0$) peut s'écrire de manière unique sous la forme :
\[
f(x) = \alpha + \frac{\beta}{x - x_0}
\quad \text{avec} \quad
\alpha = \frac{a}{c}, \quad x_0 = -\frac{d}{c}, \quad \beta = \frac{bc - ad}{c^2}.
\]
$x_0$ est la valeur interdite de $f$ (asymptote verticale) et $\beta \neq 0$ car $ad-bc \neq 0$.
\end{propriete}

\begin{experience}[Démonstration guidée sur un exemple]
Montrons que $\dfrac{3x+1}{x-2} = 3 + \dfrac{7}{x-2}$ (pour $x \neq 2$).
\begin{enumerate}
\item Astuce : on écrit le numérateur en faisant apparaître le dénominateur : $3x+1 = 3(x-2) + 6 + 1 = 3(x-2) + 7$.
\item On sépare en deux fractions : $\dfrac{3x+1}{x-2} = \dfrac{3(x-2)}{x-2} + \dfrac{7}{x-2} = 3 + \dfrac{7}{x-2}$.
\item Vérification en réduisant au même dénominateur : $3 + \dfrac{7}{x-2} = \dfrac{3(x-2)+7}{x-2} = \dfrac{3x+1}{x-2}$. L'égalité est confirmée.
\end{enumerate}
Cette technique du « faire apparaître le dénominateur au numérateur » est la méthode attendue au Probatoire.
\end{experience}

\begin{methode}[Primitiver une fonction homographique]
\begin{enumerate}
\item Écrire $f(x) = \dfrac{ax+b}{cx+d}$ et repérer la valeur interdite $x_0 = -\dfrac{d}{c}$.
\item Décomposer : écrire $ax+b = \dfrac{a}{c}(cx+d) + \left(b - \dfrac{ad}{c}\right)$, puis séparer en deux fractions pour obtenir $f(x) = \alpha + \dfrac{\beta}{x-x_0}$.
\item Vérifier la décomposition en réduisant au même dénominateur.
\item Primitiver terme à terme : une primitive de $\alpha$ est $\alpha x$ ; une primitive de $\dfrac{\beta}{x-x_0}$ est $\beta\ln|x-x_0|$.
\item Conclure sur chaque intervalle de définition : $F(x) = \alpha x + \beta \ln|x-x_0| + C$.
\end{enumerate}
\end{methode}

\begin{exoresolu}[Exemple chiffré complet]
Énoncé : Déterminer une primitive de $f(x) = \dfrac{3x+1}{x-2}$ sur $]2, +\infty[$.
\tcblower
Solution détaillée :
\begin{enumerate}
\item Valeur interdite : $x_0 = 2$.
\item Décomposition (voir activité ci-dessus) : $f(x) = 3 + \dfrac{7}{x-2}$.
\item Sur $]2, +\infty[$, on a $x - 2 > 0$, donc une primitive de $\dfrac{7}{x-2}$ est $7\ln(x-2)$.
\item Conclusion : $F(x) = 3x + 7\ln(x-2) + C$.
\end{enumerate}
Vérification par dérivation : $F'(x) = 3 + \dfrac{7}{x-2} = \dfrac{3(x-2)+7}{x-2} = \dfrac{3x+1}{x-2} = f(x)$. Le résultat est validé.
\end{exoresolu}

\begin{exercice}[Entraînement]
Exercice 1 : Déterminer une primitive de chacune des fonctions suivantes sur l'intervalle indiqué :
\begin{enumerate}
\item $f(x) = \dfrac{2x-5}{x+3}$ sur $]-3, +\infty[$ ;
\item $g(x) = \dfrac{5x+1}{2x-4}$ sur $]2, +\infty[$ ;
\item $h(x) = \dfrac{x^2+1}{x}$ sur $]0, +\infty[$ (indication : diviser terme à terme).
\end{enumerate}
\end{exercice}

\begin{corrige}[Exercice 1]
\begin{enumerate}
\item $2x-5 = 2(x+3) - 11$, donc $f(x) = 2 - \dfrac{11}{x+3}$. Sur $]-3,+\infty[$, $x+3>0$ : $F(x) = 2x - 11\ln(x+3) + C$.
\item $5x+1 = \dfrac{5}{2}(2x-4) + 10 + 1 = \dfrac{5}{2}(2x-4) + 11$, donc $g(x) = \dfrac{5}{2} + \dfrac{11}{2x-4} = \dfrac{5}{2} + \dfrac{11}{2}\times\dfrac{1}{x-2}$. Sur $]2,+\infty[$ : $G(x) = \dfrac{5}{2}x + \dfrac{11}{2}\ln(x-2) + C$.
\item $h(x) = \dfrac{x^2+1}{x} = x + \dfrac{1}{x}$. Sur $]0,+\infty[$ : $H(x) = \dfrac{x^2}{2} + \ln x + C$.
\end{enumerate}
\end{corrige}

\courssub{Primitive de $\dfrac{1}{ax+b}$ et primitives de l'exponentielle}

La décomposition précédente repose sur un résultat fondamental : la primitive de l'inverse d'une fonction affine.

\begin{propriete}[Primitive de $\frac{1}{ax+b}$]
Pour $a \neq 0$, une primitive de $x \mapsto \dfrac{1}{ax+b}$ sur tout intervalle où $ax+b$ garde un signe constant est :
\[
F(x) = \frac{1}{a} \ln|ax+b| + C.
\]
\end{propriete}

\begin{experience}[Démonstration guidée]
On pose $u(x) = ax+b$, alors $u'(x) = a$.
\begin{enumerate}
\item On fait apparaître $u'$ au numérateur : $\dfrac{1}{ax+b} = \dfrac{1}{a} \times \dfrac{a}{ax+b} = \dfrac{1}{a}\times\dfrac{u'(x)}{u(x)}$.
\item On applique la formule de la primitive de $\dfrac{u'}{u}$ : une primitive de $\dfrac{u'}{u}$ est $\ln|u|$.
\item Conclusion : une primitive de $\dfrac{1}{ax+b}$ est $\dfrac{1}{a}\ln|ax+b| + C$.
\end{enumerate}
Vérification : $\left(\dfrac{1}{a}\ln|ax+b|\right)' = \dfrac{1}{a} \times \dfrac{a}{ax+b} = \dfrac{1}{ax+b}$.
\end{experience}

\begin{attention}[Erreur classique : oublier le coefficient $\frac{1}{a}$]
C'est l'erreur la plus fréquente au Probatoire.
\begin{itemize}
\item \textbf{Faux :} une primitive de $\dfrac{1}{2x+1}$ est $\ln|2x+1|$.
\item \textbf{Vrai :} une primitive de $\dfrac{1}{2x+1}$ est $\dfrac{1}{2}\ln|2x+1| + C$.
\end{itemize}
Règle mnémotechnique : « on divise par la dérivée de l'intérieur ». En cas de doute, dérivez votre résultat : vous devez retomber sur la fonction de départ.
\end{attention}

Pour l'exponentielle, la situation est symétrique : comme la dérivée de $e^{u}$ est $u'e^{u}$, la primitive de $u'e^{u}$ est $e^{u}$.

\begin{propriete}[Primitives liées à l'exponentielle]
\begin{itemize}
\item Si $u$ est dérivable sur $I$, une primitive de $u'(x)\,e^{u(x)}$ est $e^{u(x)} + C$.
\item En particulier, pour $a \neq 0$, une primitive de $e^{ax+b}$ est $\dfrac{1}{a}\,e^{ax+b} + C$.
\end{itemize}
\end{propriete}

\begin{attention}[Même piège : le coefficient $\frac{1}{a}$]
\begin{itemize}
\item \textbf{Faux :} une primitive de $e^{3x-2}$ est $e^{3x-2}$.
\item \textbf{Vrai :} une primitive de $e^{3x-2}$ est $\dfrac{1}{3}e^{3x-2} + C$.
\end{itemize}
Vérification : $\left(\dfrac{1}{3}e^{3x-2}\right)' = \dfrac{1}{3} \times 3 \times e^{3x-2} = e^{3x-2}$.
\end{attention}

\begin{exoresolu}[Mélange homographique et exponentielle]
Énoncé : Déterminer une primitive de $f(x) = \dfrac{4}{2x-1} + 5e^{3x+2}$ sur $\left]\dfrac{1}{2}, +\infty\right[$.
\tcblower
Solution détaillée :
\begin{itemize}
\item Premier terme : $\dfrac{4}{2x-1} = 4 \times \dfrac{1}{2x-1}$. Une primitive est $4 \times \dfrac{1}{2}\ln(2x-1) = 2\ln(2x-1)$ (on a $2x-1>0$ sur l'intervalle, donc pas de valeur absolue).
\item Second terme : une primitive de $5e^{3x+2}$ est $5 \times \dfrac{1}{3}e^{3x+2} = \dfrac{5}{3}e^{3x+2}$.
\end{itemize}
Conclusion : $F(x) = 2\ln(2x-1) + \dfrac{5}{3}e^{3x+2} + C$.
\end{exoresolu}

\courssub{Technique : faire apparaître la dérivée du dénominateur}

C'est la compétence transversale majeure de cette leçon : transformer une expression pour la ramener à une forme connue ($u'/u$ ou $u'e^u$).

\begin{methode}[Stratégie « faire apparaître $u'$ »]
\begin{enumerate}
\item Identifier une fonction $u(x)$ « candidate » (souvent le dénominateur, ou l'exposant de l'exponentielle).
\item Calculer sa dérivée $u'(x)$.
\item Comparer le numérateur (ou le facteur devant $e^u$) avec $u'(x)$ : s'il ne diffère que par une constante, multiplier et diviser par cette constante.
\item Appliquer la formule correspondante.
\end{enumerate}
\end{methode}

\begin{exoresolu}[Exemple type « faire apparaître $u'$ »]
Énoncé : Déterminer une primitive de $f(x) = \dfrac{x}{x^2+4}$ sur $\mathbb{R}$.
\tcblower
Solution détaillée :
\begin{enumerate}
\item On pose $u(x) = x^2+4$, alors $u'(x) = 2x$.
\item Le numérateur est $x$ : il manque le facteur $2$ pour obtenir $u'(x)$.
\item On écrit : $f(x) = \dfrac{1}{2} \times \dfrac{2x}{x^2+4} = \dfrac{1}{2}\times\dfrac{u'(x)}{u(x)}$.
\item Comme $x^2+4 > 0$ pour tout réel $x$, une primitive est :
\[
F(x) = \frac{1}{2}\ln(x^2+4) + C.
\]
\end{enumerate}
\end{exoresolu}

\begin{exoresolu}[Cas exponentiel]
Énoncé : Déterminer une primitive de $g(x) = x\,e^{x^2}$ sur $\mathbb{R}$.
\tcblower
Solution détaillée :
On pose $u(x) = x^2$, alors $u'(x) = 2x$.
On écrit : $g(x) = \dfrac{1}{2} \times 2x \times e^{x^2} = \dfrac{1}{2}\,u'(x)\,e^{u(x)}$.
Une primitive est donc : $G(x) = \dfrac{1}{2}e^{x^2} + C$.
Vérification : $G'(x) = \dfrac{1}{2}\times 2x \times e^{x^2} = x e^{x^2} = g(x)$.
\end{exoresolu}

\courssub{Vérification d'une primitive par dérivation}

Au Probatoire, on vous demandera souvent : « Vérifier que la fonction $F$ est une primitive de $f$ sur $I$ ». C'est une application directe de la définition : il suffit de dériver $F$.

\begin{methode}[Protocole de vérification]
\begin{enumerate}
\item Calculer $F'(x)$ en dérivant terme à terme (dérivées usuelles, produit, composée).
\item Simplifier l'expression obtenue.
\item Conclure par une phrase : « Pour tout $x$ de $I$, $F'(x) = f(x)$, donc $F$ est une primitive de $f$ sur $I$. »
\end{enumerate}
\end{methode}

\begin{exoresolu}[Vérification formelle]
Énoncé : Vérifier que $F(x) = x\ln x - x$ est une primitive de $f(x) = \ln x$ sur $]0, +\infty[$.
\tcblower
Solution détaillée :
$F$ est dérivable sur $]0,+\infty[$ comme somme et produit de fonctions dérivables. On dérive le produit $x\ln x$ avec la formule $(uv)' = u'v + uv'$ :
\begin{align*}
F'(x) &= 1 \times \ln x + x \times \frac{1}{x} - 1 \\
&= \ln x + 1 - 1 \\
&= \ln x = f(x).
\end{align*}
Donc $F$ est bien une primitive de $f$ sur $]0, +\infty[$. Ce résultat est à retenir : \textbf{une primitive de $\ln x$ est $x\ln x - x$}.
\end{exoresolu}

\courssub{Application au calcul d'aires}

Le lien entre primitives et aires est fondamental : \textbf{l'aire sous la courbe d'une fonction positive se calcule à l'aide d'une primitive évaluée entre deux bornes}.

\begin{savaistu}[Naissance du logarithme : l'aire sous l'hyperbole]
C'est en cherchant à calculer l'aire sous la courbe $y = \dfrac{1}{x}$ (l'hyperbole) que Grégoire de Saint-Vincent (1647) puis Nicolas Mercator ont mis en évidence les propriétés du logarithme népérien. L'aire du domaine délimité par l'axe des abscisses, la courbe $y=\dfrac{1}{x}$ et les droites $x=1$ et $x=t$ ($t>1$) vaut exactement $\ln t$. C'est la définition géométrique historique du logarithme népérien.
\end{savaistu}

\begin{popfigure}
\centering
\begin{tikzpicture}
\begin{axis}[
    axis lines=middle,
    xlabel=$x$, ylabel=$y$,
    xmin=0, xmax=4.5, ymin=0, ymax=2.5,
    xtick={1,2,3,4}, ytick={1,2},
    width=\linewidth, height=7cm,
    samples=100
]
\addplot[popblue!20, draw=none, domain=1:3] {1/x} \closedcycle;
\addplot[popcurve, thick, domain=0.42:4.2] {1/x};
\draw[dashed, popmuted] (axis cs:3,0) -- (axis cs:3,0.333);
\node[popblue, right] at (axis cs:3.4,0.5) {$y=\dfrac{1}{x}$};
\node[below] at (axis cs:3,0) {$t$};
\node[popdark] at (axis cs:1.8,0.45) {Aire $= \ln t$};
\end{axis}
\end{tikzpicture}
\legende{Aire sous l'hyperbole $y=\dfrac{1}{x}$ entre $1$ et $t$ : $\mathcal{A} = \ln t$ unités d'aire.}
\end{popfigure}

\begin{propriete}[Aire et primitive (admise)]
Si $f$ est continue et positive sur $[a,b]$ et si $F$ est une primitive de $f$ sur $[a,b]$, alors l'aire $\mathcal{A}$ (en unités d'aire) du domaine délimité par la courbe $y=f(x)$, l'axe des abscisses et les droites $x=a$ et $x=b$ est :
\[
\mathcal{A} = F(b) - F(a).
\]
\end{propriete}

\begin{exoresolu}[Calcul d'aire concret]
Énoncé : Calculer l'aire (en unités d'aire) du domaine délimité par la courbe $y = \dfrac{2}{x+1}$, l'axe des abscisses et les droites $x=0$ et $x=2$.
\tcblower
Solution détaillée :
\begin{enumerate}
\item Sur $[0,2]$, $x+1 > 0$ donc $f(x) = \dfrac{2}{x+1} > 0$ : la formule s'applique.
\item Une primitive : $F(x) = 2\ln(x+1)$ (car une primitive de $\dfrac{1}{x+1}$ est $\ln(x+1)$).
\item Calcul : $\mathcal{A} = F(2) - F(0) = 2\ln 3 - 2\ln 1 = 2\ln 3 - 0 = 2\ln 3$.
\end{enumerate}
L'aire vaut $2\ln 3 \approx 2,20$ unités d'aire.
\end{exoresolu}

\begin{exercice}[Synthèse : primitives et aire]
Exercice 2 :
\begin{enumerate}
\item Montrer que, pour tout $x > 1$, $f(x) = \dfrac{3x-2}{x-1} = 3 + \dfrac{1}{x-1}$.
\item En déduire une primitive $F$ de $f$ sur $]1, +\infty[$.
\item Calculer l'aire du domaine délimité par la courbe $y=f(x)$, l'axe des abscisses et les droites $x=2$ et $x=4$.
\end{enumerate}
\end{exercice}

\begin{corrige}[Exercice 2]
\begin{enumerate}
\item On réduit au même dénominateur : $3 + \dfrac{1}{x-1} = \dfrac{3(x-1)+1}{x-1} = \dfrac{3x-2}{x-1} = f(x)$.
\item Sur $]1,+\infty[$, $x-1>0$ : $F(x) = 3x + \ln(x-1) + C$.
\item Sur $[2,4]$, $f(x) > 0$. $\mathcal{A} = F(4) - F(2) = \big(12 + \ln 3\big) - \big(6 + \ln 1\big) = 6 + \ln 3$ unités d'aire (soit environ $7,10$).
\end{enumerate}
\end{corrige}

\begin{aretenir}[Résumé des formules clés pour le Probatoire]
\begin{itemize}
\item Primitive de $\dfrac{u'}{u}$ : $\ln|u| + C$ (avec $u$ de signe constant sur $I$).
\item Primitive de $\dfrac{1}{ax+b}$ : $\dfrac{1}{a}\ln|ax+b| + C$ \quad ($a \neq 0$).
\item Homographique : si $f(x) = \alpha + \dfrac{\beta}{x-x_0}$, alors $F(x) = \alpha x + \beta\ln|x-x_0| + C$.
\item Primitive de $u'e^{u}$ : $e^{u} + C$.
\item Primitive de $e^{ax+b}$ : $\dfrac{1}{a}e^{ax+b} + C$ \quad ($a \neq 0$).
\item Primitive de $\ln x$ sur $]0,+\infty[$ : $x\ln x - x + C$.
\item \textbf{Réflexe de sécurité :} en cas de doute sur un coefficient ou un signe, dérivez votre primitive : vous devez retrouver la fonction de départ.
\end{itemize}
\end{aretenir}

Cette maîtrise des primitives nous donne maintenant tous les outils pour étudier complètement les fonctions $x \mapsto \ln(ax+b)$ et $x \mapsto e^{ax+b}$ : ensemble de définition, dérivée, variations, limites, tableau de variations et tracé de courbes. C'est l'objet de la prochaine section.

\coursec{Étude des fonctions x ↦ ln(ax+b) et x ↦ exp(ax+b)}

Après avoir maîtrisé le calcul des primitives des fonctions homographiques (section précédente), nous abordons maintenant l'étude complète de deux familles de fonctions composées incontournables en technique : les fonctions logarithmes affines \(x \mapsto \ln(ax+b)\) et les fonctions exponentielles affines \(x \mapsto e^{ax+b}\). Ces fonctions modélisent des phénomènes réels variés : décharge d'un condensateur, croissance bactérienne, décroissance radioactive, loi d'atténuation du son, etc. La démarche d'étude sera systématique : domaine de définition, limites aux bornes, dérivée et sens de variation, tableau de variations, courbe représentative et position par rapport aux axes.

\courssub{Fonction logarithme affine : x ↦ ln(ax+b)}

\courssub{Ensemble de définition et domaine}

La fonction logarithme népérien \(\ln u\) n'est définie que pour \(u > 0\). Pour \(f(x) = \ln(ax+b)\), il faut donc imposer la condition sur l'expression affine à l'intérieur :

\begin{definition}[Domaine de définition de x ↦ ln(ax+b)]
Soient \(a, b \in \mathbb{R}\) avec \(a \neq 0\). La fonction \(f : x \mapsto \ln(ax+b)\) est définie sur l'ensemble :
\[
D_f = \{ x \in \mathbb{R} \mid ax + b > 0 \}
\]
\begin{itemize}
    \item Si \(a > 0\) : \(ax+b > 0 \iff x > -\frac{b}{a}\). Donc \(D_f = \left] -\frac{b}{a} ; +\infty \right[\).
    \item Si \(a < 0\) : \(ax+b > 0 \iff x < -\frac{b}{a}\). Donc \(D_f = \left] -\infty ; -\frac{b}{a} \right[\).
\end{itemize}
La droite d'équation \(x = -\frac{b}{a}\) est une \textbf{asymptote verticale} (la courbe s'en approche indéfiniment sans la toucher).
\end{definition}

\begin{attention}[Piège classique : l'inégalité]
Quand on résout \(ax+b > 0\), le sens de l'inégalité change si on divise par \(a < 0\).
\textbf{Exemple :} \(f(x) = \ln(2-3x)\). \(2-3x > 0 \iff -3x > -2 \iff x < \frac{2}{3}\).
Domaine : \(\left] -\infty ; \frac{2}{3} \right[\). Asymptote verticale : \(x = \frac{2}{3}\).
\end{attention}

\courssub{Dérivée et sens de variation}

On utilise la règle de dérivation des fonctions composées : \(( \ln(u) )' = \frac{u'}{u}\).

\begin{propriete}[Dérivée de ln(ax+b)]
Soit \(f(x) = \ln(ax+b)\) sur son domaine \(D_f\).
\[
f'(x) = \frac{a}{ax+b}
\]
\emph{Preuve :} \(u(x) = ax+b \Rightarrow u'(x) = a\). Donc \(f'(x) = \frac{u'(x)}{u(x)} = \frac{a}{ax+b}\).
\end{propriete}

\begin{aretenir}[Sens de variation selon le signe de a]
Sur \(D_f\), le dénominateur \(ax+b\) est \textbf{strictement positif} (par définition du domaine).
Le signe de \(f'(x)\) est donc \textbf{exclusivement le signe de \(a\)} :
\begin{itemize}
    \item Si \(a > 0\) : \(f'(x) > 0\) sur \(D_f\) \(\Rightarrow\) \(f\) est \textbf{strictement croissante}.
    \item Si \(a < 0\) : \(f'(x) < 0\) sur \(D_f\) \(\Rightarrow\) \(f\) est \textbf{strictement décroissante}.
\end{itemize}
\end{aretenir}

\courssub{Limites aux bornes et asymptotes}

Rappel : \(\lim_{u \to 0^+} \ln u = -\infty\) et \(\lim_{u \to +\infty} \ln u = +\infty\).

\begin{itemize}
    \item \textbf{Cas \(a > 0\)} (\(D_f = \left] -\frac{b}{a} ; +\infty \right[\)) :
    \begin{align*}
        \lim_{x \to \left(-\frac{b}{a}\right)^+} (ax+b) &= 0^+ \Rightarrow \lim_{x \to \left(-\frac{b}{a}\right)^+} f(x) = -\infty \quad \text{(asymptote verticale } x = -\tfrac{b}{a}\text{)} \\
        \lim_{x \to +\infty} (ax+b) &= +\infty \Rightarrow \lim_{x \to +\infty} f(x) = +\infty
    \end{align*}
    \item \textbf{Cas \(a < 0\)} (\(D_f = \left] -\infty ; -\frac{b}{a} \right[\)) :
    \begin{align*}
        \lim_{x \to -\infty} (ax+b) &= +\infty \Rightarrow \lim_{x \to -\infty} f(x) = +\infty \\
        \lim_{x \to \left(-\frac{b}{a}\right)^-} (ax+b) &= 0^+ \Rightarrow \lim_{x \to \left(-\frac{b}{a}\right)^-} f(x) = -\infty \quad \text{(asymptote verticale } x = -\tfrac{b}{a}\text{)}
    \end{align*}
\end{itemize}

\courssub{Intersections avec les axes et points remarquables}

\begin{itemize}
    \item \textbf{Intersection avec l'axe des abscisses (zéro) :} \(f(x) = 0 \iff \ln(ax+b) = 0 \iff ax+b = 1 \iff x = \frac{1-b}{a}\).
    Ce zéro appartient à \(D_f\) car \(1 > 0\).
    \item \textbf{Intersection avec l'axe des ordonnées :} Si \(0 \in D_f\) (c'est-à-dire \(b > 0\)), alors \(f(0) = \ln b\).
    \item \textbf{Point à abscisse \(x_0 = \frac{1-b}{a}\) (le zéro) :} La tangente a pour pente \(f'(x_0) = \frac{a}{1} = a\). Équation : \(y = a\left(x - \frac{1-b}{a}\right)\).
\end{itemize}

\begin{exemplebox}[Étude complète de $f(x) = ln(3x - 6)$]
Soit \(f(x) = \ln(3x-6)\).
\begin{enumerate}
    \item \textbf{Domaine :} \(3x-6 > 0 \iff x > 2\). \(D_f = ]2; +\infty[\). Asymptote verticale : \(x=2\).
    \item \textbf{Dérivée :} \(f'(x) = \frac{3}{3x-6} = \frac{1}{x-2}\). \(a=3>0 \Rightarrow f\) strictement croissante sur \(]2;+\infty[\).
    \item \textbf{Limites :} \(\lim_{x\to 2^+} f(x) = -\infty\) ; \(\lim_{x\to +\infty} f(x) = +\infty\).
    \item \textbf{Zéro :} \(3x-6=1 \iff x = \frac{7}{3} \approx \num{2,33}\). \(f(7/3)=0\).
    \item \textbf{Ordonnée à l'origine :} \(0 \notin D_f\) (pas d'intersection avec Oy).
    \item \textbf{Tableau de variations :}
    \begin{center}
    \begin{tabular}{|c|cccc|}\hline
    \(x\) & \(2\) & & \(+\infty\) & \\ \hline
    \(f'(x)\) & & \(+\) & & \\ \hline
    \(f(x)\) & \(-\infty\) & \(\nearrow\) & \(+\infty\) & \\ \hline
    \end{tabular}
    \end{center}
\end{enumerate}
\end{exemplebox}

\begin{popfigure}
\begin{tikzpicture}
\begin{axis}[
    width=\linewidth, height=7cm,
    axis lines=middle, xlabel=$x$, ylabel=$y$,
    xmin=1.5, xmax=5, ymin=-4, ymax=4,
    xtick={2,2.333,3,4}, xticklabels={$2$, $\frac{7}{3}$, $3$, $4$},
    ytick={-3,-2,-1,1,2,3},
    domain=2.01:5, samples=200,
    popcurve,
    clip=false
]
\addplot {ln(3*x-6)};
% Asymptote verticale
\draw[popmuted, dashed, thick] (axis cs:2,-4) -- (axis cs:2,4) node[above, pos=0.9, popmuted] {$x=2$};
% Point zéro
\addplot[only marks, mark=*, mark size=2, poporange] coordinates {(2.333,0)} node[below right, poporange] {$(\frac{7}{3},0)$};
% Tangente au zéro (pente = 3)
\addplot[popgreen, dashed, domain=1.8:2.8] {3*(x-2.333)};
\end{axis}
\end{tikzpicture}
\legende{Courbe de \(f(x) = \ln(3x-6)\) : asymptote verticale \(x=2\), zéro en \(\frac{7}{3}\), tangente au zéro (pente 3).}
\end{popfigure}

\courssub{Fonction exponentielle affine : $x \mapsto e^{ax+b}$}

\courssub{Ensemble de définition et dérivée}

L'exponentielle \(e^u\) est définie pour tout réel \(u\).

\begin{definition}[Domaine et dérivée de $x \mapsto e^{ax+b}$]
Soient \(a, b \in \mathbb{R}\), \(a \neq 0\). La fonction \(g : x \mapsto e^{ax+b}\) est définie sur \(\mathbb{R}\).
Sa dérivée est :
\[
g'(x) = a \cdot e^{ax+b}
\]
\emph{Preuve :} \(v(x) = ax+b \Rightarrow v'(x) = a\). \(g'(x) = v'(x) \cdot e^{v(x)} = a e^{ax+b}\).
\end{definition}

\begin{aretenir}[Sens de variation]
Puisque \(e^{ax+b} > 0\) pour tout \(x \in \mathbb{R}\), le signe de \(g'(x)\) est \textbf{exclusivement le signe de \(a\)} :
\begin{itemize}
    \item Si \(a > 0\) : \(g'(x) > 0\) sur \(\mathbb{R}\) \(\Rightarrow\) \(g\) \textbf{strictement croissante} sur \(\mathbb{R}\).
    \item Si \(a < 0\) : \(g'(x) < 0\) sur \(\mathbb{R}\) \(\Rightarrow\) \(g\) \textbf{strictement décroissante} sur \(\mathbb{R}\).
\end{itemize}
\end{aretenir}

\courssub{Limites aux bornes et asymptotes horizontales}

Rappel : \(\lim_{u \to -\infty} e^u = 0\) et \(\lim_{u \to +\infty} e^u = +\infty\).

\begin{itemize}
    \item \textbf{Cas \(a > 0\)} :
    \begin{align*}
        \lim_{x \to -\infty} (ax+b) &= -\infty \Rightarrow \lim_{x \to -\infty} g(x) = 0^+ \quad \text{(asymptote horizontale } y=0 \text{ à gauche)} \\
        \lim_{x \to +\infty} (ax+b) &= +\infty \Rightarrow \lim_{x \to +\infty} g(x) = +\infty
    \end{align*}
    \item \textbf{Cas \(a < 0\)} :
    \begin{align*}
        \lim_{x \to -\infty} (ax+b) &= +\infty \Rightarrow \lim_{x \to -\infty} g(x) = +\infty \\
        \lim_{x \to +\infty} (ax+b) &= -\infty \Rightarrow \lim_{x \to +\infty} g(x) = 0^+ \quad \text{(asymptote horizontale } y=0 \text{ à droite)}
    \end{align*}
\end{itemize}

\courssub{Intersections avec les axes et convexité}

\begin{itemize}
    \item \textbf{Axe des ordonnées :} \(g(0) = e^b\). Point \(A(0; e^b)\). Tangente en 0 : pente \(g'(0) = a e^b\). Équation : \(y = a e^b x + e^b\).
    \item \textbf{Axe des abscisses :} \(e^{ax+b} > 0\) pour tout \(x\). \textbf{Pas d'intersection} avec l'axe des abscisses (la courbe reste strictement au-dessus).
    \item \textbf{Convexité :} \(g''(x) = a^2 e^{ax+b} > 0\). La courbe est \textbf{convexe} sur \(\mathbb{R}\) (elle se situe au-dessus de ses tangentes).
\end{itemize}

\begin{exemplebox}[Étude complète de $g(x) = e^{2x+1}$]
Soit \(g(x) = e^{2x+1}\).
\begin{enumerate}
    \item \textbf{Domaine :} \(\mathbb{R}\).
    \item \textbf{Dérivée :} \(g'(x) = 2 e^{2x+1} > 0\). \(g\) strictement croissante sur \(\mathbb{R}\).
    \item \textbf{Limites :} \(\lim_{x\to -\infty} g(x) = 0\) (asymptote \(y=0\) à gauche) ; \(\lim_{x\to +\infty} g(x) = +\infty\).
    \item \textbf{Point remarquable :} \(g(0) = e^1 = e \approx \num{2,718}\). Tangente en 0 : \(y = 2e x + e\).
    \item \textbf{Tableau de variations :}
    \begin{center}
    \begin{tabular}{|c|cccc|}\hline
    \(x\) & \(-\infty\) & & \(+\infty\) & \\ \hline
    \(g'(x)\) & & \(+\) & & \\ \hline
    \(g(x)\) & \(0\) & \(\nearrow\) & \(+\infty\) & \\ \hline
    \end{tabular}
    \end{center}
\end{enumerate}
\end{exemplebox}

\begin{popfigure}
\begin{tikzpicture}
\begin{axis}[
    width=\linewidth, height=7cm,
    axis lines=middle, xlabel=$x$, ylabel=$y$,
    xmin=-3, xmax=2, ymin=-0.5, ymax=8,
    xtick={-2,-1,0,1}, ytick={0,2.718,4,6,8},
    yticklabels={$0$, $e$, $4$, $6$, $8$},
    domain=-3:1.5, samples=200,
    popcurve,
    clip=false
]
\addplot {exp(2*x+1)};
% Asymptote horizontale
\draw[popmuted, dashed, thick] (axis cs:-3,0) -- (axis cs:2,0) node[right, popmuted] {$y=0$};
% Point (0, e)
\addplot[only marks, mark=*, mark size=2, poporange] coordinates {(0,2.718)} node[above left, poporange] {$(0,e)$};
% Tangente en 0 : y = 2e x + e
\addplot[popgreen, dashed, domain=-1.5:0.5] {2*2.718*x + 2.718};
\end{axis}
\end{tikzpicture}
\legende{Courbe de \(g(x) = e^{2x+1}\) : asymptote horizontale \(y=0\) (à gauche), point \((0,e)\), tangente en 0.}
\end{popfigure}

\courssub{Tableaux de variations types (Synthèse visuelle)}

\begin{popfigure}
\begin{center}
\begin{tabularx}{\linewidth}{|X|X|}
\hline
\rowcolor{popblueL} \textbf{Cas \(a > 0\) : \(f(x) = \ln(ax+b)\)} & \textbf{Cas \(a < 0\) : \(f(x) = \ln(ax+b)\)} \\
\hline
\(D_f = \left] -\frac{b}{a} ; +\infty \right[\) & \(D_f = \left] -\infty ; -\frac{b}{a} \right[\) \\
\hline
\begin{tabular}{|c|cccc|}\hline
\(x\) & \(-\frac{b}{a}\) & & \(+\infty\) & \\ \hline
\(f'(x)\) & & \(+\) & & \\ \hline
\(f(x)\) & \(-\infty\) & \(\nearrow\) & \(+\infty\) & \\ \hline
\end{tabular}
&
\begin{tabular}{|c|cccc|}\hline
\(x\) & \(-\infty\) & & \(-\frac{b}{a}\) & \\ \hline
\(f'(x)\) & & \(-\) & & \\ \hline
\(f(x)\) & \(+\infty\) & \(\searrow\) & \(-\infty\) & \\ \hline
\end{tabular}
\\
\hline
\rowcolor{poporangeL} \textbf{Cas \(a > 0\) : \(g(x) = e^{ax+b}\)} & \textbf{Cas \(a < 0\) : \(g(x) = e^{ax+b}\)} \\
\hline
\(D_g = \mathbb{R}\) & \(D_g = \mathbb{R}\) \\
\hline
\begin{tabular}{|c|cccc|}\hline
\(x\) & \(-\infty\) & & \(+\infty\) & \\ \hline
\(g'(x)\) & & \(+\) & & \\ \hline
\(g(x)\) & \(0\) & \(\nearrow\) & \(+\infty\) & \\ \hline
\end{tabular}
&
\begin{tabular}{|c|cccc|}\hline
\(x\) & \(-\infty\) & & \(+\infty\) & \\ \hline
\(g'(x)\) & & \(-\) & & \\ \hline
\(g(x)\) & \(+\infty\) & \(\searrow\) & \(0\) & \\ \hline
\end{tabular}
\\
\hline
\end{tabularx}
\end{center}
\legende{Tableaux de variations types pour les deux familles. Noter l'inversion du sens de variation selon le signe de \(a\).}
\end{popfigure}

\courssub{Méthode : Démarche complète d'étude d'une fonction}

\begin{methode}[Plan type rédigé pour l'étude d'une fonction f (Probatoire / Bac Tech)]
Pour toute fonction \(f\) (ici \(x \mapsto \ln(ax+b)\) ou \(x \mapsto e^{ax+b}\)), rédigez l'étude dans cet ordre :
\begin{enumerate}
    \item \textbf{Domaine de définition \(D_f\)} : Résoudre l'inéquation (pour ln) ou écrire \(\mathbb{R}\) (pour exp). Préciser l'asymptote verticale si ln.
    \item \textbf{Parité / Périodicité (si pertinent)} : Généralement ni pair, ni impair, ni périodique pour ces familles.
    \item \textbf{Limites aux bornes de \(D_f\)} : Calculer \(\lim_{x \to \text{bornes}} f(x)\). Déduire les asymptotes (verticales pour ln, horizontales pour exp).
    \item \textbf{Dérivée \(f'(x)\) et signe} : Calculer \(f'(x)\) par dérivation composée. Étudier le signe sur \(D_f\) (signe de \(a\) pour nos familles). En déduire les variations.
    \item \textbf{Tableau de variations} : Construire le tableau avec les valeurs limites et les flèches (\(\nearrow\) ou \(\searrow\)).
    \item \textbf{Points remarquables et courbe} : Zéro(s) (résoudre \(f(x)=0\)), intersection avec Oy (\(f(0)\) si \(0 \in D_f\)), tangente(s) si demandée. Tracer la courbe dans un repère orthogonal.
\end{enumerate}
\textbf{Conseil :} Justifiez chaque étape par une phrase ("Car \(a>0\), \(f'(x)>0\) donc \(f\) croissante").
\end{methode}

\courssub{Exemple d'application technique : Décharge d'un condensateur}

\begin{savaistu}[Applications réelles : RC, Radioactivité, Épidémies]
\begin{itemize}
    \item \textbf{Électricité (circuit RC) :} La tension aux bornes d'un condensateur qui se décharge suit \(u(t) = U_0 e^{-t/\tau}\) (\(\tau = RC\) constante de temps). C'est une exponentielle affine \(a = -1/\tau < 0\) : décroissante, asymptote \(y=0\).
    \item \textbf{Physique nucléaire :} La décroissance radioactive \(N(t) = N_0 e^{-\lambda t}\) (loi de Rutherford-Soddy). Demi-vie \(T_{1/2} = \ln 2 / \lambda\).
    \item \textbf{Épidémiologie / Biologie :} Croissance exponentielle initiale d'une épidémie \(I(t) = I_0 e^{kt}\) (\(k>0\)) ou croissance logistique (saturation).
    \item \textbf{Acoustique :} Niveau sonore en décibels \(L = 20 \ln(P/P_0)\) : fonction logarithme affine.
\end{itemize}
\end{savaistu}

\begin{exoresolu}[Décharge d'un condensateur]
Un condensateur de capacité \(C = 100~\mu\text{F}\) se décharge dans une résistance \(R = 10~\text{k}\Omega\). La tension \(u(t)\) (en volts) à l'instant \(t\) (en secondes) est modélisée par \(u(t) = 12 e^{-t}\) pour \(t \ge 0\).
\begin{enumerate}
    \item Étudier la fonction \(u\) sur \([0; +\infty[\) (domaine, variations, limites).
    \item Déterminer le temps nécessaire pour que la tension passe sous 1 V.
    \item Tracer la courbe et la tangente à l'origine.
\end{enumerate}
\tcblower
\textbf{Solution :}
\begin{enumerate}
    \item \textbf{Domaine :} \([0; +\infty[\) (temps positif). \textbf{Dérivée :} \(u'(t) = -12 e^{-t} < 0\). \(u\) est \textbf{strictement décroissante}. \textbf{Limites :} \(u(0)=12\) ; \(\lim_{t\to +\infty} u(t) = 0\) (asymptote horizontale \(y=0\)).
    \item \(u(t) < 1 \iff 12 e^{-t} < 1 \iff e^{-t} < \frac{1}{12} \iff -t < \ln(1/12) = -\ln 12 \iff t > \ln 12\).
    \(\ln 12 \approx \num{2,485} \text{ s}\). Au bout d'environ \textbf{2,5 secondes}, la tension est inférieure à 1 V.
    \item \textbf{Tangente à l'origine (t=0) :} \(u(0)=12\), \(u'(0)=-12\). Équation : \(y = -12t + 12\).
\end{enumerate}
\end{exoresolu}

\begin{popfigure}
\begin{tikzpicture}
\begin{axis}[
    width=\linewidth, height=7cm,
    axis lines=middle, xlabel=$t\,(\text{s})$, ylabel=$u\,(\text{V})$,
    xmin=0, xmax=5, ymin=0, ymax=13,
    xtick={0,1,2,2.485,3,4,5}, xticklabels={$0$, $1$, $2$, $\ln 12$, $3$, $4$, $5$},
    ytick={0,1,12}, yticklabels={$0$, $1\,\text{V}$, $12\,\text{V}$},
    domain=0:5, samples=200,
    popcurve,
]
\addplot {12*exp(-x)};
% Seuil 1V
\draw[poporange, dashed] (axis cs:0,1) -- (axis cs:2.485,1) node[midway, above] {Seuil 1 V};
\draw[poporange, dashed] (axis cs:2.485,1) -- (axis cs:2.485,0) node[below] {$\ln 12$};
% Tangente à l'origine
\addplot[popgreen, dashed, domain=0:1] {-12*x + 12};
\end{axis}
\end{tikzpicture}
\legende{Décharge du condensateur : courbe \(u(t)=12e^{-t}\), tangente à l'origine, temps pour atteindre 1 V.}
\end{popfigure}

\courssub{Organigramme de la démarche d'étude}

\begin{popfigure}
\begin{tikzpicture}[
    node distance=1.2cm and 2cm,
    box/.style={draw=popblue, thick, rounded corners, fill=popblueL, text width=3.5cm, align=center, minimum height=1.2cm},
    diamond/.style={draw=poporange, thick, aspect=2, fill=poporangeL, text width=3cm, align=center, inner sep=2pt},
    arrow/.style={-Stealth, thick, popdark}
]
\node[box, align=center] (start){1. Déterminer le domaine $D_f$\\(Inéquation $ax+b>0$ ou $\mathbb{R}$)};
\node[diamond, below of=start] (type) {Type de fonction ?};
\node[box, below left=1.5cm and 1cm of type, align=center] (ln){2a. $\ln(ax+b)$ :\\Asymptote verticale $x=-b/a$\\Limites $\to \pm\infty$};
\node[box, below right=1.5cm and 1cm of type, align=center] (exp){2b. $e^{ax+b}$ :\\Asymptote horizontale $y=0$\\Limites $0$ et $+\infty$};
\node[box, below=1.5cm of ln, align=center] (der){3. Calculer $f'(x)$\\(Règle chaîne : $\frac{a}{ax+b}$ ou $a e^{ax+b}$)};
\node[diamond, below of=der] (sign) {Signe de $a$ ?};
\node[box, below left=1.2cm and 0.5cm of sign, align=center] (apos){$a>0$ : $f'>0$\\$\Rightarrow$ Croissante};
\node[box, below right=1.2cm and 0.5cm of sign, align=center] (aneg){$a<0$ : $f'<0$\\$\Rightarrow$ Décroissante};
\node[box, below=1.5cm of apos, align=center] (tab){4. Tableau de variations\\(Bornes, flèches $\nearrow/\searrow$)};
\node[box, below of=tab, align=center] (end){5. Points remarquables\\Zéros, $f(0)$, Tangentes\\$\to$ Tracé de la courbe};
\draw[arrow] (start) -- (type);
\draw[arrow] (type) -- node[left, pos=0.5, popmuted] {ln} (ln);
\draw[arrow] (type) -- node[right, pos=0.5, popmuted] {exp} (exp);
\draw[arrow] (ln) -- (der);
\draw[arrow] (exp) -- (der);
\draw[arrow] (der) -- (sign);
\draw[arrow] (sign) -- node[left, pos=0.5, popmuted] {$+$} (apos);
\draw[arrow] (sign) -- node[right, pos=0.5, popmuted] {$-$} (aneg);
\draw[arrow] (apos) -- (tab);
\draw[arrow] (aneg) -- (tab);
\draw[arrow] (tab) -- (end);
\end{tikzpicture}
\legende{Organigramme de la démarche complète d'étude d'une fonction $x \mapsto \ln(ax+b)$ ou $x \mapsto e^{ax+b}$.}
\end{popfigure}

\begin{exercice}[Entraînement : Étude complète]
Étudier complètement les fonctions suivantes (domaine, limites, dérivée, variations, tableau, courbe, points remarquables) :
\begin{enumerate}
    \item $f(x) = \ln(4 - 2x)$
    \item $g(x) = e^{-3x+2}$
    \item $h(x) = \ln(x^2 - 4)$ \textit{(Attention : ce n'est pas de la forme $ax+b$, mais on utilise la même méthode avec $u(x)=x^2-4$)}
\end{enumerate}
\end{exercice}

\begin{corrige}[Exercice n°1 : $f(x) = ln(4-2x)$]
\begin{enumerate}
    \item \textbf{Domaine :} $4-2x > 0 \iff -2x > -4 \iff x < 2$. $D_f = ]-\infty; 2[$. Asymptote verticale $x=2$.
    \item \textbf{Dérivée :} $f'(x) = \frac{-2}{4-2x} = \frac{-1}{2-x}$. $a=-2<0 \Rightarrow f'(x)<0$. $f$ \textbf{décroissante} sur $]-\infty; 2[$.
    \item \textbf{Limites :} $\lim_{x\to -\infty} f(x) = +\infty$ ; $\lim_{x\to 2^-} f(x) = -\infty$.
    \item \textbf{Zéro :} $4-2x=1 \iff x=1,5$. $f(1,5)=0$.
    \item \textbf{Ordonnée à l'origine :} $f(0) = \ln 4 \approx \num{1,386}$.
    \item \textbf{Tableau :}
    \begin{center}
    \begin{tabular}{|c|cccc|}\hline
    $x$ & $-\infty$ & & $2$ & \\ \hline
    $f'(x)$ & & $-$ & & \\ \hline
    $f(x)$ & $+\infty$ & $\searrow$ & $-\infty$ & \\ \hline
    \end{tabular}
    \end{center}
\end{enumerate}
\end{corrige}

\begin{corrige}[Exercice n°2 : $g(x) = e^{-3x+2}$]
\begin{enumerate}
    \item \textbf{Domaine :} $\mathbb{R}$.
    \item \textbf{Dérivée :} $g'(x) = -3 e^{-3x+2} < 0$. $g$ \textbf{décroissante} sur $\mathbb{R}$.
    \item \textbf{Limites :} $\lim_{x\to -\infty} g(x) = +\infty$ ; $\lim_{x\to +\infty} g(x) = 0$ (asymptote $y=0$ à droite).
    \item \textbf{Ordonnée à l'origine :} $g(0) = e^2 \approx \num{7,389}$. Tangente en 0 : $y = -3e^2 x + e^2$.
    \item \textbf{Tableau :}
    \begin{center}
    \begin{tabular}{|c|cccc|}\hline
    $x$ & $-\infty$ & & $+\infty$ & \\ \hline
    $g'(x)$ & & $-$ & & \\ \hline
    $g(x)$ & $+\infty$ & $\searrow$ & $0$ & \\ \hline
    \end{tabular}
    \end{center}
\end{enumerate}
\end{corrige}

\begin{corrige}[Exercice n°3 : $h(x) = ln(x^2-4)$]
\begin{enumerate}
    \item \textbf{Domaine :} $x^2-4 > 0 \iff (x-2)(x+2) > 0 \iff x < -2$ ou $x > 2$.
    $D_h = ]-\infty; -2[ \cup ]2; +\infty[$. Asymptotes verticales : $x=-2$ et $x=2$.
    \item \textbf{Dérivée :} $h'(x) = \frac{2x}{x^2-4}$.
    \item \textbf{Signe de la dérivée et variations :}
    \begin{itemize}
        \item Sur $]-\infty; -2[$ : $2x < 0$ et $x^2-4 > 0 \Rightarrow h'(x) < 0$ ($h$ décroissante).
        \item Sur $]2; +\infty[$ : $2x > 0$ et $x^2-4 > 0 \Rightarrow h'(x) > 0$ ($h$ croissante).
    \end{itemize}
    \item \textbf{Limites :}
    \begin{align*}
        \lim_{x\to -\infty} h(x) &= +\infty, \quad \lim_{x\to -2^-} h(x) = -\infty \\
        \lim_{x\to 2^+} h(x) &= -\infty, \quad \lim_{x\to +\infty} h(x) = +\infty
    \end{align*}
    \item \textbf{Zéros :} $x^2-4 = 1 \iff x^2 = 5 \iff x = -\sqrt{5}$ (dans $]-\infty;-2[$) ou $x = \sqrt{5}$ (dans $]2;+\infty[$).
    \item \textbf{Ordonnée à l'origine :} $0 \notin D_h$.
    \item \textbf{Tableaux de variations :}
    \begin{center}
    \begin{tabular}{|c|cccc|}\hline
    $x$ & $-\infty$ & & $-2$ & \\ \hline
    $h'(x)$ & & $-$ & & \\ \hline
    $h(x)$ & $+\infty$ & $\searrow$ & $-\infty$ & \\ \hline
    \end{tabular}
    \qquad
    \begin{tabular}{|c|cccc|}\hline
    $x$ & $2$ & & $+\infty$ & \\ \hline
    $h'(x)$ & & $+$ & & \\ \hline
    $h(x)$ & $-\infty$ & $\nearrow$ & $+\infty$ & \\ \hline
    \end{tabular}
    \end{center}
\end{enumerate}
\end{corrige}

\begin{attention}[Erreurs fatales à éviter au Probatoire]
\begin{itemize}
    \item \textbf{Oublier le domaine de définition} pour $\ln(ax+b)$ (note éliminatoire fréquente).
    \item \textbf{Ne pas justifier le signe de la dérivée} : écrire "$f'(x) = \frac{a}{ax+b} > 0$" sans dire "car $ax+b>0$ sur $D_f$ et $a>0$".
    \item \textbf{Confondre asymptote verticale et horizontale} : $\ln$ a une asymptote \textbf{verticale} ($x=cte$), $\exp$ a une asymptote \textbf{horizontale} ($y=0$).
    \item \textbf{Dire que $\exp$ coupe l'axe des abscisses} : $e^u > 0$ toujours, jamais de zéro.
    \item \textbf{Mauvaise lecture du tableau de variations} : les flèches ($\nearrow, \searrow$) vont dans la ligne de $f$, les signes ($+, -$) dans la ligne de $f'$.
\end{itemize}
\end{attention}

\coursec{Applications à des situations concrètes}

Après avoir maîtrisé l'étude complète des fonctions $x \mapsto \ln(ax+b)$ et $x \mapsto \exp(ax+b)$, nous disposons maintenant de tous les outils pour \textbf{modéliser des situations réelles} issues de la gestion, de l'économie, de la physique ou de la biologie. Le passage du modèle mathématique à la réalité technique exige une rigueur particulière : il ne suffit pas de résoudre une équation, il faut \emph{interpréter} le nombre trouvé, vérifier sa cohérence avec le contexte et connaître les limites du modèle utilisé.

\courssub{Modélisation de la croissance exponentielle}

De nombreux phénomènes techniques se caractérisent par le fait que la \textbf{vitesse d'évolution} d'une grandeur est \emph{proportionnelle} à la valeur de cette grandeur à l'instant considéré. C'est le cas de l'évolution d'une clientèle, de la croissance d'une population bactérienne, ou du capital placé à intérêts composés.

\begin{propriete}[Loi exponentielle : propriété fondamentale (admise)]
Soit une grandeur $y$ fonction du temps $t$. Si le taux de variation instantané de $y$ est proportionnel à $y$ lui-même, c'est-à-dire si $y'(t) = k \cdot y(t)$ avec $k \in \mathbb{R}^*$, alors $y$ suit une \textbf{loi exponentielle} :
\[
y(t) = y_0 \cdot e^{kt}
\]
où $y_0 = y(0)$ est la valeur initiale.
\begin{itemize}
    \item Si $k > 0$ : \textbf{croissance exponentielle} ($k$ est le taux de croissance instantané).
    \item Si $k < 0$ : \textbf{décroissance exponentielle} ($|k|$ est le taux de décroissance instantané).
\end{itemize}
Au Probatoire, cette propriété est \textbf{admise} (pas de démonstration exigible), mais il faut savoir l'\textbf{appliquer} et en \textbf{interpréter} les paramètres.
\end{propriete}

\begin{savaistu}[Les banques camerounaises et les intérêts composés]
Au Cameroun, comme partout dans le monde, les établissements de crédit (Afriland First Bank, UBA, BICEC, etc.) calculent les intérêts composés selon une formule de type exponentiel. Si vous placez un capital $C_0$ à un taux annuel $t$ (exprimé en décimal, ex : 5\,\% $\to$ 0,05), le capital au bout de $n$ années est $C_n = C_0 (1+t)^n$. En posant $k = \ln(1+t)$, on retrouve la forme continue $C(t) = C_0 e^{kt}$. C'est cette modélisation continue qui permet aux gestionnaires de comparer des placements ou d'estimer le coût réel d'un crédit pour l'achat de matériel d'atelier.
\end{savaistu}

\begin{exemplebox}[Croissance d'une clientèle]
Une entreprise de maintenance informatique à Douala démarre avec $500$ clients. On estime que sa clientèle suit un modèle de croissance exponentielle continue de coefficient $k = 0,3$ par mois. On modélise le nombre de clients $N(t)$ après $t$ mois par :
\[
N(t) = 500 \cdot e^{0,3t}
\]
\begin{itemize}
    \item $500$ : clientèle initiale (valeur en $t=0$).
    \item $0,3$ : coefficient de croissance instantané mensuel, \textbf{donné par l'énoncé} (il ne se confond pas avec le taux discret de $30\,\%$ par mois, qui correspondrait à $k = \ln(1,3) \approx 0,262$).
    \item Domaine de validité : $t \ge 0$ (le temps ne peut pas être négatif dans ce contexte).
\end{itemize}
\end{exemplebox}

\begin{popfigure}
\begin{tikzpicture}
\begin{axis}[
    width=\linewidth,
    height=8cm,
    axis lines=middle,
    xlabel={$t$ (mois)},
    ylabel={$N(t)$ (clients)},
    domain=0:10,
    samples=200,
    xmin=0, xmax=10.5,
    ymin=0, ymax=11000,
    xtick={0,2,4,6,8,10},
    ytick={0,2000,4000,6000,8000,10000},
    yticklabels={0,2000,4000,6000,8000,10\,000},
    grid=major,
    grid style={popline},
    clip=false,
    every axis x label/.style={at={(ticklabel* cs:1.02)},anchor=west},
    every axis y label/.style={at={(ticklabel* cs:1.02)},anchor=south},
]
% Courbe principale
\addplot[popcurve, thick, domain=0:10] {500*exp(0.3*x)};
% Seuil 10 000
\addplot[poporange, dashed, thick] coordinates {(0,10000) (10,10000)} node[pos=0, left, poporange, font=\small] {Seuil 10\,000};
% Intersection : 500*e^(0.3t)=10000 -> t=ln(20)/0.3 ~ 9.98
\draw[popgreen, dashed, thick] (9.98,0) -- (9.98,10000);
\node[popgreen, anchor=north west, font=\small] at (9.98,0) {$t \approx 10$ mois};
% Point initial
\addplot[only marks, mark=*, mark size=2pt, popblue] coordinates {(0,500)} node[above right, font=\small] {$(0\,; 500)$};
\end{axis}
\end{tikzpicture}
\legende{Croissance de la clientèle $N(t)=500 e^{0,3t}$. Lecture graphique du temps nécessaire pour atteindre 10\,000 clients ($t \approx 10$ mois).}
\end{popfigure}

\begin{exoresolu}[Temps de doublement de la clientèle]
\textbf{Énoncé :} En utilisant le modèle $N(t) = 500 e^{0,3t}$, déterminer au bout de combien de mois la clientèle aura doublé. Donner une valeur approchée au dixième de mois près.
\tcblower
\textbf{Solution détaillée :}
\begin{enumerate}
    \item \textbf{Traduction mathématique :} On cherche $t$ tel que $N(t) = 2 \times 500 = 1000$.
    \item \textbf{Équation :} $500 e^{0,3t} = 1000 \iff e^{0,3t} = 2$.
    \item \textbf{Résolution :} On applique le logarithme népérien (fonction réciproque de l'exponentielle) aux deux membres :
    \[
    \ln\left(e^{0,3t}\right) = \ln(2) \iff 0,3t = \ln(2)
    \]
    \item \textbf{Calcul numérique :} $\ln(2) \approx 0,6931$.
    \[
    t = \frac{\ln(2)}{0,3} \approx \frac{0,6931}{0,3} \approx 2,3
    \]
    \item \textbf{Interprétation :} La clientèle double environ tous les \textbf{2,3 mois} (environ 2 mois et 9 jours).
    \item \textbf{Vérification de cohérence :} $N(2,3) = 500 e^{0,69} \approx 500 \times 1,994 \approx 997 \approx 1000$. Le résultat est cohérent.
\end{enumerate}
\end{exoresolu}

\courssub{Modélisation de la décroissance exponentielle}

Le même formalisme s'applique quand la grandeur diminue : refroidissement d'une pièce métallique (loi de Newton), dépréciation comptable d'une machine, décharge d'un condensateur, désintégration radioactive.

\begin{exemplebox}[Dépréciation d'un tour à commande numérique]
Un atelier achète un tour CNC pour $25\,000\,000$ FCFA. Sa valeur vénale $V(t)$ (en millions de FCFA) après $t$ années suit une loi de décroissance exponentielle de coefficient $k = \ln(0,85) \approx -0,1625$ (ce qui correspond à une baisse de $15\,\%$ par an en modèle discret).
\[
V(t) = 25 \cdot e^{-0,1625 t} \quad (t \ge 0)
\]
\begin{itemize}
    \item \textbf{Valeur après 3 ans :} $V(3) = 25 e^{-0,4875} \approx 25 \times 0,614 \approx 15,35$ millions de FCFA.
    \item \textbf{Durée au bout de laquelle la valeur atteint 10\,\% de la valeur initiale (soit 2,5 millions de FCFA) :} on résout
    \[
    25 e^{-0,1625 t} = 2,5 \iff e^{-0,1625 t} = 0,1 \iff -0,1625\, t = \ln(0,1) = -\ln(10)
    \]
    \[
    t = \frac{\ln(10)}{0,1625} \approx \frac{2,3026}{0,1625} \approx 14,2 \text{ ans.}
    \]
\end{itemize}
\end{exemplebox}

\courssub{Modélisation par fonctions rationnelles : coûts moyens et seuils de rentabilité}

En gestion de production, le \textbf{coût total} $CT(x)$ pour produire $x$ unités se décompose souvent en coûts fixes $CF$ (loyer, assurance, salaires fixes) et coûts variables proportionnels à la production. Avec un coût variable unitaire constant $v$, on a $CT(x) = CF + v \cdot x$.
Le \textbf{coût moyen unitaire} est alors une fonction rationnelle (homographique) :
\[
CM(x) = \frac{CT(x)}{x} = \frac{CF}{x} + v, \quad x > 0
\]
Son étude (sens de variation, asymptotes) permet de déterminer le volume de production rentable.

\begin{popfigure}
\begin{tikzpicture}
\begin{axis}[
    width=\linewidth,
    height=8cm,
    axis lines=middle,
    xlabel={$x$ (unités produites)},
    ylabel={FCFA/unité},
    domain=1:200,
    samples=200,
    xmin=0, xmax=210,
    ymin=0, ymax=2100,
    xtick={0,50,100,150,200},
    ytick={50,300,500,1000,1500,2000},
    grid=major,
    grid style={popline},
    clip=false,
    every axis x label/.style={at={(ticklabel* cs:1.02)},anchor=west},
    every axis y label/.style={at={(ticklabel* cs:1.02)},anchor=south},
]
% Courbe CM(x) = 50 + 2000/x
\addplot[popcurve, thick, domain=1:200] {50 + 2000/x};
% Asymptote horizontale y = 50
\addplot[poporange, dashed, thick, domain=0:210] {50} node[pos=0.75, above, poporange, font=\small] {Asymptote $y=50$};
% Prix de vente P = 300
\addplot[popgreen, dashed, thick, domain=0:210] {300} node[pos=0.6, above, popgreen, font=\small] {Prix de vente $P=300$};
% Intersection CM(x) = 300 -> x = 8
\draw[poppink, dashed, thick] (8,0) -- (8,300);
\node[poppink, anchor=north west, font=\small] at (8,0) {$x=8$ (seuil)};
% Point (10 ; 250)
\addplot[only marks, mark=*, mark size=2pt, popblue] coordinates {(10,250)} node[above right, font=\small] {$(10\,; 250)$};
\end{axis}
\end{tikzpicture}
\legende{Coût moyen $CM(x) = 50 + \dfrac{2000}{x}$ (coûts fixes 2\,000 FCFA, coût variable 50 FCFA/unité). Le coût moyen décroît et se rapproche de l'asymptote $y=50$. Seuil de rentabilité $x=8$ pour un prix de vente de 300 FCFA.}
\end{popfigure}

\begin{exemplebox}[Analyse de rentabilité]
Soit $CM(x) = 50 + \dfrac{2000}{x}$ (coûts fixes $2\,000$ FCFA, coût variable $50$ FCFA/unité), définie pour $x > 0$. Le prix de vente unitaire est fixé à $P = 300$ FCFA.
\begin{enumerate}
    \item \textbf{Seuil de rentabilité (point mort) :} on résout $CM(x) = P$.
    \[
    50 + \frac{2000}{x} = 300 \iff \frac{2000}{x} = 250 \iff x = \frac{2000}{250} = 8
    \]
    Il faut produire et vendre \textbf{au moins 8 unités} pour couvrir les coûts (à 8 unités, le bénéfice est nul ; la rentabilité stricte commence à 9 unités).
    \item \textbf{Bénéfice unitaire :} $B_u(x) = P - CM(x) = 250 - \dfrac{2000}{x}$.
    \item \textbf{Bénéfice total :} $B_T(x) = x \cdot B_u(x) = 250x - 2000$. C'est une fonction affine strictement croissante.
    \item \textbf{Interprétation de l'asymptote :} quand $x \to +\infty$, $\dfrac{2000}{x} \to 0$, donc $CM(x) \to 50$. Le coût moyen se rapproche du coût variable unitaire sans jamais descendre en dessous : la droite $y=50$ est asymptote horizontale. Le bénéfice unitaire se rapproche alors de $250$ FCFA.
\end{enumerate}
\end{exemplebox}

\courssub{Détermination de seuils : résolution d'inéquations et exploitation graphique}

Dans un contexte technique, on cherche souvent un \textbf{seuil} : « À partir de quand le bénéfice dépasse-t-il 5\,000 FCFA ? », « Combien de temps avant que la température ne descende sous 50 °C ? ». Cela revient à résoudre une inéquation faisant intervenir $\ln$ ou $\exp$.

\begin{methode}[Modéliser et résoudre un problème concret en 4 étapes]
\begin{enumerate}
    \item \textbf{Variables et fonction :} identifier la variable (souvent le temps $t$ ou la quantité $x$), définir la fonction qui modélise la situation et \textbf{préciser son domaine de définition} (souvent $x>0$ ou $t \ge 0$).
    \item \textbf{Traduction mathématique :} écrire l'équation ou l'inéquation correspondant à la question posée (ex : $N(t) \ge 10\,000$, $CM(x) \le P$).
    \item \textbf{Résolution algébrique :} résoudre en utilisant la stricte croissance de $\ln$ et de $\exp$ sur leurs domaines (le sens de l'inégalité est conservé quand on applique $\ln$ ou $\exp$ ; il est \textbf{inversé} si l'on divise par un nombre négatif). Vérifier que la solution appartient au domaine de définition.
    \item \textbf{Interprétation et conclusion :} revenir au langage courant, donner l'unité, arrondir de manière \textbf{cohérente} (à l'entier supérieur pour un délai d'attente, à l'entier inférieur pour une quantité maximale) et \textbf{commenter le résultat}.
\end{enumerate}
\end{methode}

\begin{attention}[Piège majeur : la réponse « mathématique » sans contexte]
\textbf{Erreur fréquente :} l'élève écrit « $t = \dfrac{\ln(20)}{0,3} \approx 9,98$ » et s'arrête là.
\textbf{Correction :} au Probatoire, la conclusion \textbf{doit} être rédigée : « La clientèle atteindra 10\,000 clients au bout de \textbf{10 mois} (arrondi à l'entier supérieur, car au bout de 9 mois le seuil n'est pas encore atteint). »
Autres pièges :
\begin{itemize}
    \item Oublier l'unité (mois, années, FCFA, °C, unités produites).
    \item Arrondir $9,98$ à $9$ mois : le seuil ne serait pas atteint.
    \item Oublier la condition d'existence : $ax+b > 0$ pour $\ln(ax+b)$.
    \item Confondre un coefficient de croissance continue $k = 0,3$ avec un taux discret de $30\,\%$ (qui correspond à $k = \ln(1,3) \approx 0,262$).
    \item Oublier d'inverser le sens de l'inégalité lors d'une division par un nombre négatif (ex : $-0,05t < -2,28 \iff t > 45,5$).
\end{itemize}
\end{attention}

\begin{exoresolu}[Seuil de température : refroidissement d'une pièce]
\textbf{Énoncé :} Une pièce en acier sort du four à $800$ °C. Elle refroidit dans un atelier à $20$ °C selon la loi : $T(t) = 20 + 780 e^{-0,05 t}$ ($t$ en minutes). Au bout de combien de temps la température descend-elle \textbf{sous} $100$ °C (seuil de manipulation sans gants) ?
\tcblower
\textbf{Solution détaillée :}
\begin{enumerate}
    \item \textbf{Fonction :} $T(t) = 20 + 780 e^{-0,05 t}$, définie sur $[0\,; +\infty[$. Comme $T'(t) = -39 e^{-0,05t} < 0$, $T$ est strictement décroissante : l'inéquation $T(t) < 100$ aura une solution de la forme $t > t_0$.
    \item \textbf{Inéquation :} $T(t) < 100 \iff 20 + 780 e^{-0,05 t} < 100$.
    \item \textbf{Résolution :}
    \begin{align*}
    780 e^{-0,05 t} &< 80 \\
    e^{-0,05 t} &< \frac{80}{780} = \frac{4}{39} \\
    -0,05 t &< \ln\left(\frac{4}{39}\right) \quad \text{(ln est croissante : sens conservé)} \\
    -0,05 t &< \ln(4) - \ln(39) \approx 1,3863 - 3,6636 \approx -2,2773 \\
    t &> \frac{-2,2773}{-0,05} \quad \text{(division par un négatif : sens inversé)} \\
    t &> 45,5
    \end{align*}
    \item \textbf{Interprétation :} la température passe sous $100$ °C après environ \textbf{45,5 minutes}. En pratique, l'opérateur attendra \textbf{46 minutes} (arrondi à la minute supérieure, par sécurité).
    \item \textbf{Exploitation graphique :} sur la courbe de $T$, on trace la droite horizontale $y=100$. L'abscisse du point d'intersection donne le seuil ; la courbe étant décroissante, la zone $T < 100$ se situe à \textbf{droite} de ce point.
\end{enumerate}
\end{exoresolu}

\courssub{Critique du modèle et rédaction de la conclusion argumentée}

Un modèle mathématique est une \emph{simplification} de la réalité. Au Probatoire, on attend de l'élève qu'il soit capable de \textbf{discuter la validité} du modèle utilisé.

\begin{aretenir}[Rédiger une conclusion argumentée : structure type]
Une bonne conclusion technique contient :
\begin{enumerate}
    \item \textbf{La réponse chiffrée} (avec unité et arrondi justifié).
    \item \textbf{La référence au modèle} utilisé (ex : « Selon le modèle exponentiel $N(t)=500e^{0,3t}$... »).
    \item \textbf{Les réserves / limites} (hypothèses non vérifiées dans la réalité).
    \item \textbf{Une phrase de synthèse} en langage courant, adressée au décideur (chef d'atelier, client, banquier).
\end{enumerate}
\end{aretenir}

\begin{savaistu}[Limites de la modélisation exponentielle en entreprise]
Une croissance exponentielle $N(t) = N_0 e^{kt}$ suppose des ressources illimitées (espace, matières premières, main-d'œuvre, marché). Dans la réalité, la croissance ralentit à l'approche d'une \textbf{capacité maximale} (saturation du marché, limite de production de l'atelier) : la courbe réelle prend une forme en « S ». De même, la loi de refroidissement utilisée plus haut n'est qu'une approximation valable pour des écarts de température modérés. Un bon technicien sait donc toujours \emph{dater} et \emph{réviser} ses prévisions.
\end{savaistu}

\begin{exemplebox}[Conclusion type pour l'exemple de la clientèle]
\textbf{Conclusion :} Selon le modèle de croissance exponentielle continue de coefficient $0,3$ par mois ($N(t)=500e^{0,3t}$), l'entreprise atteindra le seuil de 10\,000 clients au bout de \textbf{10 mois}.

\emph{Réserves :} ce modèle suppose un marché illimité et une capacité de traitement sans contrainte. En réalité, la saturation du marché de Douala ou la limite de l'effectif des techniciens fera probablement ralentir la croissance avant cette échéance. Il est recommandé de réviser les prévisions chaque trimestre.
\end{exemplebox}

\begin{exercice}[Entraînement : Investissement et seuils]
Une PME du secteur bois (menuiserie industrielle) à Yaoundé souhaite acquérir une nouvelle scie numérique coûtant $15\,000\,000$ FCFA.
\begin{enumerate}
    \item L'amortissement est linéaire sur 5 ans (valeur résiduelle nulle). Écrire la fonction $V(t)$ donnant la valeur comptable (en millions de FCFA) après $t$ années ($0 \le t \le 5$). Tracer sa représentation graphique.
    \item La machine permet un gain de productivité estimé à $4\,000\,000$ FCFA par an. Au bout de combien d'années l'investissement est-il rentabilisé (somme des gains $\ge$ coût initial) ?
    \item On envisage un modèle de dépréciation exponentielle pour la valeur de revente : $V_{exp}(t) = 15 e^{-0,223 t}$ (en millions de FCFA). Comparer les deux modèles à $t=3$ ans. Lequel donne la valeur de revente la plus élevée ?
    \item Rédiger une conclusion argumentée à l'attention du directeur financier.
\end{enumerate}
\end{exercice}

\begin{corrige}[Exercice : Investissement et seuils]
\begin{enumerate}
    \item \textbf{Amortissement linéaire :} perte de valeur annuelle $= \dfrac{15}{5} = 3$ millions de FCFA par an.
    \[
    V(t) = 15 - 3t \quad \text{(en millions de FCFA)}, \quad t \in [0\,;5]
    \]
    La représentation graphique est le segment reliant les points $(0\,;15)$ et $(5\,;0)$.
    \item \textbf{Rentabilité :} le gain cumulé après $t$ années est $G(t) = 4t$ (en millions). On résout :
    \[
    G(t) \ge 15 \iff 4t \ge 15 \iff t \ge 3,75
    \]
    L'investissement est rentabilisé au bout de \textbf{3,75 ans}, soit \textbf{3 ans et 9 mois}.
    \item \textbf{Comparaison à $t=3$ :}
    \begin{itemize}
        \item Modèle linéaire : $V(3) = 15 - 3 \times 3 = 6$ millions de FCFA.
        \item Modèle exponentiel : $V_{exp}(3) = 15 e^{-0,223 \times 3} = 15 e^{-0,669} \approx 15 \times 0,512 \approx 7,68$ millions de FCFA.
    \end{itemize}
    Le modèle exponentiel donne une valeur de revente \textbf{plus élevée} (7,68 M contre 6 M de FCFA) : il traduit une décote moins sévère les premières années.
    \item \textbf{Conclusion :} Selon l'amortissement linéaire comptable, la scie est rentabilisée en 3 ans et 9 mois, ce qui est inférieur à sa durée d'amortissement (5 ans) : l'investissement est donc rentable. À 3 ans, la valeur comptable est de 6 M FCFA, alors que la valeur de marché estimée (modèle exponentiel) est d'environ 7,68 M FCFA : en cas de revente à cette date, l'entreprise réaliserait une plus-value d'environ 1,68 M FCFA par rapport à la valeur comptable. Le directeur financier doit toutefois noter que la valeur de marché réelle dépend de l'état de la machine et de la demande ; ces estimations sont à réviser chaque année.
\end{enumerate}
\end{corrige}

\newpage

\coursec{Activité d'intégration}

\courssub{Objectifs de l'activité}

À la fin de cette activité, tu seras capable de :
\begin{itemize}
\item \cle{modéliser} une situation économique réelle à l'aide des fonctions usuelles (exponentielle, logarithme, fonction homographique) ;
\item étudier le \cle{sens de variation} d'une fonction exponentielle et d'une fonction rationnelle ;
\item résoudre une \cle{équation} et une \cle{inéquation} faisant intervenir $\ln$ et $\exp$ ;
\item déterminer une \cle{asymptote horizontale} et interpréter sa signification concrète ;
\item \cle{rédiger et justifier} une démarche complète sous forme d'un rapport de conseil, comme le ferait un technicien en gestion.
\end{itemize}

\courssub{Situation}

Mama Ngo Bell tient un dépôt de boissons bien connu au marché Mokolo à Yaoundé. Depuis qu'elle a ouvert une page sur les réseaux sociaux pour prendre les commandes des bars et des buvettes du quartier, sa clientèle ne cesse de grandir. Son neveu, élève en Terminale technique, a relevé les chiffres et propose les modèles suivants.

\textbf{Évolution de la clientèle.} Le nombre de clients réguliers du dépôt, $t$ mois après le lancement de la page, est modélisé par :
\[ N(t) = 400\,e^{0{,}25t} \qquad (t \geq 0). \]
Mama Ngo Bell estime qu'au-delà de \num{10000} clients réguliers, elle ne pourra plus tout gérer seule : il lui faudra embaucher un employé supplémentaire.

\textbf{Coût de stockage.} Le dépôt stocke les caisses de boissons dans un hangar loué à côté. Pour $x$ caisses stockées (avec $x \geq 100$), le coût moyen de stockage par caisse, en francs CFA, est donné par :
\[ C(x) = 120 + \frac{36000}{x}. \]
Pour que l'activité reste rentable, Mama Ngo Bell s'est fixé une règle : le coût moyen de stockage par caisse ne doit pas dépasser $200$ FCFA.

\textbf{Problème posé :} dans combien de mois Mama Ngo Bell devra-t-elle embaucher un employé supplémentaire ? Quel volume minimal de stockage garantit-il la rentabilité du dépôt ?

\courssub{Consignes}

Le travail se fait par groupes de trois ou quatre élèves. Chaque groupe rédigera ses réponses soigneusement, en justifiant chaque étape par une propriété du cours.

\begin{enumerate}
\item \textbf{Croissance de la clientèle.}
\begin{enumerate}
\item Calculer $N(0)$. Que représente concrètement ce nombre ?
\item Justifier, en utilisant la dérivée, que la fonction $N$ est strictement croissante sur $[0\,;\,+\infty[$.
\item Calculer la clientèle après $6$ mois. Arrondir à l'unité.
\end{enumerate}
\item \textbf{Le mois de l'embauche.} Résoudre dans $[0\,;\,+\infty[$ l'équation $N(t) = \num{10000}$, puis l'inéquation $N(t) > \num{10000}$. En déduire à partir de quel mois la clientèle dépassera \num{10000} personnes.
\item \textbf{Étude du coût moyen de stockage.}
\begin{enumerate}
\item Montrer que pour tout $x \geq 100$ : $C'(x) = -\dfrac{36000}{x^2}$, et en déduire le sens de variation de $C$ sur $[100\,;\,+\infty[$.
\item Calculer $\displaystyle\lim_{x \to +\infty} C(x)$ et interpréter graphiquement le résultat (asymptote). Que signifie-t-il concrètement pour Mama Ngo Bell ?
\item Dresser le tableau de variations complet de $C$ sur $[100\,;\,+\infty[$.
\end{enumerate}
\item \textbf{Le volume rentable.} Résoudre dans $[100\,;\,+\infty[$ l'inéquation $C(x) \leq 200$. En déduire le nombre minimal de caisses que le dépôt doit stocker pour respecter la règle de rentabilité.
\item \textbf{Rapport de conseil.} Rédiger un court rapport (10 à 15 lignes) adressé à Mama Ngo Bell, dans lequel tu lui conseilles le mois d'embauche de l'employé supplémentaire et le volume de stockage minimal à prévoir. Chaque conseil doit être justifié par un résultat mathématique obtenu dans les questions précédentes.
\end{enumerate}

\courssub{Ressources à mobiliser}

\begin{aretenir}[Boîte à outils de la leçon]
\begin{itemize}
\item \textbf{Dérivées :} $(e^{ax+b})' = a\,e^{ax+b}$ ; $\left(\dfrac{1}{x}\right)' = -\dfrac{1}{x^2}$ ; $(u+v)' = u'+v'$ ; $(k\,u)' = k\,u'$ ($k$ constante).
\item \textbf{Signe de l'exponentielle :} pour tout réel $u$, $e^{u} > 0$.
\item \textbf{Équations et inéquations :} $e^{A} = B \iff A = \ln B$ (avec $B>0$) ; $\ln A = B \iff A = e^{B}$ (avec $A>0$). Les fonctions $\ln$ et $\exp$ sont \textbf{strictement croissantes} : elles conservent le sens des inégalités.
\item \textbf{Limite :} $\displaystyle\lim_{x \to +\infty} \frac{k}{x} = 0$ ($k$ constante). Si $\displaystyle\lim_{x \to +\infty} f(x) = \ell$, la droite d'équation $y = \ell$ est asymptote horizontale à la courbe de $f$ en $+\infty$.
\item \textbf{Sens de variation :} le signe de $f'$ donne les variations de $f$ : si $f' > 0$ sur un intervalle, $f$ y est strictement croissante ; si $f' < 0$, $f$ y est strictement décroissante.
\item \textbf{Inéquations :} multiplier ou diviser les deux membres par un nombre \textbf{strictement positif} conserve le sens de l'inégalité.
\end{itemize}
\end{aretenir}

\courssub{Démarche et corrigé commenté}

\begin{exoresolu}[Corrigé complet de l'activité]
\textbf{Question 1.} a) Calculer $N(0)$. b) Justifier que $N$ est strictement croissante. c) Calculer $N(6)$.

\textbf{Question 2.} Résoudre $N(t) = \num{10000}$ puis $N(t) > \num{10000}$.

\textbf{Question 3.} a) Calculer $C'(x)$ et conclure sur les variations. b) Limite en $+\infty$ et asymptote. c) Tableau de variations.

\textbf{Question 4.} Résoudre $C(x) \leq 200$ sur $[100\,;\,+\infty[$.

\textbf{Question 5.} Rédiger le rapport de conseil.
\tcblower
\textbf{Question 1 — Croissance de la clientèle.}

a) $N(0) = 400\,e^{0} = 400 \times 1 = 400$. Au lancement de la page, le dépôt comptait $400$ clients réguliers.

b) $N$ est de la forme $t \mapsto k\,e^{at}$ avec $k = 400$ et $a = 0{,}25$ ; elle est donc dérivable sur $[0\,;\,+\infty[$ et :
\[ N'(t) = 400 \times 0{,}25\,e^{0{,}25t} = 100\,e^{0{,}25t}. \]
Or pour tout réel $t$, $e^{0{,}25t} > 0$, donc $N'(t) = 100\,e^{0{,}25t} > 0$. La fonction $N$ est donc \textbf{strictement croissante} sur $[0\,;\,+\infty[$ : la clientèle augmente continuellement.

c) $N(6) = 400\,e^{0{,}25 \times 6} = 400\,e^{1{,}5} \approx 400 \times 4{,}4817 \approx \num{1793}$. Après $6$ mois, le dépôt compte environ \num{1793} clients réguliers.

\textbf{Question 2 — Le mois de l'embauche.}

On résout l'équation $N(t) = \num{10000}$ :
\begin{align*}
400\,e^{0{,}25t} &= \num{10000}\\
e^{0{,}25t} &= \frac{\num{10000}}{400} = 25\\
0{,}25t &= \ln 25 \quad (\text{car } e^{A} = B \iff A = \ln B, \text{ avec } B = 25 > 0)\\
t &= \frac{\ln 25}{0{,}25} = 4\ln 25 \approx 4 \times 3{,}2189 \approx 12{,}88.
\end{align*}
Comme $N$ est strictement croissante sur $[0\,;\,+\infty[$, elle conserve l'ordre : l'inéquation $N(t) > \num{10000}$ équivaut à $t > 4\ln 25 \approx 12{,}88$. La clientèle dépassera donc \num{10000} personnes au cours du \textbf{13\textsuperscript{e} mois}.

\textbf{Question 3 — Étude du coût moyen.}

a) On écrit $C(x) = 120 + 36000 \times \dfrac{1}{x}$. La fonction $x \mapsto \dfrac{1}{x}$ est dérivable sur $[100\,;\,+\infty[$ (qui ne contient pas $0$), donc $C$ est dérivable sur $[100\,;\,+\infty[$ et :
\[ C'(x) = 0 + 36000 \times \left(-\frac{1}{x^2}\right) = -\frac{36000}{x^2}. \]
Pour tout $x \geq 100$, $x^2 > 0$, donc $C'(x) < 0$ : la fonction $C$ est \textbf{strictement décroissante} sur $[100\,;\,+\infty[$. Plus on stocke de caisses, plus le coût moyen par caisse diminue.

b) $\displaystyle\lim_{x \to +\infty} \frac{36000}{x} = 0$, donc $\displaystyle\lim_{x \to +\infty} C(x) = 120$. La droite d'équation $y = 120$ est \textbf{asymptote horizontale} à la courbe de $C$ en $+\infty$. Concrètement : le coût moyen par caisse ne pourra jamais descendre en dessous de $120$ FCFA, qui représente le coût fixe incompressible (loyer du hangar, gardiennage).

c) On calcule $C(100) = 120 + \dfrac{36000}{100} = 120 + 360 = 480$ FCFA. Tableau de variations :

\begin{center}
\begin{tabular}{|c|ccc|}\hline
$x$ & $100$ & & $+\infty$ \\ \hline
$C'(x)$ & & $-$ & \\ \hline
$C$ & $480$ & $\searrow$ & $120$ \\ \hline
\end{tabular}
\end{center}

\textbf{Question 4 — Le volume rentable.}

On résout dans $[100\,;\,+\infty[$ :
\begin{align*}
C(x) \leq 200 &\iff 120 + \frac{36000}{x} \leq 200\\
&\iff \frac{36000}{x} \leq 80\\
&\iff 36000 \leq 80x \quad (\text{car } x \geq 100 > 0 : \text{ multiplier par } x \text{ conserve le sens})\\
&\iff x \geq \frac{36000}{80} = 450.
\end{align*}
Comme $450 \geq 100$, la solution est compatible avec le domaine : l'ensemble des solutions est $[450\,;\,+\infty[$. Le dépôt doit donc stocker \textbf{au moins $450$ caisses} pour que le coût moyen par caisse ne dépasse pas $200$ FCFA. Vérification : $C(450) = 120 + \dfrac{36000}{450} = 120 + 80 = 200$ FCFA.

\textbf{Question 5 — Rapport de conseil (modèle de rédaction).}

\emph{À l'attention de Mama Ngo Bell, dépôt de boissons, marché Mokolo.}

\emph{Madame, l'étude mathématique de votre activité conduit aux recommandations suivantes. Votre clientèle, modélisée par $N(t) = 400\,e^{0{,}25t}$, est strictement croissante car $N'(t) = 100\,e^{0{,}25t} > 0$ pour tout $t \geq 0$. La résolution de l'équation $N(t) = \num{10000}$ donne $t = 4\ln 25 \approx 12{,}88$ mois : nous vous conseillons donc de recruter un employé supplémentaire \textbf{au début du 13\textsuperscript{e} mois}, moment où votre clientèle dépassera \num{10000} personnes. Concernant le stockage, le coût moyen par caisse $C(x) = 120 + \frac{36000}{x}$ est strictement décroissant et ne peut descendre sous $120$ FCFA (asymptote horizontale d'équation $y = 120$). L'inéquation $C(x) \leq 200$ montre qu'il faut stocker \textbf{au minimum $450$ caisses} pour respecter votre seuil de rentabilité de $200$ FCFA par caisse. Nous vous recommandons donc d'aménager le hangar pour un volume d'au moins $450$ caisses.}
\end{exoresolu}

\begin{attention}[Pièges fréquents dans cette activité]
\begin{itemize}
\item Ne pas oublier que $e^{0{,}25t} = 25$ se résout par $0{,}25t = \ln 25$ et non $t = \dfrac{25}{0{,}25}$ : on ne « descend » pas l'exposant sans le logarithme.
\item Dans l'inéquation $\dfrac{36000}{x} \leq 80$, on multiplie par $x$ qui est \textbf{strictement positif} (ce sont des caisses, $x \geq 100$) : le sens de l'inégalité est conservé. Le justifier dans la rédaction.
\item Le résultat $t \approx 12{,}88$ doit être \textbf{interprété} : on répond « au cours du 13\textsuperscript{e} mois », pas « au bout de $12{,}88$ mois » dans le rapport.
\item Ne pas confondre le coût moyen $C(x)$ (par caisse) avec un coût total.
\item Dans le tableau de variations, la valeur $120$ en $+\infty$ est une \textbf{limite} : elle n'est jamais atteinte, la courbe s'en approche sans la toucher.
\end{itemize}
\end{attention}

\courssub{Bilan de l'activité}

Cette activité a permis de mobiliser l'ensemble des savoir-faire de la leçon dans un contexte économique authentique :
\begin{itemize}
\item l'\textbf{exponentielle népérienne} modélise une croissance rapide (clientèle) et sa dérivée $a\,e^{ax+b}$, toujours du signe de $a$, justifie la monotonie ;
\item le \textbf{logarithme népérien} est l'outil indispensable pour « extraire » l'inconnue d'un exposant : résoudre $e^{at} = k$ (avec $k > 0$) revient à calculer $t = \dfrac{\ln k}{a}$ ;
\item la \textbf{fonction homographique} $x \mapsto a + \dfrac{b}{x}$ modélise un coût moyen : elle est décroissante quand $b > 0$ et admet l'asymptote horizontale d'équation $y = a$, qui s'interprète comme un coût plancher ;
\item la résolution d'\textbf{équations et d'inéquations} avec ces fonctions fournit des seuils de décision concrets (mois d'embauche, volume minimal) ;
\item enfin, la \textbf{rédaction d'un rapport justifié} montre que les mathématiques de Terminale technique sont un véritable outil d'aide à la décision pour un commerçant ou un gestionnaire.
\end{itemize}

\begin{savaistu}[Les mathématiques au service des commerçants]
Les modèles de croissance exponentielle sont utilisés chaque jour par les banques et les opérateurs de mobile money au Cameroun pour prévoir l'évolution du nombre de clients. De même, la notion de coût moyen décroissant (les « économies d'échelle ») explique pourquoi les grossistes du marché Mokolo ou de Douala proposent des prix dégressifs : plus le volume stocké est grand, plus le coût par unité diminue, jusqu'à un plancher incompressible.
\end{savaistu}

\newpage

\coursec{Méthodes et exercices résolus}

\courssub{Méthodes et exercices résolus}

\begin{methode}[Étude complète d'une fonction rationnelle / homographique]
\textbf{Objectif :} Déterminer le domaine, les limites, la dérivée, le tableau de variations, les asymptotes et tracer la courbe.
\begin{enumerate}
    \item \textbf{Domaine de définition :} Identifier les valeurs interdites (zéros du dénominateur pour une fraction rationnelle). $D_f = \mathbb{R} \setminus \{ \text{zéros du dénominateur} \}$.
    \item \textbf{Parité / symétrie :} Vérifier si $f(-x)=f(x)$ (paire) ou $f(-x)=-f(x)$ (impaire). Pour une homographique $f(x)=\dfrac{ax+b}{cx+d}$ avec $c\neq 0$, la courbe admet le centre de symétrie $I\left(-\dfrac{d}{c}\,;\,\dfrac{a}{c}\right)$.
    \item \textbf{Limites aux bornes du domaine :} Calculer $\lim\limits_{x\to \pm\infty}f(x)$ et les limites aux valeurs interdites. En déduire les asymptotes :
    \begin{itemize}
        \item Horizontale $y=\ell$ si la limite en $\pm\infty$ est un nombre fini $\ell$.
        \item Verticale $x=x_0$ si la limite en $x_0$ est infinie.
    \end{itemize}
    \item \textbf{Dérivée :} Calculer $f'(x)$ (formule du quotient). Pour une homographique, on peut aussi écrire $f(x)=\dfrac{a}{c}+\dfrac{bc-ad}{c(cx+d)}$ par division euclidienne, ce qui simplifie la dérivation.
    \item \textbf{Signe de la dérivée et tableau de variations :} Résoudre $f'(x)>0$ et $f'(x)<0$ sur $D_f$. Construire le tableau en séparant les valeurs interdites par une double barre.
    \item \textbf{Points remarquables :} Intersections avec les axes ($f(0)$, $f(x)=0$), centre de symétrie.
    \item \textbf{Tracé de la courbe :} Repère orthonormé, asymptotes en pointillés, courbe respectant les variations.
\end{enumerate}
\end{methode}

\begin{exoresolu}[Étude de $f(x)=\dfrac{2x-1}{x+3}$ et tracé de $\mathcal{C}_f$]
Soit la fonction $f$ définie sur $\mathbb{R}\setminus\{-3\}$ par $f(x)=\dfrac{2x-1}{x+3}$.
\begin{enumerate}
    \item Étudier les variations de $f$ et dresser son tableau de variations.
    \item Déterminer les asymptotes de la courbe $\mathcal{C}_f$.
    \item Montrer que $\mathcal{C}_f$ admet un centre de symétrie $I$ et donner ses coordonnées.
    \item Tracer $\mathcal{C}_f$ dans un repère orthonormé (unité : 2 cm).
\end{enumerate}
\tcblower
\textbf{Solution détaillée}

\textbf{1. Domaine et dérivée}

Le dénominateur s'annule en $x=-3$, donc $D_f = \mathbb{R}\setminus\{-3\}$.

$f$ est un quotient de fonctions dérivables sur $D_f$, avec $u(x)=2x-1$, $u'(x)=2$, $v(x)=x+3$, $v'(x)=1$ :
\[ f'(x) = \frac{u'v-uv'}{v^2} = \frac{2(x+3) - (2x-1)\times 1}{(x+3)^2} = \frac{2x+6-2x+1}{(x+3)^2} = \frac{7}{(x+3)^2}. \]
Pour tout $x\in D_f$, $(x+3)^2 > 0$, donc $f'(x) > 0$.

$f$ est \textbf{strictement croissante} sur $]-\infty;-3[$ et sur $]-3;+\infty[$.

\begin{attention}[Croissance sur chaque intervalle]
On dit que $f$ est croissante \emph{sur chacun} des intervalles $]-\infty;-3[$ et $]-3;+\infty[$, mais \textbf{jamais} sur leur réunion : par exemple $f(-4)=7$ et $f(0)=-\tfrac13$, donc $-4<0$ mais $f(-4)>f(0)$.
\end{attention}

\textbf{2. Limites et asymptotes}
\begin{itemize}
    \item En l'infini, on conserve les termes de plus haut degré : $\lim\limits_{x\to \pm\infty}f(x) = \lim\limits_{x\to \pm\infty}\dfrac{2x}{x} = 2$. Donc la droite $\mathcal{H} : y = 2$ est \textbf{asymptote horizontale} à $\mathcal{C}_f$.
    \item En $-3$ : le numérateur vaut $2\times(-3)-1=-7\neq 0$. $\lim\limits_{x\to -3^-}f(x) = +\infty$ (car $x+3\to 0^-$) et $\lim\limits_{x\to -3^+}f(x) = -\infty$ (car $x+3\to 0^+$). Donc la droite $\mathcal{V} : x = -3$ est \textbf{asymptote verticale}.
\end{itemize}

\textbf{3. Centre de symétrie}

Pour une homographique $f(x)=\dfrac{ax+b}{cx+d}$, le centre de symétrie est $I\left(-\dfrac{d}{c}\,;\,\dfrac{a}{c}\right)$. Ici $a=2$, $b=-1$, $c=1$, $d=3$, donc $I(-3\,;\,2)$.

\emph{Vérification :} pour tout $h\neq 0$,
\[ f(-3+h)+f(-3-h) = \frac{2(-3+h)-1}{h} + \frac{2(-3-h)-1}{-h} = \frac{-7+2h}{h} + \frac{7+2h}{h} = \frac{4h}{h} = 4 = 2\times 2. \]
Donc $\dfrac{f(-3+h)+f(-3-h)}{2}=2$ : le point $I(-3;2)$ est bien centre de symétrie de $\mathcal{C}_f$.

\textbf{4. Tableau de variations}
\begin{center}
\begin{tabular}{|c|cccccc|}\hline
$x$ & $-\infty$ & & $-3$ & & & $+\infty$ \\ \hline
$f'(x)$ & & $+$ & $\|$ & $\|$ & $+$ & \\ \hline
$f$ & $2$ & $\nearrow$ & $+\infty$ & $-\infty$ & $\nearrow$ & $2$ \\ \hline
\end{tabular}
\end{center}
(La double barre $\|$ matérialise la valeur interdite $-3$.)

\textbf{5. Points particuliers}

$f(0) = -\dfrac{1}{3} \approx -0,33$ : point $A\left(0\,;\,-\tfrac13\right)$.

$f(x)=0 \iff 2x-1=0 \iff x=0,5$ : point $B(0,5\,;\,0)$.

\textbf{6. Tracé}

Repère orthonormé $(O;\vec{\imath},\vec{\jmath})$, unité 2 cm. Tracer $\mathcal{V}: x=-3$ et $\mathcal{H}: y=2$ en pointillés, marquer $I(-3;2)$, $A$, $B$, puis tracer les deux branches de l'hyperbole, croissantes, symétriques par rapport à $I$, et s'approchant des asymptotes.

\begin{popfigure}
\begin{tikzpicture}[scale=0.55]
\draw[popline,->] (-7,0) -- (4,0) node[below left]{$x$};
\draw[popline,->] (0,-4) -- (0,7) node[below left]{$y$};
\draw[popmuted,dashed] (-3,-4) -- (-3,7) node[above]{$\mathcal{V}$};
\draw[popmuted,dashed] (-7,2) -- (4,2) node[right]{$\mathcal{H}$};
\draw[popblue,thick,domain=-6.8:-3.9,samples=100] plot (\x,{(2*\x-1)/(\x+3)});
\draw[popblue,thick,domain=-2.1:3.8,samples=100] plot (\x,{(2*\x-1)/(\x+3)});
\fill[poporange] (-3,2) circle (2.5pt) node[above right]{$I$};
\fill[poporange] (0,-0.333) circle (2.5pt) node[below right]{$A$};
\fill[poporange] (0.5,0) circle (2.5pt) node[above right]{$B$};
\end{tikzpicture}
\legende{Courbe de $f(x)=\dfrac{2x-1}{x+3}$, ses asymptotes et son centre de symétrie.}
\end{popfigure}
\end{exoresolu}

\begin{methode}[Résolution d'équations faisant intervenir $\ln$ et $\exp$]
\textbf{Principe :} $\ln$ et $\exp$ sont des bijections réciproques : $\exp(\ln x)=x$ pour $x>0$ et $\ln(\exp x)=x$ pour tout $x\in\mathbb{R}$.
\begin{enumerate}
    \item \textbf{Conditions d'existence (C.E.) :} pour tout $\ln(u(x))$, imposer $u(x) > 0$. Noter l'ensemble de définition $\mathcal{D}$.
    \item \textbf{Simplification :} utiliser les propriétés algébriques :
    $\ln(ab)=\ln a+\ln b$ ($a>0$, $b>0$), $\ln\dfrac{a}{b}=\ln a-\ln b$, $\ln(a^n)=n\ln a$ ($a>0$) ;
    $\exp(a+b)=\exp a \times \exp b$, $\exp(a-b)=\dfrac{\exp a}{\exp b}$.
    \item \textbf{Passer à l'exponentielle (si $\ln$) ou au logarithme (si $\exp$) :}
    \begin{itemize}
        \item $\ln A = \ln B \iff A=B$ (avec $A>0$, $B>0$) ;
        \item $\ln A = k \iff A = e^k$ ;
        \item $\exp A = \exp B \iff A=B$ ;
        \item $\exp A = k$ avec $k>0$ $\iff A = \ln k$ ; si $k\le 0$, l'équation $\exp A = k$ n'a \textbf{aucune solution} car $\exp A > 0$.
    \end{itemize}
    \item \textbf{Résoudre l'équation algébrique obtenue.}
    \item \textbf{Vérifier l'appartenance à $\mathcal{D}$ :} éliminer les solutions parasites.
    \item \textbf{Conclusion :} écrire l'ensemble des solutions $\mathcal{S}$.
\end{enumerate}
\end{methode}

\begin{exoresolu}[Équations mixtes $\ln$ et $\exp$]
Résoudre dans $\mathbb{R}$ les équations suivantes :
\begin{enumerate}
    \item $\ln(2x-1) + \ln(x+2) = \ln(15)$
    \item $e^{2x} - 3e^x + 2 = 0$
    \item $\ln(x+1) = 2 - \ln(x-1)$
\end{enumerate}
\tcblower
\textbf{Solution détaillée}

\textbf{Équation 1 : $\ln(2x-1) + \ln(x+2) = \ln(15)$}
\begin{itemize}
    \item \textbf{C.E. :} $\begin{cases} 2x-1 > 0 \\ x+2 > 0 \end{cases} \iff \begin{cases} x > 1/2 \\ x > -2 \end{cases} \iff x > \dfrac12$. Donc $\mathcal{D} = \left]\dfrac12\,;\, +\infty\right[$.
    \item \textbf{Transformation :} $\ln\big[(2x-1)(x+2)\big] = \ln(15)$.
    \item \textbf{Équivalence :} $(2x-1)(x+2) = 15$.
    \item \textbf{Résolution :} $2x^2 + 4x - x - 2 = 15 \iff 2x^2 + 3x - 17 = 0$.
    $\Delta = 3^2 - 4\times 2\times(-17) = 9 + 136 = 145 > 0$.
    $x_1 = \dfrac{-3-\sqrt{145}}{4} \approx -3,76$ ; $x_2 = \dfrac{-3+\sqrt{145}}{4} \approx 2,26$.
    \item \textbf{Vérification des C.E. :} $x_1 \approx -3,76 \notin \mathcal{D}$ (rejeté) ; $x_2 \approx 2,26 \in \mathcal{D}$ (conservé).
\end{itemize}
$\mathcal{S} = \left\{ \dfrac{-3+\sqrt{145}}{4} \right\}$.

\textbf{Équation 2 : $e^{2x} - 3e^x + 2 = 0$}
\begin{itemize}
    \item \textbf{Changement de variable :} on pose $X = e^x$ (donc $X > 0$). Comme $e^{2x}=(e^x)^2=X^2$, l'équation devient $X^2 - 3X + 2 = 0$.
    \item \textbf{Résolution :} $\Delta = 9-8=1$, $X=1$ ou $X=2$. Les deux valeurs sont strictement positives : valides.
    \item \textbf{Retour à $x$ :} $e^x = 1 \iff x = \ln 1 = 0$ ; $e^x = 2 \iff x = \ln 2$.
\end{itemize}
$\mathcal{S} = \{ 0\,;\ \ln 2 \}$.

\textbf{Équation 3 : $\ln(x+1) = 2 - \ln(x-1)$}
\begin{itemize}
    \item \textbf{C.E. :} $\begin{cases} x+1 > 0 \\ x-1 > 0 \end{cases} \iff x > 1$. Donc $\mathcal{D} = ]1\,;\, +\infty[$.
    \item \textbf{Regroupement :} $\ln(x+1) + \ln(x-1) = 2 \iff \ln\big[(x+1)(x-1)\big] = 2$.
    \item \textbf{Passage à l'exponentielle :} $(x+1)(x-1) = e^2 \iff x^2 - 1 = e^2 \iff x^2 = e^2 + 1$.
    \item \textbf{Solutions algébriques :} $x = \sqrt{e^2+1}$ ou $x = -\sqrt{e^2+1}$.
    \item \textbf{Vérification des C.E. :} $\sqrt{e^2+1} > \sqrt{e^2} = e > 1$ (conservé) ; $-\sqrt{e^2+1} < 0 < 1$ (rejeté).
\end{itemize}
$\mathcal{S} = \left\{ \sqrt{e^2+1} \right\}$.
\end{exoresolu}

\begin{methode}[Résolution d'inéquations faisant intervenir $\ln$ et $\exp$]
\textbf{Principe :} $\ln$ et $\exp$ sont \textbf{strictement croissantes} sur leurs domaines : elles \emph{conservent l'ordre}.
\begin{enumerate}
    \item \textbf{C.E. :} déterminer $\mathcal{D}$ (arguments des $\ln$ strictement positifs).
    \item \textbf{Isoler le $\ln$ ou l'$\exp$ :} se ramener à $\ln(u(x)) \geqslant \ln(v(x))$, ou $\exp(u(x)) \geqslant \exp(v(x))$, ou $\ln(u(x)) \geqslant k$, ou $\exp(u(x)) \geqslant k$.
    \item \textbf{Utiliser la stricte croissance :}
    \begin{itemize}
        \item $\ln A \geqslant \ln B \iff \begin{cases} A > 0 \\ B > 0 \\ A \geqslant B \end{cases}$ (idem avec $>$, $\leqslant$, $<$) ;
        \item $\ln A \geqslant k \iff \begin{cases} A > 0 \\ A \geqslant e^k \end{cases}$ ;
        \item $\exp A \geqslant \exp B \iff A \geqslant B$ (aucune condition de signe sur $A$ et $B$) ;
        \item $\exp A \geqslant k$ : si $k>0$, $\iff A \geqslant \ln k$ ; si $k\leqslant 0$, l'inéquation est \textbf{toujours vraie} car $\exp A > 0$.
    \end{itemize}
    \item \textbf{Résoudre le système d'inéquations algébriques obtenu.}
    \item \textbf{Intersecter avec $\mathcal{D}$} (souvent déjà pris en compte par les conditions $A>0$, $B>0$).
    \item \textbf{Conclusion :} écrire $\mathcal{S}$ sous forme d'intervalle ou de réunion d'intervalles.
\end{enumerate}
\end{methode}

\begin{exoresolu}[Inéquations $\ln$ et $\exp$]
Résoudre dans $\mathbb{R}$ :
\begin{enumerate}
    \item $\ln(x^2 - 3) \geqslant 0$
    \item $e^{2x-1} < e^{x+2}$
    \item $\ln(2x+1) < \ln(x+3)$
\end{enumerate}
\tcblower
\textbf{Solution détaillée}

\textbf{Inéquation 1 : $\ln(x^2 - 3) \geqslant 0$}
\begin{itemize}
    \item \textbf{C.E. :} $x^2 - 3 > 0 \iff x^2 > 3 \iff x < -\sqrt{3}$ ou $x > \sqrt{3}$. Donc $\mathcal{D} = ]-\infty\,;\,-\sqrt{3}[ \cup ]\sqrt{3}\,;\,+\infty[$.
    \item \textbf{Monotonie :} $\ln$ est croissante et $0=\ln 1$, donc $\ln(x^2-3) \geqslant \ln 1 \iff \begin{cases} x^2-3 > 0 \\ x^2-3 \geqslant 1 \end{cases}$.
    \item \textbf{Système :} $\begin{cases} x^2 > 3 \\ x^2 \geqslant 4 \end{cases}$. Or $x^2\geqslant 4 \Rightarrow x^2 > 3$, donc le système équivaut à $x^2 \geqslant 4$.
    \item \textbf{Résolution :} $x^2 \geqslant 4 \iff x \leqslant -2$ ou $x \geqslant 2$.
\end{itemize}
$\mathcal{S} = ]-\infty\,;\,-2] \cup [2\,;\,+\infty[$.

\textbf{Inéquation 2 : $e^{2x-1} < e^{x+2}$}
\begin{itemize}
    \item $\exp$ est strictement croissante sur $\mathbb{R}$ ; aucune C.E. n'est requise.
    \item $e^{2x-1} < e^{x+2} \iff 2x-1 < x+2 \iff x < 3$.
\end{itemize}
$\mathcal{S} = ]-\infty\,;\, 3[$.

\textbf{Inéquation 3 : $\ln(2x+1) < \ln(x+3)$}
\begin{itemize}
    \item \textbf{C.E. :} $\begin{cases} 2x+1 > 0 \\ x+3 > 0 \end{cases} \iff \begin{cases} x > -1/2 \\ x > -3 \end{cases} \iff x > -\dfrac12$. Donc $\mathcal{D} = \left]-\dfrac12\,;\, +\infty\right[$.
    \item \textbf{Monotonie :} $\ln$ croissante : $\ln(2x+1) < \ln(x+3) \iff \begin{cases} 2x+1 > 0 \\ x+3 > 0 \\ 2x+1 < x+3 \end{cases}$.
    \item \textbf{Système :} $\begin{cases} x > -1/2 \\ x < 2 \end{cases} \iff -\dfrac12 < x < 2$.
\end{itemize}
$\mathcal{S} = \left]-\dfrac12\,;\, 2\right[$.
\end{exoresolu}

\begin{methode}[Recherche de primitives de fonctions homographiques et usuelles]
\textbf{Formules à connaître par cœur} (sur un intervalle où les expressions sont définies) :
\begin{itemize}
    \item Une primitive de $\dfrac{u'}{u}$ est $\ln|u|$ \quad ($u$ dérivable, $u$ ne s'annule pas) ;
    \item Une primitive de $u'\, e^{u}$ est $e^{u}$ ;
    \item Une primitive de $\dfrac{1}{x}$ est $\ln|x|$ ;
    \item Une primitive de $e^x$ est $e^x$ ; une primitive de la constante $k$ est $kx$.
\end{itemize}
\textbf{Démarche pour une fonction homographique $f(x)=\dfrac{ax+b}{cx+d}$ ($c\neq 0$) :}
\begin{enumerate}
    \item Effectuer la \textbf{division euclidienne} du numérateur par le dénominateur : $ax+b = \dfrac{a}{c}(cx+d) + \left(b-\dfrac{ad}{c}\right)$.
    \item Écrire $f(x) = \dfrac{a}{c} + \dfrac{bc-ad}{c(cx+d)}$.
    \item Primitiver terme à terme : une primitive de $\dfrac{a}{c}$ est $\dfrac{a}{c}x$ ; une primitive de $\dfrac{bc-ad}{c(cx+d)}$ est $\dfrac{bc-ad}{c^2}\ln|cx+d|$ (car la dérivée de $cx+d$ est $c$, d'où le facteur $\dfrac{1}{c}$ supplémentaire).
    \item Ajouter la constante $C$ et préciser l'intervalle de définition.
\end{enumerate}
\textbf{Astuce de vérification :} dériver la primitive trouvée et contrôler qu'on retombe sur $f$.
\end{methode}

\begin{exoresolu}[Primitives : cas homographique et formes composées]
Déterminer une primitive sur l'intervalle indiqué :
\begin{enumerate}
    \item $f(x) = \dfrac{3x+2}{x-1}$ sur $]1\,;\,+\infty[$.
    \item $g(x) = (2x-1)e^{x^2-x}$ sur $\mathbb{R}$.
    \item $h(x) = \dfrac{2x}{x^2+4}$ sur $\mathbb{R}$.
\end{enumerate}
\tcblower
\textbf{Solution détaillée}

\textbf{1. $f(x) = \dfrac{3x+2}{x-1}$ sur $]1\,;\,+\infty[$}
\begin{itemize}
    \item \textbf{Division euclidienne :} $3x+2 = 3(x-1) + 5$, donc $f(x) = 3 + \dfrac{5}{x-1}$.
    \item \textbf{Primitive :} $F(x) = 3x + 5\ln|x-1| + C$.
    \item Sur $]1\,;\,+\infty[$, $x-1 > 0$, donc $|x-1| = x-1$ : $F(x) = 3x + 5\ln(x-1) + C$.
    \item \emph{Vérification :} $F'(x) = 3 + \dfrac{5}{x-1} = \dfrac{3(x-1)+5}{x-1} = \dfrac{3x+2}{x-1} = f(x)$. C'est correct.
\end{itemize}

\textbf{2. $g(x) = (2x-1)e^{x^2-x}$ sur $\mathbb{R}$}
\begin{itemize}
    \item On reconnaît la forme $u'\,e^{u}$ avec $u(x) = x^2 - x$, donc $u'(x) = 2x - 1$ : c'est exactement le facteur devant l'exponentielle.
    \item \textbf{Primitive :} $G(x) = e^{x^2-x} + C$.
    \item \emph{Vérification :} $G'(x) = (2x-1)e^{x^2-x} = g(x)$. C'est correct.
\end{itemize}

\textbf{3. $h(x) = \dfrac{2x}{x^2+4}$ sur $\mathbb{R}$}
\begin{itemize}
    \item On reconnaît la forme $\dfrac{u'}{u}$ avec $u(x) = x^2+4$, donc $u'(x) = 2x$. De plus $u(x) = x^2+4 > 0$ pour tout réel $x$.
    \item \textbf{Primitive :} $H(x) = \ln(x^2+4) + C$ (la valeur absolue est inutile car $x^2+4>0$).
    \item \emph{Vérification :} $H'(x) = \dfrac{2x}{x^2+4} = h(x)$. C'est correct.
\end{itemize}
\end{exoresolu}

\begin{methode}[Étude complète des fonctions $x \mapsto \ln(ax+b)$ et $x \mapsto \exp(ax+b)$]
\textbf{Plan d'étude attendu au Probatoire / Baccalauréat technique :}
\begin{enumerate}
    \item \textbf{Domaine $D_f$ :}
    \begin{itemize}
        \item $f(x)=\ln(ax+b)$ : $ax+b > 0 \iff x > -\dfrac{b}{a}$ si $a>0$, ou $x < -\dfrac{b}{a}$ si $a<0$ ;
        \item $f(x)=\exp(ax+b)$ : $D_f = \mathbb{R}$.
    \end{itemize}
    \item \textbf{Limites aux bornes de $D_f$ :}
    \begin{itemize}
        \item $\ln(ax+b)$ : limite $-\infty$ à la borne où $ax+b\to 0^+$ (asymptote verticale) ; limite infinie de signe opposé à l'autre borne ;
        \item $\exp(ax+b)$ : limites $0$ et $+\infty$ en $\pm\infty$ selon le signe de $a$ ; la limite nulle donne l'asymptote horizontale $y=0$ (axe des abscisses).
    \end{itemize}
    \item \textbf{Dérivée :} $f'(x) = \dfrac{a}{ax+b}$ pour $\ln(ax+b)$ ; $f'(x) = a\,e^{ax+b}$ pour $\exp(ax+b)$.
    \item \textbf{Signe de $f'$ et variations :}
    \begin{itemize}
        \item $\ln(ax+b)$ : sur $D_f$, $ax+b>0$, donc $f'$ a le signe de $a$ ;
        \item $\exp(ax+b)$ : $e^{ax+b}>0$ toujours, donc $f'$ a le signe de $a$.
    \end{itemize}
    Dans les deux cas : croissante si $a>0$, décroissante si $a<0$.
    \item \textbf{Tableau de variations} (une seule branche).
    \item \textbf{Points particuliers :} intersections avec les axes ($f(0)$, $f(x)=0$).
    \item \textbf{Tracé de la courbe.}
\end{enumerate}
\end{methode}

\begin{exoresolu}[Étude de $f(x)=\ln(3-2x)$ et $g(x)=e^{1-2x}$]
\textbf{Partie A.} Soit $f(x) = \ln(3-2x)$.
\begin{enumerate}
    \item Déterminer $D_f$, les limites aux bornes, la dérivée et le tableau de variations.
    \item Résoudre $f(x) = 0$ puis $f(x) > 0$.
    \item Tracer $\mathcal{C}_f$.
\end{enumerate}
\textbf{Partie B.} Soit $g(x) = e^{1-2x}$.
\begin{enumerate}
    \item Déterminer les limites, la dérivée et le tableau de variations.
    \item Résoudre $g(x) = 1$ puis $g(x) \geqslant 2$.
    \item Tracer $\mathcal{C}_g$ dans le même repère que $\mathcal{C}_f$ (unité : 2 cm).
\end{enumerate}
\tcblower
\textbf{Solution détaillée}

\textbf{Partie A : $f(x) = \ln(3-2x)$}
\begin{enumerate}
    \item \textbf{Domaine :} $3-2x > 0 \iff 2x < 3 \iff x < 1,5$. Donc $D_f = ]-\infty\,;\, 1,5[$.

    \textbf{Limites :} $\lim\limits_{x\to -\infty} (3-2x) = +\infty$, donc $\lim\limits_{x\to -\infty} f(x) = +\infty$.

    $\lim\limits_{x\to 1,5^-} (3-2x) = 0^+$, donc $\lim\limits_{x\to 1,5^-} f(x) = -\infty$ : la droite $x = 1,5$ est \textbf{asymptote verticale}.

    \textbf{Dérivée :} $f'(x) = \dfrac{-2}{3-2x}$. Sur $D_f$, $3-2x > 0$, donc $f'(x) < 0$ : $f$ est \textbf{strictement décroissante} sur $]-\infty\,;\,1,5[$.
    \begin{center}
    \begin{tabular}{|c|ccc|}\hline
    $x$ & $-\infty$ & & $1,5$ \\ \hline
    $f'(x)$ & & $-$ & \\ \hline
    $f$ & $+\infty$ & $\searrow$ & $-\infty$ \\ \hline
    \end{tabular}
    \end{center}
    \item \textbf{Équation $f(x)=0$ :} $\ln(3-2x)=0 \iff 3-2x=e^0=1 \iff x=1$. Point $A(1\,;\,0)$.

    \textbf{Inéquation $f(x)>0$ :} $\ln(3-2x) > 0 \iff 3-2x > 1 \iff x < 1$. Donc $\mathcal{S} = ]-\infty\,;\, 1[$.

    \textbf{Ordonnée à l'origine :} $f(0)=\ln 3 \approx 1,1$ : point $B(0\,;\,1,1)$.
\end{enumerate}

\textbf{Partie B : $g(x) = e^{1-2x}$}
\begin{enumerate}
    \item \textbf{Domaine :} $D_g = \mathbb{R}$.

    \textbf{Limites :} $\lim\limits_{x\to -\infty} (1-2x) = +\infty$, donc $\lim\limits_{x\to -\infty} g(x) = +\infty$.

    $\lim\limits_{x\to +\infty} (1-2x) = -\infty$, donc $\lim\limits_{x\to +\infty} g(x) = 0$ : l'axe des abscisses $y=0$ est \textbf{asymptote horizontale} en $+\infty$.

    \textbf{Dérivée :} $g'(x) = -2\, e^{1-2x}$. Comme $e^{1-2x} > 0$, on a $g'(x) < 0$ sur $\mathbb{R}$ : $g$ est \textbf{strictement décroissante} sur $\mathbb{R}$.
    \begin{center}
    \begin{tabular}{|c|ccc|}\hline
    $x$ & $-\infty$ & & $+\infty$ \\ \hline
    $g'(x)$ & & $-$ & \\ \hline
    $g$ & $+\infty$ & $\searrow$ & $0$ \\ \hline
    \end{tabular}
    \end{center}
    \item \textbf{Équation $g(x)=1$ :} $e^{1-2x}=1=e^0 \iff 1-2x=0 \iff x=0,5$. Point $C(0,5\,;\,1)$.

    \textbf{Inéquation $g(x)\geqslant 2$ :} $e^{1-2x} \geqslant 2 \iff 1-2x \geqslant \ln 2 \iff -2x \geqslant \ln 2 - 1 \iff x \leqslant \dfrac{1-\ln 2}{2}$ (on divise par $-2<0$, le sens change). $\dfrac{1-\ln 2}{2}\approx 0,15$. Donc $\mathcal{S} = \left]-\infty\,;\, \dfrac{1-\ln 2}{2}\right]$.

    \textbf{Ordonnée à l'origine :} $g(0)=e \approx 2,7$ : point $D(0\,;\,2,7)$.
\end{enumerate}

\textbf{Tracé commun :} repère orthonormé, unité 2 cm. Tracer en pointillés l'asymptote verticale $x=1,5$. $\mathcal{C}_f$ : branche logarithmique décroissante passant par $A(1;0)$ et $B(0\,;\,1,1)$, s'approchant de l'asymptote $x=1,5$. $\mathcal{C}_g$ : branche exponentielle décroissante passant par $C(0,5\,;\,1)$ et $D(0\,;\,2,7)$, s'approchant de l'axe $Ox$.

\begin{popfigure}
\begin{tikzpicture}[scale=0.9]
\draw[popline,->] (-3.5,0) -- (3,0) node[below left]{$x$};
\draw[popline,->] (0,-3) -- (0,4) node[below left]{$y$};
\draw[popmuted,dashed] (1.5,-3) -- (1.5,4) node[above]{$x=1,5$};
\draw[popblue,thick,domain=-3.2:1.35,samples=150] plot (\x,{ln(3-2*\x)});
\draw[poporange,thick,domain=-1.2:2.8,samples=150] plot (\x,{exp(1-2*\x)});
\fill[popdark] (1,0) circle (1.8pt) node[below right]{$A$};
\fill[popdark] (0,1.099) circle (1.8pt) node[above left]{$B$};
\fill[popdark] (0.5,1) circle (1.8pt) node[below right]{$C$};
\fill[popdark] (0,2.718) circle (1.8pt) node[above left]{$D$};
\node[popblue] at (-2.2,2.6) {$\mathcal{C}_f$};
\node[poporange] at (2.2,0.35) {$\mathcal{C}_g$};
\end{tikzpicture}
\legende{Courbes de $f(x)=\ln(3-2x)$ (bleu) et $g(x)=e^{1-2x}$ (orange).}
\end{popfigure}
\end{exoresolu}

\begin{methode}[Modélisation concrète : croissance et décroissance exponentielles]
\textbf{Schémas classiques à reconnaître :}
\begin{enumerate}
    \item \textbf{Croissance exponentielle} (population, capital à intérêts composés, culture de bactéries) : $f(t) = f_0\, e^{kt}$ avec $k>0$.
    \begin{itemize}
        \item Temps de doublement : $T_2 = \dfrac{\ln 2}{k}$ ;
        \item $f(t)=M \iff t = \dfrac{1}{k}\ln\dfrac{M}{f_0}$.
    \end{itemize}
    \item \textbf{Décroissance exponentielle} (radioactivité, décharge d'un condensateur, refroidissement) : $f(t) = f_0\, e^{-kt}$ avec $k>0$.
    \begin{itemize}
        \item Demi-vie (ou période) : $T_{1/2} = \dfrac{\ln 2}{k}$.
    \end{itemize}
    \item \textbf{Démarche type « problème concret » :}
    \begin{enumerate}
        \item Identifier la fonction modèle $f(t)$ et la variable $t$ (avec son unité).
        \item Extraire les données : valeur initiale $f_0$, puis un point $(t_1, f_1)$ pour déterminer $k$ : $k = \dfrac{1}{t_1}\ln\dfrac{f_1}{f_0}$.
        \item Répondre à la question posée (valeur à un instant donné, temps pour atteindre un seuil, demi-vie ou temps de doublement).
        \item Rédiger une \textbf{phrase de conclusion} avec l'unité.
    \end{enumerate}
\end{enumerate}
\end{methode}

\begin{exoresolu}[Application : désintégration radioactive]
Un échantillon d'iode 131 a une masse initiale $m_0 = 100$ mg. Sa masse suit la loi $m(t) = 100\, e^{-kt}$, où $t$ est en jours. On sait qu'après 8 jours, il reste 50 mg.
\begin{enumerate}
    \item Déterminer la constante $k$ et la demi-vie $T_{1/2}$.
    \item Au bout de combien de jours la masse devient-elle inférieure à 10 mg ?
\end{enumerate}
\tcblower
\textbf{Solution détaillée}
\begin{enumerate}
    \item \textbf{Calcul de $k$ :} $m(8) = 50 \iff 100\, e^{-8k} = 50 \iff e^{-8k} = \dfrac12$.
    En passant au logarithme : $-8k = \ln\dfrac12 = -\ln 2$, donc $k = \dfrac{\ln 2}{8} \approx 0,0866\ \text{jour}^{-1}$.

    \textbf{Demi-vie :} $T_{1/2} = \dfrac{\ln 2}{k} = \dfrac{\ln 2}{\ln 2 / 8} = 8$ jours. (Cohérent avec l'énoncé : la masse passe de 100 mg à 50 mg en 8 jours.)
    \item \textbf{Masse inférieure à 10 mg :} $100\, e^{-kt} < 10 \iff e^{-kt} < 0,1 \iff -kt < \ln 0,1 = -\ln 10$.
    En divisant par $-k < 0$ (le sens de l'inégalité change) : $t > \dfrac{\ln 10}{k} = \dfrac{8\ln 10}{\ln 2} \approx \dfrac{8\times 2,303}{0,693} \approx 26,6$ jours.

    \textbf{Conclusion :} la masse devient inférieure à 10 mg à partir du 27\textsuperscript{e} jour (premier entier supérieur à 26,6).
\end{enumerate}
\end{exoresolu}

\begin{exoresolu}[Application : évolution d'une température]
Dans un atelier, une pièce sortant du four à $T_0 = 200$ degrés refroidit selon la loi $T(t) = 20 + 180\, e^{-0,05t}$, où $T$ est la température en degrés et $t$ le temps en minutes.
\begin{enumerate}
    \item Vers quelle température la pièce tend-elle quand $t$ devient grand ? Interpréter.
    \item Au bout de combien de minutes la température passe-t-elle sous 50 degrés ?
\end{enumerate}
\tcblower
\textbf{Solution détaillée}
\begin{enumerate}
    \item $\lim\limits_{t\to +\infty} e^{-0,05t} = 0$, donc $\lim\limits_{t\to +\infty} T(t) = 20$. La pièce se refroidit vers la température ambiante de l'atelier, 20 degrés (asymptote horizontale $y=20$).
    \item $T(t) < 50 \iff 20 + 180\, e^{-0,05t} < 50 \iff e^{-0,05t} < \dfrac{30}{180} = \dfrac16$.
    En passant au logarithme : $-0,05t < \ln\dfrac16 = -\ln 6$, donc $t > \dfrac{\ln 6}{0,05} = 20\ln 6 \approx 20\times 1,792 \approx 35,8$ minutes.

    \textbf{Conclusion :} la température de la pièce passe sous 50 degrés au bout d'environ 36 minutes.
\end{enumerate}
\end{exoresolu}

\begin{attention}[Les 5 erreurs qui coûtent des points au Probatoire]
\begin{enumerate}
    \item \textbf{Oublier les conditions d'existence pour $\ln$ :} ne pas écrire $\mathcal{D}$, ou ne pas vérifier que les solutions trouvées appartiennent à $\mathcal{D}$. Exemple : $\ln(x-1)=0 \iff x-1=1 \iff x=2$ (et non $x=1$, qui n'est même pas dans $\mathcal{D}$).
    \item \textbf{Confondre $\ln(x^2)$ et $2\ln x$ :} l'égalité $\ln(x^2) = 2\ln x$ n'est valable que pour $x>0$. Pour $x\neq 0$, on a $\ln(x^2)=2\ln|x|$.
    \item \textbf{Perdre les conditions dans les inéquations avec $\ln$ :} $\ln A < \ln B \iff \begin{cases} A>0 \\ B>0 \\ A<B \end{cases}$. Les conditions $A>0$ et $B>0$ doivent figurer dans le système final.
    \item \textbf{Oublier la dérivée intérieure dans les formes composées :} une primitive de $\dfrac{1}{2x+1}$ est $\dfrac12\ln|2x+1|$ (et non $\ln|2x+1|$) ; de même $\big(e^{3x+1}\big)' = 3e^{3x+1}$.
    \item \textbf{Tableau de variations mal construit :}
    \begin{itemize}
        \item valeur interdite non séparée par une double barre ;
        \item flèches ($\nearrow$, $\searrow$) placées sur la ligne de $f'$ au lieu de celle de $f$ ;
        \item limites non calculées ou asymptotes non déduites des limites ;
        \item division par un nombre négatif dans une inéquation sans changer le sens de l'inégalité.
    \end{itemize}
\end{enumerate}
\textbf{Réflexe de relecture :} « Ai-je écrit $\mathcal{D}$ ? Ai-je vérifié mes solutions dans $\mathcal{D}$ ? Mes flèches sont-elles sur la bonne ligne ? Ai-je noté la constante $C$ pour les primitives ? »
\end{attention}

\newpage

\coursec{Exercices}

\courssub{Je vérifie mes connaissances}

\begin{exercice}[QCM : Domaine et limites d'une fonction homographique]
Parmi les propositions suivantes, une seule est exacte pour la fonction $f$ définie sur $\mathbb{R}\setminus\{-2\}$ par $f(x)=\dfrac{3x+1}{x+2}$.
\begin{enumerate}
    \item $\displaystyle\lim_{x\to -2^+} f(x) = -\infty$ et $\displaystyle\lim_{x\to +\infty} f(x) = 3$.
    \item $\displaystyle\lim_{x\to -2^-} f(x) = +\infty$ et $\displaystyle\lim_{x\to -\infty} f(x) = 0$.
    \item $\displaystyle\lim_{x\to -2^+} f(x) = +\infty$ et $\displaystyle\lim_{x\to +\infty} f(x) = 3$.
    \item $\displaystyle\lim_{x\to -2^-} f(x) = -\infty$ et $\displaystyle\lim_{x\to +\infty} f(x) = 1$.
\end{enumerate}
\end{exercice}

\begin{exercice}[Vrai ou Faux : Propriétés du logarithme népérien]
Dire si les affirmations suivantes sont vraies ou fausses. Justifier brièvement chaque réponse.
\begin{enumerate}
    \item $\ln(1) = e$.
    \item $\forall a>0,\ \forall b>0,\ \ln\left(\dfrac{a}{b}\right) = \ln a - \ln b$.
    \item $\ln(e^2) = 2$.
    \item La fonction $x \mapsto \ln(x)$ est strictement décroissante sur $]0;+\infty[$.
\end{enumerate}
\end{exercice}

\begin{exercice}[Calculs directs : Exponentielle et dérivées]
Effectuer les calculs suivants sans calculatrice.
\begin{enumerate}
    \item Simplifier $e^{\ln 5} \times e^{2\ln 2}$.
    \item Calculer la dérivée de la fonction $g(x) = e^{3x-1}$ sur $\mathbb{R}$.
    \item Calculer $\displaystyle\int_1^2 \dfrac{4}{x}\,dx$ (donner une valeur exacte).
    \item Résoudre l'équation $e^{2x} = e^5$.
\end{enumerate}
\end{exercice}

\begin{exercice}[Définitions et primitives]
\begin{enumerate}
    \item Donner l'ensemble de définition de la fonction $h(x) = \ln(2x-6)$.
    \item Écrire sous la forme $a\ln b$ (avec $a,b$ entiers) l'expression $\ln 8 - \ln 2$.
    \item Déterminer une primitive sur $]1;+\infty[$ de la fonction $k(x) = \dfrac{2x+1}{x-1}$.
    \item Compléter : La fonction exponentielle népérienne est la fonction réciproque de la fonction \ldots
\end{enumerate}
\end{exercice}

\courssub{Je m'entraîne}

\begin{exercice}[Étude complète d'une fonction homographique]
On considère la fonction $f$ définie sur $\mathbb{R}\setminus\{-3\}$ par $f(x) = \dfrac{2x-5}{x+3}$.
\begin{enumerate}
    \item Déterminer les limites de $f$ aux bornes de son ensemble de définition. En déduire les asymptotes (verticale et horizontale) à la courbe $\mathcal{C}_f$.
    \item Calculer $f'(x)$ et dresser le tableau de variations de $f$.
    \item Montrer que la courbe $\mathcal{C}_f$ admet un centre de symétrie. En donner les coordonnées.
    \item Tracer $\mathcal{C}_f$ dans un repère orthonormé (unité : 2 cm).
\end{enumerate}
\end{exercice}

\begin{exercice}[Équations et inéquations logarithmiques et exponentielles]
Résoudre dans $\mathbb{R}$ les équations et inéquations suivantes.
\begin{enumerate}
    \item $\ln(x+2) + \ln(x-2) = \ln 5$.
    \item $e^{2x} - 7e^x + 12 = 0$.
    \item $\ln(3x-1) \ge 0$.
    \item $(\ln x - 1)(e^x - 1) \le 0$ sur $]0;+\infty[$.
\end{enumerate}
\end{exercice}

\begin{exercice}[Primitives de fonctions homographiques et usuelles]
\begin{enumerate}
    \item Décomposer la fonction $f(x) = \dfrac{4x+3}{x-1}$ (définie sur $\mathbb{R}\setminus\{1\}$) en une somme de termes simples.
    \item Déterminer l'ensemble des primitives de $f$ sur $]1;+\infty[$ et sur $]-\infty;1[$.
    \item Vérifier par dérivation que l'expression trouvée est correcte.
    \item Calculer $\displaystyle\int_2^3 \dfrac{4x+3}{x-1}\,dx$.
\end{enumerate}
\end{exercice}

\begin{exercice}[Étude des fonctions $x \mapsto \ln(ax+b)$ et $x \mapsto e^{ax+b}$]
\begin{enumerate}
    \item Étudier la fonction $g$ définie par $g(x) = \ln(2x+4)$ : domaine de définition, limites, dérivée, tableau de variations, équation de la tangente au point d'abscisse $0$. Tracer sa courbe $\mathcal{C}_g$.
    \item Étudier la fonction $h$ définie par $h(x) = e^{-x+3}$ : domaine de définition, limites, dérivée, tableau de variations. Tracer sa courbe $\mathcal{C}_h$ dans le même repère que $\mathcal{C}_g$.
    \item Les courbes $\mathcal{C}_g$ et $\mathcal{C}_h$ sont-elles symétriques par rapport à la droite $y=x$ ? Justifier.
\end{enumerate}
\end{exercice}

\courssub{Je me prépare au Probatoire}

\begin{exercice}[Problème de synthèse : Dépréciation d'une machine]
Une entreprise achète une machine-outil pour un montant de $5\,000\,000$ FCFA. Sa valeur $V(t)$ (en FCFA) au bout de $t$ années suit la loi de dépréciation exponentielle :
\[ V(t) = 5\,000\,000 \times e^{-0,2t} \quad (t \ge 0). \]
\begin{enumerate}
    \item Calculer la valeur de la machine après 3 ans (arrondir au FCFA le plus proche).
    \item Au bout de combien d'années la valeur de la machine sera-t-elle inférieure à $1\,000\,000$ FCFA ? (Donner l'entier $t$ minimal).
    \item On définit le taux de dépréciation annuel moyen entre l'année $t$ et $t+1$ par $\tau(t) = \dfrac{V(t)-V(t+1)}{V(t)}$.
    \begin{enumerate}
        \item Exprimer $\tau(t)$ en fonction de $t$ et montrer qu'il est constant.
        \item Interpréter ce résultat dans le contexte de l'amortissement exponentiel.
    \end{enumerate}
    \item L'entreprise souhaite revendre la machine lorsque sa valeur aura baissé de $60\%$ par rapport à son prix d'achat. Déterminer la date de revente (arrondir au dixième d'année).
\end{enumerate}
\end{exercice}

\begin{exercice}[Coût moyen et seuil de rentabilité en menuiserie]
Un atelier de menuiserie produit $x$ meubles par mois ($x > 0$). Le coût total de production mensuel (en milliers de FCFA) est modélisé par $CT(x) = 80x + 45\,000$.
Le coût moyen unitaire est $C(x) = \dfrac{CT(x)}{x}$.
\begin{enumerate}
    \item Exprimer $C(x)$ sous la forme d'une fonction homographique. Préciser son ensemble de définition.
    \item Étudier les variations de $C$ sur $]0;+\infty[$. Construire son tableau de variations.
    \item Déterminer les limites de $C$ en $0^+$ et en $+\infty$. Interpréter économiquement ces deux limites.
    \item L'atelier vend chaque meuble $150\,000$ FCFA. Déterminer le nombre minimal de meubles à produire par mois pour que le coût moyen unitaire soit inférieur ou égal au prix de vente (seuil de rentabilité).
    \item Tracer la courbe de $C$ et la droite $y=150$ dans un repère orthonormé (unité : 1 cm pour 10 meubles en abscisses, 1 cm pour 20 milliers FCFA en ordonnées). Représenter graphiquement la zone de rentabilité.
\end{enumerate}
\end{exercice}

\begin{exercice}[Étude de fonction et résolution graphique -- Type Probatoire]
Soit la fonction $f$ définie sur $]-\infty; 2[ \cup ]2; +\infty[$ par $f(x) = \dfrac{x+1}{x-2}$.
\begin{enumerate}
    \item \textbf{Étude algébrique}
    \begin{enumerate}
        \item Calculer $f'(x)$. En déduire les variations de $f$ sur chaque intervalle de définition.
        \item Déterminer les limites de $f$ aux bornes de son ensemble de définition. Écrire les équations des asymptotes.
        \item Montrer que le point $\Omega(2\,;\,1)$ est le centre de symétrie de la courbe $\mathcal{C}_f$.
    \end{enumerate}
    \item \textbf{Résolution d'inéquation}
    Résoudre l'inéquation $f(x) \le 3$.
    \item \textbf{Partie graphique}
    On trace $\mathcal{C}_f$ dans un repère orthonormé.
    \begin{enumerate}
        \item Placer le point $\Omega$ et tracer les asymptotes.
        \item Tracer la droite $\mathcal{D}$ d'équation $y = 3x - 5$. Montrer que $\mathcal{D}$ est la tangente à $\mathcal{C}_f$ au point d'abscisse $3$.
        \item Déterminer graphiquement (et vérifier par le calcul) les abscisses des points d'intersection de $\mathcal{C}_f$ et de $\mathcal{D}$.
    \end{enumerate}
\end{enumerate}
\end{exercice}

\newpage

\coursec{Corrigés des exercices}

\courssub{Je vérifie mes connaissances}

\begin{corrige}[Exercice 1 — QCM : Domaine et limites d'une fonction homographique]
La bonne réponse est la proposition \textbf{1}.

\emph{Justification.} $f(x)=\dfrac{3x+1}{x+2}$ est une fonction homographique définie sur $\mathbb{R}\setminus\{-2\}$.

\textbf{Limite en $-2$ par valeurs supérieures :} quand $x\to -2^+$, le numérateur tend vers $3\times(-2)+1=-5<0$ et le dénominateur $x+2\to 0^+$. Par quotient :
\[ \lim_{x\to -2^+} f(x) = \frac{-5}{0^+} = -\infty. \]

\textbf{Limite en $+\infty$ :} pour une fraction rationnelle, la limite en l'infini est le quotient des termes de plus haut degré :
\[ \lim_{x\to +\infty} f(x) = \lim_{x\to +\infty} \frac{3x}{x} = 3. \]

\textbf{Élimination des autres propositions :} en $-2^-$, $x+2\to 0^-$, donc $\displaystyle\lim_{x\to -2^-}f(x)=\frac{-5}{0^-}=+\infty$ (et non $-\infty$) ; la limite en $-\infty$ vaut $3$ (et non $0$ ni $1$). Seule la proposition 1 est entièrement exacte.
\end{corrige}

\begin{corrige}[Exercice 2 — Vrai ou Faux : Propriétés du logarithme népérien]
\begin{enumerate}
    \item \textbf{Faux.} $\ln(1)=0$ (et non $e$). C'est $\ln(e)=1$.
    \item \textbf{Vrai.} C'est la propriété fondamentale du logarithme d'un quotient : pour tous $a>0$ et $b>0$, $\ln\left(\dfrac{a}{b}\right)=\ln a - \ln b$.
    \item \textbf{Vrai.} Pour tout réel $x$, $\ln(e^x)=x$, donc $\ln(e^2)=2$.
    \item \textbf{Faux.} La fonction $\ln$ est \emph{strictement croissante} sur $]0;+\infty[$, car sa dérivée $\dfrac{1}{x}$ est strictement positive sur cet intervalle.
\end{enumerate}
\end{corrige}

\begin{corrige}[Exercice 3 — Calculs directs : Exponentielle et dérivées]
\begin{enumerate}
    \item $e^{\ln 5}\times e^{2\ln 2} = 5 \times e^{\ln(2^2)} = 5\times 4 = \boxed{20}$.
    \item $g(x)=e^{3x-1}$ est de la forme $e^u$ avec $u(x)=3x-1$, $u'(x)=3$. Donc :
    \[ g'(x) = 3e^{3x-1}. \]
    \item Une primitive de $x\mapsto \dfrac{4}{x}$ sur $]0;+\infty[$ est $x\mapsto 4\ln x$. Donc :
    \[ \int_1^2 \frac{4}{x}\,dx = \Big[4\ln x\Big]_1^2 = 4\ln 2 - 4\ln 1 = 4\ln 2. \]
    \item Par injectivité de l'exponentielle : $e^{2x}=e^5 \iff 2x = 5 \iff x=\dfrac{5}{2}$. \quad $\mathcal{S}=\left\{\dfrac{5}{2}\right\}$.
\end{enumerate}
\end{corrige}

\begin{corrige}[Exercice 4 — Définitions et primitives]
\begin{enumerate}
    \item $h(x)=\ln(2x-6)$ est définie si et seulement si $2x-6>0 \iff x>3$. Donc $\mathcal{D}_h = ]3;+\infty[$.
    \item $\ln 8 - \ln 2 = \ln\left(\dfrac{8}{2}\right) = \ln 4 = \ln(2^2) = 2\ln 2$.
    \item On décompose : $k(x)=\dfrac{2x+1}{x-1}$. Or $2x+1 = 2(x-1)+3$, donc :
    \[ k(x) = 2 + \frac{3}{x-1}. \]
    Sur $]1;+\infty[$, $x-1>0$, une primitive est donc :
    \[ K(x) = 2x + 3\ln(x-1). \]
    \emph{Vérification :} $K'(x)=2+\dfrac{3}{x-1}=\dfrac{2(x-1)+3}{x-1}=\dfrac{2x+1}{x-1}=k(x)$. \checkmark
    \item La fonction exponentielle népérienne est la fonction réciproque de la fonction \textbf{logarithme népérien}.
\end{enumerate}
\end{corrige}

\courssub{Je m'entraîne}

\begin{corrige}[Exercice 5 — Étude complète d'une fonction homographique]
$f(x)=\dfrac{2x-5}{x+3}$, définie sur $\mathbb{R}\setminus\{-3\}$.

\textbf{1. Limites et asymptotes.}

En $-3$ : le numérateur tend vers $2\times(-3)-5=-11\neq 0$.
\begin{itemize}
    \item Quand $x\to -3^+$, $x+3\to 0^+$ : $\displaystyle\lim_{x\to -3^+}f(x)=\frac{-11}{0^+}=-\infty$.
    \item Quand $x\to -3^-$, $x+3\to 0^-$ : $\displaystyle\lim_{x\to -3^-}f(x)=\frac{-11}{0^-}=+\infty$.
\end{itemize}
La droite d'équation $\boxed{x=-3}$ est \textbf{asymptote verticale} à $\mathcal{C}_f$.

En l'infini : $\displaystyle\lim_{x\to\pm\infty}f(x)=\lim_{x\to\pm\infty}\frac{2x}{x}=2$.
La droite d'équation $\boxed{y=2}$ est \textbf{asymptote horizontale} à $\mathcal{C}_f$ en $-\infty$ et en $+\infty$.

\textbf{2. Dérivée et tableau de variations.}

$f$ est de la forme $\dfrac{u}{v}$ avec $u(x)=2x-5$, $v(x)=x+3$ :
\[ f'(x) = \frac{2(x+3)-(2x-5)\times 1}{(x+3)^2} = \frac{2x+6-2x+5}{(x+3)^2} = \frac{11}{(x+3)^2}. \]
Pour tout $x\neq -3$, $f'(x)>0$ : $f$ est strictement croissante sur $]-\infty;-3[$ et sur $]-3;+\infty[$.

\begin{center}
\begin{tabular}{|c|ccccc|}\hline
$x$ & $-\infty$ & & $-3$ & & $+\infty$ \\ \hline
$f'(x)$ & & $+$ & \textbf{||} & $+$ & \\ \hline
$f$ & $2$ & $\nearrow$ & \textbf{||} $\ -\infty$ / $+\infty$\ \textbf{||} & $\nearrow$ & $2$ \\ \hline
\end{tabular}
\end{center}

\textbf{3. Centre de symétrie.}

Montrons que $\Omega(-3\,;\,2)$ est centre de symétrie, c'est-à-dire que pour tout $h\neq 0$ :
$f(-3+h)+f(-3-h)=2\times 2 = 4$.
\begin{align*}
f(-3+h) &= \frac{2(-3+h)-5}{h} = \frac{2h-11}{h} = 2-\frac{11}{h},\\
f(-3-h) &= \frac{2(-3-h)-5}{-h} = \frac{-2h-11}{-h} = 2+\frac{11}{h}.
\end{align*}
Donc $f(-3+h)+f(-3-h) = 4$ : le point $\boxed{\Omega(-3\,;\,2)}$ est bien le centre de symétrie de $\mathcal{C}_f$.

\emph{Remarque :} le centre de symétrie d'une hyperbole homographique est toujours le point d'intersection de ses deux asymptotes.

\textbf{4. Tracé.} On place les asymptotes $x=-3$ et $y=2$, le centre $\Omega(-3;2)$, puis quelques points : $A(-2;-9)$, $B(-1;-\tfrac{7}{2})$, $C(0;-\tfrac{5}{3})$, $D(2;-\tfrac{1}{5})$, $E(\tfrac{5}{2};0)$ (intersection avec l'axe des abscisses), $F(-4;13)$, $G(-5;\tfrac{15}{2})$. On trace les deux branches d'hyperbole, symétriques par rapport à $\Omega$, en suivant les asymptotes.
\end{corrige}

\begin{corrige}[Exercice 6 — Équations et inéquations logarithmiques et exponentielles]
\textbf{1.} $\ln(x+2)+\ln(x-2)=\ln 5$.

\emph{Conditions :} $x+2>0$ et $x-2>0$, donc $x>2$.

Sur $]2;+\infty[$ : $\ln\big[(x+2)(x-2)\big]=\ln 5 \iff x^2-4=5 \iff x^2=9 \iff x=3$ ou $x=-3$.

Seul $x=3$ vérifie la condition $x>2$. \quad $\mathcal{S}=\{3\}$.

\textbf{2.} $e^{2x}-7e^x+12=0$. Posons $X=e^x$ ($X>0$) : $X^2-7X+12=0$.

$\Delta = 49-48=1$ ; $X_1=\dfrac{7-1}{2}=3$, $X_2=\dfrac{7+1}{2}=4$ (toutes deux $>0$, acceptables).

$e^x=3 \iff x=\ln 3$ ; $e^x=4 \iff x=\ln 4$. \quad $\mathcal{S}=\{\ln 3\,;\,\ln 4\}$.

\textbf{3.} $\ln(3x-1)\ge 0$.

\emph{Condition :} $3x-1>0 \iff x>\dfrac{1}{3}$.

$\ln(3x-1)\ge 0 = \ln 1 \iff 3x-1\ge 1$ (croissance de $\ln$) $\iff x\ge \dfrac{2}{3}$.

Comme $\dfrac{2}{3}>\dfrac{1}{3}$ : \quad $\mathcal{S}=\left[\dfrac{2}{3};+\infty\right[$.

\textbf{4.} $(\ln x - 1)(e^x-1)\le 0$ sur $]0;+\infty[$.

\emph{Signe de chaque facteur :}
\begin{itemize}
    \item $\ln x - 1 \ge 0 \iff x\ge e$ ; donc $\ln x - 1 < 0$ sur $]0;e[$, nul en $e$, positif sur $]e;+\infty[$.
    \item $e^x - 1 \ge 0 \iff x\ge 0$ ; sur $]0;+\infty[$, $e^x-1>0$ toujours.
\end{itemize}
Le produit a donc le signe de $\ln x - 1$ : il est $\le 0$ si et seulement si $x\le e$.

$\mathcal{S}=]0;e]$.
\end{corrige}

\begin{corrige}[Exercice 7 — Primitives de fonctions homographiques et usuelles]
\textbf{1. Décomposition.} On écrit le numérateur en fonction de $x-1$ :
$4x+3 = 4(x-1)+7$, donc :
\[ f(x) = \frac{4(x-1)+7}{x-1} = 4 + \frac{7}{x-1}. \]

\textbf{2. Primitives.} Une primitive de $4$ est $4x$ ; une primitive de $\dfrac{7}{x-1}$ est $7\ln|x-1|$.
\begin{itemize}
    \item Sur $]1;+\infty[$, $x-1>0$ : $F(x)=4x+7\ln(x-1)+C$, $C\in\mathbb{R}$.
    \item Sur $]-\infty;1[$, $x-1<0$ : $F(x)=4x+7\ln(1-x)+C$, $C\in\mathbb{R}$.
\end{itemize}

\textbf{3. Vérification par dérivation} (sur $]1;+\infty[$) :
\[ F'(x) = 4 + \frac{7}{x-1} = \frac{4(x-1)+7}{x-1} = \frac{4x+3}{x-1} = f(x). \checkmark \]

\textbf{4. Calcul de l'intégrale.} L'intervalle $[2;3]\subset\,]1;+\infty[$ :
\begin{align*}
\int_2^3 \frac{4x+3}{x-1}\,dx &= \Big[4x+7\ln(x-1)\Big]_2^3\\
&= \big(12+7\ln 2\big) - \big(8+7\ln 1\big) = 4+7\ln 2.
\end{align*}
Valeur exacte : $\boxed{4+7\ln 2}$ (soit environ $8{,}85$).
\end{corrige}

\begin{corrige}[Exercice 8 — Étude des fonctions $x\mapsto\ln(ax+b)$ et $x\mapsto e^{ax+b}$]
\textbf{1. Étude de $g(x)=\ln(2x+4)$.}

\emph{Domaine :} $2x+4>0 \iff x>-2$ ; $\mathcal{D}_g=]-2;+\infty[$.

\emph{Limites :} quand $x\to -2^+$, $2x+4\to 0^+$ donc $\displaystyle\lim_{x\to -2^+}g(x)=-\infty$ (asymptote verticale $x=-2$) ; quand $x\to+\infty$, $\displaystyle\lim_{x\to+\infty}g(x)=+\infty$.

\emph{Dérivée :} $g$ est de la forme $\ln u$ avec $u(x)=2x+4$ :
\[ g'(x)=\frac{2}{2x+4}=\frac{1}{x+2}>0 \ \text{sur } ]-2;+\infty[. \]
$g$ est strictement croissante sur $]-2;+\infty[$.

\begin{center}
\begin{tabular}{|c|ccc|}\hline
$x$ & $-2$ & & $+\infty$ \\ \hline
$g'(x)$ & \textbf{||} & $+$ & \\ \hline
$g$ & \textbf{||} $-\infty$ & $\nearrow$ & $+\infty$ \\ \hline
\end{tabular}
\end{center}

\emph{Tangente en $0$ :} $g(0)=\ln 4$, $g'(0)=\dfrac{1}{2}$. Équation :
\[ y = g'(0)(x-0)+g(0) = \frac{1}{2}x + \ln 4. \]

\emph{Tracé :} points remarquables : $g(-1)=\ln 2\approx 0{,}69$ ; $g(0)=\ln 4\approx 1{,}39$ ; $g(-\tfrac{3}{2})=\ln 1 = 0$ (intersection avec l'axe des abscisses) ; asymptote verticale $x=-2$.

\textbf{2. Étude de $h(x)=e^{-x+3}$.}

\emph{Domaine :} $\mathcal{D}_h=\mathbb{R}$ (l'exponentielle est définie partout).

\emph{Limites :} quand $x\to+\infty$, $-x+3\to-\infty$ donc $\displaystyle\lim_{x\to+\infty}h(x)=0$ (asymptote horizontale $y=0$) ; quand $x\to-\infty$, $-x+3\to+\infty$ donc $\displaystyle\lim_{x\to-\infty}h(x)=+\infty$.

\emph{Dérivée :} $h'(x)=-e^{-x+3}<0$ pour tout $x$ : $h$ est strictement décroissante sur $\mathbb{R}$.

\begin{center}
\begin{tabular}{|c|ccc|}\hline
$x$ & $-\infty$ & & $+\infty$ \\ \hline
$h'(x)$ & & $-$ & \\ \hline
$h$ & $+\infty$ & $\searrow$ & $0$ \\ \hline
\end{tabular}
\end{center}

Points : $h(0)=e^3\approx 20{,}09$ ; $h(3)=1$ ; $h(4)=e^{-1}\approx 0{,}37$.

\textbf{3. Symétrie par rapport à $y=x$ ?}

Deux courbes sont symétriques par rapport à la droite $y=x$ si et seulement si les fonctions sont réciproques l'une de l'autre. Cherchons la réciproque de $g$ :
\[ y=\ln(2x+4) \iff e^y = 2x+4 \iff x = \frac{e^y-4}{2}. \]
Donc $g^{-1}(x)=\dfrac{e^x-4}{2}$, qui n'est pas égale à $h(x)=e^{-x+3}$.

\textbf{Conclusion :} $\mathcal{C}_g$ et $\mathcal{C}_h$ ne sont \emph{pas} symétriques par rapport à la droite $y=x$.
\end{corrige}

\courssub{Je me prépare au Probatoire}

\begin{corrige}[Exercice 9 — Problème de synthèse : Dépréciation d'une machine]
\textbf{1.} $V(3) = 5\,000\,000 \times e^{-0,2\times 3} = 5\,000\,000 \times e^{-0,6}$.

$e^{-0,6}\approx 0{,}548812$, donc $V(3)\approx 5\,000\,000\times 0{,}548812 \approx \boxed{2\,744\,058\ \text{FCFA}}$.

\textbf{2.} On résout $V(t)<1\,000\,000$ :
\begin{align*}
5\,000\,000\, e^{-0,2t} < 1\,000\,000 &\iff e^{-0,2t} < \frac{1}{5} = 0{,}2\\
&\iff -0{,}2t < \ln(0{,}2) \quad (\ln \text{ est croissante})\\
&\iff t > \frac{-\ln(0{,}2)}{0{,}2} = \frac{\ln 5}{0{,}2} = 5\ln 5 \approx 8{,}05.
\end{align*}
L'entier minimal est $\boxed{t=9}$ : c'est au bout de \textbf{9 ans} que la valeur passe sous $1\,000\,000$ FCFA.

\emph{Vérification :} $V(8)\approx 5\,000\,000\,e^{-1,6}\approx 1\,009\,483>10^6$ ; $V(9)\approx 5\,000\,000\,e^{-1,8}\approx 826\,494<10^6$. \checkmark

\textbf{3. a)} $\tau(t)=\dfrac{V(t)-V(t+1)}{V(t)} = 1 - \dfrac{V(t+1)}{V(t)}$.

Or $\dfrac{V(t+1)}{V(t)}=\dfrac{5\,000\,000\,e^{-0,2(t+1)}}{5\,000\,000\,e^{-0,2t}}=e^{-0,2}$.

Donc $\tau(t)=1-e^{-0,2}\approx 0{,}1813$, \textbf{indépendant de $t$} : le taux est constant.

\textbf{3. b)} \emph{Interprétation :} chaque année, la machine perd environ $\boxed{18{,}13\,\%}$ de sa valeur \emph{restante} (et non de sa valeur d'achat). C'est la caractéristique de l'amortissement exponentiel (ou « dégressif à taux constant ») : la dépréciation absolue diminue chaque année, mais le taux reste le même.

\textbf{4.} Une baisse de $60\,\%$ signifie $V(t)=40\,\%\times 5\,000\,000 = 2\,000\,000$ :
\[ e^{-0,2t} = 0{,}4 \iff -0{,}2t = \ln(0{,}4) \iff t = \frac{-\ln(0{,}4)}{0{,}2} = 5\ln(2{,}5)\approx 4{,}58. \]
La machine sera revendue au bout d'environ $\boxed{4{,}6\ \text{ans}}$ (soit environ 4 ans et 7 mois).

\textbf{Barème indicatif (sur 5 pts) :} Q1 : 0,75 pt ; Q2 : 1,25 pt (inéquation + passage au $\ln$ + entier) ; Q3 : 1,5 pt (0,75 + 0,75) ; Q4 : 1 pt ; présentation : 0,5 pt.

\textbf{Points de vigilance :} citer la croissance de $\ln$ pour justifier le sens de l'inégalité ; ne pas confondre « baisse de 60 \% » avec « valeur de 60 \% » ; vérifier l'entier demandé par encadrement ($t=8$ ne convient pas).
\end{corrige}

\begin{corrige}[Exercice 10 — Coût moyen et seuil de rentabilité en menuiserie]
\textbf{1.} $C(x)=\dfrac{CT(x)}{x}=\dfrac{80x+45\,000}{x}=80+\dfrac{45\,000}{x}$ (en milliers de FCFA).

C'est une fonction homographique (somme d'une constante et d'une fonction inverse), définie pour $x>0$ : $\mathcal{D}_C=]0;+\infty[$.

\textbf{2. Variations.} $C'(x)=-\dfrac{45\,000}{x^2}<0$ pour tout $x>0$ : $C$ est \textbf{strictement décroissante} sur $]0;+\infty[$.

\begin{center}
\begin{tabular}{|c|ccc|}\hline
$x$ & $0$ & & $+\infty$ \\ \hline
$C'(x)$ & \textbf{||} & $-$ & \\ \hline
$C$ & \textbf{||} $+\infty$ & $\searrow$ & $80$ \\ \hline
\end{tabular}
\end{center}

\textbf{3. Limites et interprétation.}
\begin{itemize}
    \item $\displaystyle\lim_{x\to 0^+}C(x)=+\infty$ : si l'atelier produit très peu de meubles, les coûts fixes ($45\,000$ milliers FCFA) sont répartis sur très peu d'unités, et le coût moyen explose.
    \item $\displaystyle\lim_{x\to+\infty}C(x)=80$ : en produisant beaucoup, le coût moyen se rapproche du coût variable unitaire de $80\,000$ FCFA par meuble, sans jamais descendre en dessous (asymptote horizontale $y=80$).
\end{itemize}

\textbf{4. Seuil de rentabilité.} Le prix de vente est $150\,000$ FCFA $=150$ milliers FCFA. On résout :
\begin{align*}
C(x)\le 150 &\iff 80+\frac{45\,000}{x}\le 150 \iff \frac{45\,000}{x}\le 70\\
&\iff x \ge \frac{45\,000}{70} \approx 642{,}86 \quad (x>0).
\end{align*}
L'atelier doit produire au minimum $\boxed{643\ \text{meubles par mois}}$ pour être rentable.

\textbf{5. Graphique.} On trace la courbe de $C$ (décroissante, asymptotes $x=0$ et $y=80$) et la droite horizontale $y=150$. Elles se coupent au point d'abscisse $x=\dfrac{45\,000}{70}\approx 643$. La \textbf{zone de rentabilité} est la partie du graphique située à droite de ce point d'intersection, où la courbe de $C$ passe \emph{sous} la droite $y=150$ (on la hachure ou on la colorie).

\textbf{Barème indicatif (sur 5 pts) :} Q1 : 0,5 pt ; Q2 : 1 pt ; Q3 : 1 pt (limites + interprétation) ; Q4 : 1,5 pt ; Q5 : 1 pt.

\textbf{Points de vigilance :} travailler dans une unité cohérente (milliers de FCFA) ; justifier le changement de sens éventuel de l'inégalité (ici $x>0$, le sens est conservé) ; donner une réponse \emph{entière} et la formuler dans le contexte (nombre de meubles).
\end{corrige}

\begin{corrige}[Exercice 11 — Étude de fonction et résolution graphique (type Probatoire)]
$f(x)=\dfrac{x+1}{x-2}$ sur $]-\infty;2[\cup]2;+\infty[$.

\textbf{1. a) Dérivée et variations.} Avec $u(x)=x+1$, $v(x)=x-2$ :
\[ f'(x)=\frac{1\times(x-2)-(x+1)\times 1}{(x-2)^2}=\frac{x-2-x-1}{(x-2)^2}=\frac{-3}{(x-2)^2}. \]
Pour tout $x\neq 2$, $f'(x)<0$ : $f$ est strictement décroissante sur $]-\infty;2[$ et sur $]2;+\infty[$.

\textbf{1. b) Limites et asymptotes.}
\begin{itemize}
    \item En $2^-$ : numérateur $\to 3>0$, dénominateur $\to 0^-$ : $\displaystyle\lim_{x\to 2^-}f(x)=-\infty$.
    \item En $2^+$ : $\displaystyle\lim_{x\to 2^+}f(x)=+\infty$. \quad Asymptote verticale : $\boxed{x=2}$.
    \item En $\pm\infty$ : $\displaystyle\lim_{x\to\pm\infty}f(x)=\lim_{x\to\pm\infty}\frac{x}{x}=1$. \quad Asymptote horizontale : $\boxed{y=1}$.
\end{itemize}

\begin{center}
\begin{tabular}{|c|ccccc|}\hline
$x$ & $-\infty$ & & $2$ & & $+\infty$ \\ \hline
$f'(x)$ & & $-$ & \textbf{||} & $-$ & \\ \hline
$f$ & $1$ & $\searrow$ & \textbf{||} $+\infty$ / $-\infty$ \textbf{||} & $\searrow$ & $1$ \\ \hline
\end{tabular}
\end{center}

\textbf{1. c) Centre de symétrie $\Omega(2;1)$.} Pour tout $h\neq 0$ :
\begin{align*}
f(2+h)&=\frac{3+h}{h}=1+\frac{3}{h},\\
f(2-h)&=\frac{3-h}{-h}=1-\frac{3}{h}.
\end{align*}
Donc $f(2+h)+f(2-h)=2=2\times 1$ : $\Omega(2;1)$ est bien le centre de symétrie de $\mathcal{C}_f$.

\textbf{2. Résolution de $f(x)\le 3$.}
\[ \frac{x+1}{x-2}\le 3 \iff \frac{x+1-3(x-2)}{x-2}\le 0 \iff \frac{-2x+7}{x-2}\le 0. \]
Le numérateur s'annule en $x=\dfrac{7}{2}$, le dénominateur en $x=2$ (valeur interdite). Tableau de signes : le quotient est négatif à l'extérieur des racines.
\[ \mathcal{S}=\,]-\infty;2[\ \cup\ \left[\frac{7}{2};+\infty\right[. \]

\textbf{3. a)} On place $\Omega(2;1)$, on trace en pointillés les asymptotes $x=2$ et $y=1$, puis les deux branches de l'hyperbole.

\textbf{3. b)} \emph{Point de vigilance — cohérence de l'énoncé.} Calculons la tangente à $\mathcal{C}_f$ au point d'abscisse $3$ : $f(3)=\dfrac{4}{1}=4$ et $f'(3)=\dfrac{-3}{(3-2)^2}=-3$. La tangente en $3$ a pour équation :
\[ y = f'(3)(x-3)+f(3) = -3(x-3)+4 = -3x+13. \]
La droite $\mathcal{D} : y=3x-5$ n'est donc \textbf{pas} la tangente en $3$ (son coefficient directeur est $3\neq -3$) : l'énoncé comporte une coquille. La tangente correcte au point d'abscisse $3$ est $\boxed{y=-3x+13}$. La droite $\mathcal{D}$ est en réalité une \emph{sécante} à $\mathcal{C}_f$ (voir question suivante). Au Probatoire, si une telle situation se présente, le candidat doit vérifier la tangence par le double test $f(a)=y_{\mathcal{D}}(a)$ et $f'(a)=$ coefficient directeur, et signaler le résultat.

\textbf{3. c) Intersections de $\mathcal{C}_f$ et de $\mathcal{D}$.} On résout :
\begin{align*}
\frac{x+1}{x-2}=3x-5 &\iff x+1=(3x-5)(x-2) \quad (x\neq 2)\\
&\iff x+1 = 3x^2-11x+10\\
&\iff 3x^2-12x+9=0 \iff x^2-4x+3=0.
\end{align*}
$\Delta = 16-12=4$ ; $x_1=\dfrac{4-2}{2}=1$, $x_2=\dfrac{4+2}{2}=3$.

\emph{Vérification :} $f(1)=\dfrac{2}{-1}=-2$ et $3\times 1-5=-2$ \checkmark ; $f(3)=4$ et $3\times 3-5=4$ \checkmark.

Les points d'intersection sont $\boxed{A(1;-2)}$ et $\boxed{B(3;4)}$, que l'on lit également sur le graphique.

\textbf{Barème indicatif (sur 7 pts) :} 1.a) 1 pt ; 1.b) 1,5 pt ; 1.c) 1 pt ; 2) 1,25 pt ; 3.a) 0,5 pt ; 3.b) 0,75 pt ; 3.c) 1 pt.

\textbf{Points de vigilance :} interdire $x=2$ dans toute résolution ; pour le centre de symétrie, écrire explicitement la relation $f(a+h)+f(a-h)=2b$ ; pour la tangence, vérifier \emph{deux} conditions (le point appartient à la courbe \emph{et} les coefficients directeurs sont égaux) ; ne pas conclure sur les variations « sur $\mathbb{R}\setminus\{2\}$ » d'un seul tenant, mais bien intervalle par intervalle.
\end{corrige}

\newpage

\coursec{Fiche bilan}

\begin{popfigure}
\begin{tikzpicture}[
    mindmap,
    grow cyclic,
    every node/.style={concept, execute at begin node=\small\bfseries},
    concept color=popblue,
    level 1/.append style={level distance=4.5cm, sibling angle=360/6},
    level 2/.append style={level distance=3cm, sibling angle=360/4},
    branch/.style={concept color=#1, edge from parent/.style={draw=#1, thick}}
]
\node[concept, text width=3.5cm, minimum size=3.5cm] {Fonctions\\ Usuelles\\ (Tle Tech)}
    child[branch=poporange] {node[concept, text width=2.8cm] {Fonctions\\ Rationnelles\\ \& Homographiques}
        child {node[concept, text width=2.2cm] {Domaine, asymptotes}}
        child {node[concept, text width=2.2cm] {Variations, tableau}}
        child {node[concept, text width=2.2cm] {Primitives $\frac{u'}{u}$}}
    }
    child[branch=popgreen] {node[concept, text width=2.8cm] {Logarithme\\ Népérien $\ln$}
        child {node[concept, text width=2.2cm] {Déf: $\int_1^x \frac{dt}{t}$}}
        child {node[concept, text width=2.2cm] {Propriétés algébriques}}
        child {node[concept, text width=2.2cm] {Dérivée $\frac{1}{x}$}}
        child {node[concept, text width=2.2cm] {Limites, variations}}
    }
    child[branch=poppurple] {node[concept, text width=2.8cm] {Exponentielle\\ Népérienne $\exp$}
        child {node[concept, text width=2.2cm] {Réciproque de $\ln$}}
        child {node[concept, text width=2.2cm] {$\exp'=\exp$}}
        child {node[concept, text width=2.2cm] {Propriétés $e^{a+b}=e^ae^b$}}
        child {node[concept, text width=2.2cm] {Limites, variations}}
    }
    child[branch=poppink] {node[concept, text width=2.8cm] {Équations \&\\ Inéquations $\ln/\exp$}
        child {node[concept, text width=2.2cm] {Équivalences $\ln A = \ln B$}}
        child {node[concept, text width=2.2cm] {Conditions d'existence}}
        child {node[concept, text width=2.2cm] {Résolution graphique/algébrique}}
    }
    child[branch=popgold] {node[concept, text width=2.8cm] {Fonctions\\ Composées $\ln(ax+b)$, $\exp(ax+b)$}
        child {node[concept, text width=2.2cm] {Domaine ($ax+b>0$)}}
        child {node[concept, text width=2.2cm] {Dérivée, variations}}
        child {node[concept, text width=2.2cm] {Asymptotes, limites}}
        child {node[concept, text width=2.2cm] {Tableau de variations}}
    }
    child[branch=popdark] {node[concept, text width=2.8cm] {Applications\\ Techniques}
        child {node[concept, text width=2.2cm] {Croissance/décroissance (RC, RL)}}
        child {node[concept, text width=2.2cm] {Modélisation population}}
        child {node[concept, text width=2.2cm] {pH, échelle de Richter}}
        child {node[concept, text width=2.2cm] {Intégrales, aires}}
    };
\end{tikzpicture}
\legende{Carte mentale : Fonctions usuelles — Terminale Technique}
\end{popfigure}

\begin{bilan}

\courssub{Définitions et propriétés fondamentales}

\begin{definition}[Fonction homographique]
Une fonction de la forme $f(x)=\frac{ax+b}{cx+d}$ avec $ad-bc \neq 0$. Domaine : $\mathbb{R}\setminus\{-d/c\}$. Asymptotes : verticale $x=-\frac{d}{c}$, horizontale $y=\frac{a}{c}$.
\end{definition}

\begin{definition}[Logarithme népérien]
$\ln : ]0,+\infty[ \to \mathbb{R}$ définie par $\ln(x)=\int_1^x \frac{dt}{t}$. C'est une bijection strictement croissante. $\ln(1)=0$, $\ln(e)=1$.
\end{definition}

\begin{definition}[Exponentielle népérienne]
$\exp : \mathbb{R} \to ]0,+\infty[$ est la fonction réciproque de $\ln$. $\forall x\in\mathbb{R}, \ln(\exp(x))=x$ et $\forall x>0, \exp(\ln(x))=x$. On note $e^x=\exp(x)$.
\end{definition}

\courssub{Formules à connaître par cœur (Probatoire)}

\begin{center}
\begin{tabularx}{\linewidth}{|l|X|X|}\hline
\rowcolor{popblueL}
\textbf{Objet} & \textbf{Formule / Propriété} & \textbf{Condition} \\ \hline
Dérivée $\ln u$ & $(\ln u)' = \frac{u'}{u}$ & $u>0$ \\ \hline
Dérivée $\exp u$ & $(\exp u)' = u'\exp(u)$ & $u$ dérivable \\ \hline
Primitive $\frac{u'}{u}$ & $\int \frac{u'}{u} = \ln|u|+C$ & $u\neq 0$ \\ \hline
Primitive $\exp u$ & $\int u'\exp(u) = \exp(u)+C$ & $u$ dérivable \\ \hline
$\ln(ab)$ & $\ln a + \ln b$ & $a>0, b>0$ \\ \hline
$\ln(a/b)$ & $\ln a - \ln b$ & $a>0, b>0$ \\ \hline
$\ln(a^n)$ & $n\ln a$ & $a>0, n\in\mathbb{Z}$ \\ \hline
$\exp(a+b)$ & $\exp(a)\exp(b)$ & $\forall a,b\in\mathbb{R}$ \\ \hline
$\exp(a-b)$ & $\frac{\exp(a)}{\exp(b)}$ & $\forall a,b\in\mathbb{R}$ \\ \hline
$\exp(n a)$ & $(\exp(a))^n$ & $n\in\mathbb{Z}$ \\ \hline
Limites $\ln$ & $\lim_{x\to 0^+}\ln x = -\infty$, $\lim_{x\to+\infty}\ln x = +\infty$ &  \\ \hline
Limites $\exp$ & $\lim_{x\to -\infty}\exp x = 0$, $\lim_{x\to+\infty}\exp x = +\infty$ &  \\ \hline
Comparaison & $\lim_{x\to+\infty}\frac{\ln x}{x}=0$, $\lim_{x\to+\infty}\frac{x^n}{\exp x}=0$ & $n\in\mathbb{N}$ \\ \hline
\end{tabularx}
\end{center}

\courssub{Méthodes express}

\begin{methode}[Étudier une fonction homographique $f(x)=\frac{ax+b}{cx+d}$]
1. Domaine : $cx+d\neq 0$. \\
2. Asymptotes : verticale $x=-\frac{d}{c}$, horizontale $y=\frac{a}{c}$. \\
3. Dérivée : $f'(x)=\frac{ad-bc}{(cx+d)^2}$. Signe constant = signe de $ad-bc$. \\
4. Tableau de variations (2 branches). \\
5. Intersections axes : $f(0)=\frac{b}{d}$, $f(x)=0 \Leftrightarrow ax+b=0$.
\end{methode}

\begin{methode}[Résoudre une équation avec $\ln$ ou $\exp$]
1. \textbf{Conditions d'existence} (arguments $>0$ pour $\ln$). \\
2. Utiliser l'injectivité : $\ln A = \ln B \Leftrightarrow A=B>0$ ; $\exp A = \exp B \Leftrightarrow A=B$. \\
3. Simplifier (passer au $\exp$ ou au $\ln$). \\
4. Résoudre l'équation algébrique obtenue. \\
5. \textbf{Vérifier les solutions} dans les conditions initiales.
\end{methode}

\begin{methode}[Résoudre une inéquation avec $\ln$ ou $\exp$]
1. Conditions d'existence. \\
2. Utiliser la monotonie (croissante) : $\ln A < \ln B \Leftrightarrow 0<A<B$ ; $\exp A < \exp B \Leftrightarrow A<B$. \\
3. Résoudre l'inéquation algébrique. \\
4. Intersecter avec les conditions d'existence.
\end{methode}

\begin{methode}[Étudier $f(x)=\ln(ax+b)$ ou $f(x)=\exp(ax+b)$]
1. Domaine : $ax+b>0$ (pour $\ln$) ou $\mathbb{R}$ (pour $\exp$). \\
2. Dérivée : $f'(x)=\frac{a}{ax+b}$ ou $f'(x)=a\exp(ax+b)$. \\
3. Signe de $f'$ = signe de $a$ (pour $\ln$) ou signe de $a$ (pour $\exp$). \\
4. Limites aux bornes du domaine (asymptotes). \\
5. Tableau de variations. \\
6. Éventuel point d'inflexion ($f''$).
\end{methode}

\end{bilan}

\begin{aretenir}[Checklist avant le Probatoire]
$\square$ Je sais déduire le domaine, les asymptotes et le tableau de variations d'une fonction homographique. \\
$\square$ Je connais par cœur la définition de $\ln$ comme intégrale et ses propriétés algébriques ($\ln(ab)$, $\ln(a/b)$, $\ln(a^n)$). \\
$\square$ Je sais que $\exp$ est la réciproque de $\ln$ et j'applique $\exp(a+b)=\exp(a)\exp(b)$. \\
$\square$ Je dérivé correctement $\ln(u)$ ($\frac{u'}{u}$) et $\exp(u)$ ($u'\exp(u)$) \textbf{avec la dérivée de $u$}. \\
$\square$ Je donne les primitives : $\int \frac{u'}{u}=\ln|u|+C$ et $\int u'\exp(u)=\exp(u)+C$. \\
$\square$ Je résous $\ln A = \ln B$ en posant $A=B>0$ (conditions d'existence \textbf{obligatoires}). \\
$\square$ Je résous $\exp A = \exp B$ en posant $A=B$. \\
$\square$ Je résous les inéquations en utilisant la croissance stricte de $\ln$ et $\exp$. \\
$\square$ J'étudie complètement $x\mapsto \ln(ax+b)$ : domaine ($ax+b>0$), dérivée $\frac{a}{ax+b}$, variations, limites, asymptote verticale. \\
$\square$ J'étudie complètement $x\mapsto \exp(ax+b)$ : domaine $\mathbb{R}$, dérivée $a\exp(ax+b)$, variations, limites, asymptote horizontale. \\
$\square$ Je modélise une situation concrète (charge/décharge condensateur, croissance bactérienne, pH) par une exponentielle ou un logarithme et j'interprète les paramètres.
\end{aretenir}

\findecours

\end{document}
